
\documentclass[11pt,a4paper]{article}
\usepackage[spanish]{babel}
\usepackage[utf8]{inputenc}
\usepackage[T1]{fontenc}
\usepackage{geometry}
\usepackage{enumitem}
\geometry{margin=2.2cm}
\setlist[enumerate]{itemsep=0.15em, topsep=0.25em}
\title{Banco de preguntas de Dirección Integrada de Proyectos}
\author{}
\date{}
\begin{document}
\maketitle

\subsection*{Pregunta 1}
¿Cuál de las siguientes afirmaciones es CORRECTA?

\begin{enumerate}[label=\alph*)]
  \item \textbf{Sólo puede haber un camino crítico.}
  \item Un proyecto nunca puede tener una holgura negativa.
  \item Las tareas críticas cambiarán cada vez que cambie la fecha de finalización.
  \item El camino crítico ayuda a comprobar cuánto tiempo llevará el proyecto.
\end{enumerate}

\subsection*{Pregunta 2}
Un proyecto tiene un SPI (Índice de desempeño del cronograma) de 0,8 y un CPI (índice de desempeño del costo) de 1,1. Esto indica...

\begin{enumerate}[label=\alph*)]
  \item Que el proyecto está progresando más rápidamente y con un costo mayor que lo planificado
  \item \textbf{Que el proyecto está progresando más lentamente y con un costo menor que el planificado.}
  \item Que el proyecto está progresando más lentamente y con un costo mayor que el planificado
  \item Que el proyecto está progresando más rápidamente y con un costo menor que el planificado
\end{enumerate}

\subsection*{Pregunta 3}
¿Cuál de siguientes opciones describe MEJOR las principales restricciones de un proyecto?

\begin{enumerate}[label=\alph*)]
  \item Alcance, número de recursos y coste
  \item \textbf{Alcance, coste y tiempo}
  \item Alcance, calidad y tiempo
  \item Tiempo, coste y número de cambios
\end{enumerate}

\subsection*{Pregunta 4}
¿Durante cuál fase del ciclo de vida del proyecto tienen los interesados la mayor habilidad para influir sobre el producto final del proyecto y el costo final?

\begin{enumerate}[label=\alph*)]
  \item Ejecución del proyecto.
  \item \textbf{Iniciación del proyecto y aprobación.}
  \item Planificación del proyecto.
  \item Cierre del proyecto.
\end{enumerate}

\subsection*{Pregunta 5}
Durante la ejecución del proyecto un miembro de su equipo se dirige al Director del proyecto porque no está de cuál es el trabajo que debe de realizar en el proyecto. ¿Cuál de los siguientes documentos contiene la descripción de los paquetes de trabajo?

\begin{enumerate}[label=\alph*)]
  \item El plan de gestión de alcance
  \item Lista de actividades
  \item El Enunciado del Alcance del Proyecto
  \item El Diccionario de la EDP
\end{enumerate}

\subsection*{Pregunta 6}
Respecto a la teoría de la Cadena Crítica ¿qué afirmación es FALSA?

\begin{enumerate}[label=\alph*)]
  \item \textbf{No se incorporan "colchones" en la estimación de la duración de las actividades.}
  \item El proyecto se protege incorporando colchones selectivamente en posiciones críticas.
  \item Se planifican las actividades para empezar tan pronto como sea posible
  \item Se trata de evitar el "multitasking". Los recursos tienden a aplicarse al 100\% a una tarea, por períodos cortos de tiempo
\end{enumerate}

\subsection*{Pregunta 7}
Se considera que contratar un seguro es MÁS QUE OTRA COSA un ejemplo de:

\begin{enumerate}[label=\alph*)]
  \item Evitar el riesgo
  \item Mitigar el riesgo
  \item \textbf{Transferir el riesgo}
  \item Aceptar el riesgo
\end{enumerate}

\subsection*{Pregunta 8}
Cuáles de los siguientes aspectos de liderazgo es MÁS importante para un director de proyecto:

\begin{enumerate}[label=\alph*)]
  \item \textbf{Comunicación el liderazgo}
  \item Gestión de los grupos de interés
  \item Control del proyecto
  \item Experiencia técnica
\end{enumerate}

\subsection*{Pregunta 9}
Cuando se trata de gestionar y reducir los riesgos asociados con un proyecto, los eventos de riesgo pueden identificarse utilizando:

\begin{enumerate}[label=\alph*)]
  \item Una estructura de desglose de los riesgos.
  \item Una lista de nomenclatura de riesgos.
  \item Un análisis de simulación de riesgos de Monte Carlo.
  \item Una taxonomía de las causas del riesgo y su prevención.
\end{enumerate}

\subsection*{Pregunta 10}
Una desventaja de la utilización del Diagrama de Gantt es:

\begin{enumerate}[label=\alph*)]
  \item No se puede realizar considerando la utilización de los recursos
  \item Las holguras de las actividades
  \item \textbf{No muestra claramente las relaciones de interdependencia en proyectos complejos}
  \item No se puede representar con claridad el camino crítico
\end{enumerate}

\subsection*{Pregunta 11}
Durante la ejecución del proyecto, un miembro del equipo de proyectos informa al director de proyectos que un paquete de trabajo no ha cumplido con la métrica de calidad y que cree que no es posible cumplirla. El director de proyectos se reúne con todas las partes involucradas para analizar la situación. ¿En qué parte del proceso de gestión de la calidad está involucrado el director del proyecto?

\begin{enumerate}[label=\alph*)]
  \item Realizar el Aseguramiento de Calidad
  \item Control del Proyecto
  \item Realizar el Contro l de Calidad
  \item \textbf{Plan de Gestión de la Calidad 327}
\end{enumerate}

\subsection*{Pregunta 12}
¿Durante qué proceso de la gestión de los riesgos se determina que se va hacer una transferencia del riesgo?

\begin{enumerate}[label=\alph*)]
  \item Identificar los Riesgos
  \item Realizar el Análisis Cuantitativo de Riesgos
  \item \textbf{Planificar la Respuesta a los Riesgos}
  \item Controlar los Riesgos
\end{enumerate}

\subsection*{Pregunta 13}
El acta de constitución del proyecto:

\begin{enumerate}[label=\alph*)]
  \item \textbf{Confiere al director de proyecto la autoridad para asignar recursos al proyecto}
  \item Integra toda la información necesaria para dirigir el proyecto
  \item Desarrolla la matriz de asignación de responsabilidades para gestionar los recursos
  \item Es un documento formal que se desarrolla como parte de los procesos de la fase planificación
\end{enumerate}

\subsection*{Pregunta 14}
El director de proyectos está intentando recordar el método de comunicación preferido de un interesado. ¿Dónde puede encontrar esa información?

\begin{enumerate}[label=\alph*)]
  \item \textbf{Diagrama RACI}
  \item Matriz de evaluación del compromiso de los interesados
  \item Plan de gestión de los interesados
  \item Plan de gestión de los recursos humanos
\end{enumerate}

\subsection*{Pregunta 15}
En la Estructura de Descomposición del Proyecto:

\begin{enumerate}[label=\alph*)]
  \item \textbf{Se descompone el proyecto en elementos de trabajo hasta llegar a un último nivel de paquetes de trabajo}
  \item El nivel más bajo se denomina paquete de trabajo y tiene una duración mínima de una jornada de trabajo (8 horas)
  \item Se describen en detalle para cada paquete de trabajo, los responsables, plazos, actividades y entregables
  \item Se establece una estructura de trabajo que permite hacer una identificación de riesgos del proyecto
\end{enumerate}

\subsection*{Pregunta 16}
EXAMEN DIP\_2020 /EXAMEN DIP\_2020 ¿Cuáles son las principales diferencias entre el PERT y el CPM?

\begin{enumerate}[label=\alph*)]
  \item El número de actividades a controlar
  \item \textbf{La forma de establecer la duración de las actividades}
  \item El tiempo total de duración del proyecto
  \item La forma de relacionar las actividades
\end{enumerate}

\subsection*{Pregunta 17}
Si un evento de riesgo tiene 90 por ciento de posibilidades de ocurrir y su consecuencia sería de US\$10.000, ¿qué representa la cifra de US\$9.000?

\begin{enumerate}[label=\alph*)]
  \item \textbf{Valor monetario esperado}
  \item Valor del riesgo
  \item Valor actual
  \item Presupuesto de contingencias
\end{enumerate}

\subsection*{Pregunta 18}
Un índice de cumplimiento del programa (ICP en castellano SPI en inglés) de 0,76 significa:

\begin{enumerate}[label=\alph*)]
  \item Te has excedido del presupuesto
  \item Estás progresando a 76 por ciento del ritmo que se planificó originalmente
  \item Estás progresando a 24 por ciento del ritmo que se planificó originalmente
  \item Te has adelantado al cronograma Johana se encuentra definiendo un proceso del área de comunicaciones en el que se describe quién necesita información, cuándo la
\end{enumerate}

\subsection*{Pregunta 19}
Usted está realizando el análisis de interesados y ha identificado la categoría de clasificación e impacto potencial de cada interesado. ¿En qué se deberán enfocar usted y su equipo a continuación?

\begin{enumerate}[label=\alph*)]
  \item Determinar cómo los interesados clave responderán a varias situaciones.
  \item Verificar con el patrocinador del proyecto para asegurar la alineación con el acta del proyecto
  \item Distribuir información a los interesados adecuados.
  \item \textbf{Desarrollar el registro de interesados.}
\end{enumerate}

\subsection*{Pregunta 20}
Basándose en el ratio beneficio/coste, ¿cuál de los siguientes proyectos escogería?

\begin{enumerate}[label=\alph*)]
  \item \textbf{Ratio coste/beneficio de 0,1 y coste de 1000\$}
  \item Ratio beneficio/coste de 0,4 y beneficio de 3000\$
  \item Ratio coste/beneficio de 1,1 y beneficio 20000\$
  \item Ratio beneficio/coste de 0,6 y beneficio de 6000\$.
\end{enumerate}

\subsection*{Pregunta 21}
¿Con respecto al proceso Controlar el Alcance, cuál de las siguientes es CORRECTA?

\begin{enumerate}[label=\alph*)]
  \item Una definición efectiva del alcance puede llevar a un enunciado del alcance del proyecto más completo.
  \item \textbf{El proceso Controlar el Alcance debe hacerse antes de la planificación del alcance.}
  \item La forma más efectiva de controlar el alcance es controlar el cronograma . 190 © 201 3
\end{enumerate}

\subsection*{Pregunta 22}
¿Cuál de las siguientes técnicas NO es una técnica de reducción de la programación? :

\begin{enumerate}[label=\alph*)]
  \item Reducción de recursos
  \item Uso de horas extra
  \item Fast-tracking
  \item Reducción del alcance
\end{enumerate}

\subsection*{Pregunta 23}
¿Cuál de los siguientes eventos de riesgo es el que tiene la MAYOR probabilidad de interferir con el cumplimiento del objetivo del cronograma del proyecto?

\begin{enumerate}[label=\alph*)]
  \item Retrasos en la obtención de las aprobaciones necesarias
  \item Incrementos sustanciales en el costo de los materiales adquiridos
  \item Disputas contractuales que tienen como consecuencia reclamos de aumentos de pago
  \item \textbf{Retraso de la reunión de revisión posterior a la implementación que se tenía prevista}
\end{enumerate}

\subsection*{Pregunta 24}
¿Cuándo tiene un Grupo de Interés una MAYOR influencia en un proyecto?

\begin{enumerate}[label=\alph*)]
  \item \textbf{A lo largo de todo el proyecto}
  \item Al inicio del proyecto
  \item Al final del proyecto
  \item A la mitad del proyecto
\end{enumerate}

\subsection*{Pregunta 25}
¿Cómo definiría la Cadena Crítica?

\begin{enumerate}[label=\alph*)]
  \item Ninguna de las que se mencionan
  \item Secuencia de eventos dependientes que limitan la finalización temprana de un proyecto; incluye dependencias de recursos
  \item Camino más largo de los posibles teniendo en cuenta las restricciones temporales
  \item \textbf{Secuencia de eventos dependientes que limitan la finalización temprana de un proyecto; incluye dependencias de recursos así como de tareas}
\end{enumerate}

\subsection*{Pregunta 26}
Durante la fase de inicio del proyecto, buscamos los factores críticos de éxito para determinar cómo vamos a:

\begin{enumerate}[label=\alph*)]
  \item Cerrar el proyecto.
  \item Monitorear el desempeño de las tareas.
  \item Cumplir con los hitos principales.
  \item \textbf{Lograr las metas del proyecto.}
\end{enumerate}

\subsection*{Pregunta 27}
El director del proyecto acaba de recibir un cambio del cliente que no afecta el cronograma del proyecto y es fácil de completar . ¿Qué debe hacer PRIMERO el director del proyecto?

\begin{enumerate}[label=\alph*)]
  \item Hacer que el cambio suceda lo más pronto posible.
  \item Contactar al patrocinador del proyecto para obtener permiso.
  \item Dirigirse al comité de control de cambios.
  \item Evaluar el impacto sobre otras restricciones del proyecto.
\end{enumerate}

\subsection*{Pregunta 28}
El Director del Proyecto puede usar \_\_\_\_\_\_\_\_\_\_\_ para estar seguro de que los miembros del equipo del proyecto entienden con claridad cuál es el trabajo incluido en los paquetes de trabajo

\begin{enumerate}[label=\alph*)]
  \item La programación
  \item El Enunciado del Alcance del Proyecto
  \item El alcance del producto
  \item El Diccionario de la EDP
\end{enumerate}

\subsection*{Pregunta 29}
El objetivo clave de la gestión de los interesados es:

\begin{enumerate}[label=\alph*)]
  \item Comunicación.
  \item Coordinación.
  \item Satisfacción.
  \item \textbf{Relaciones.}
\end{enumerate}

\subsection*{Pregunta 30}
En la fase de planificación de un proyecto para el desarrollo de un nuevo producto se especifica que éste ha de servir tanto para jóvenes como para adultos. ¿De qué tipo de requisito se trata?

\begin{enumerate}[label=\alph*)]
  \item \textbf{Funcional}
  \item No funcional
  \item De producto
  \item De proyecto
\end{enumerate}

\subsection*{Pregunta 31}
En una reunión con los interesados durante la iniciación del proyecto, el patrocinador del proyecto deja claro que este utilizará la tecnología actual de la compañía, recursos humanos existentes y que el presupuesto no excederá \$500,000. ¿Qué es lo que el patrocinador está describiendo?

\begin{enumerate}[label=\alph*)]
  \item Supuestos.
  \item Límites.
  \item \textbf{Restricciones.}
  \item Alcance.
\end{enumerate}

\subsection*{Pregunta 32}
Está buscando una manera de cuantificar los riesgos de su proyecto. Todos los siguientes son métodos cuantitativos del análisis de riesgo con EXCEPCIÓN de:

\begin{enumerate}[label=\alph*)]
  \item Análisis del valor monetario esperado.
  \item Análisis de sensibilidad.
  \item \textbf{Análisis de probabilidad e impacto.}
  \item Modelar y simular.
\end{enumerate}

\subsection*{Pregunta 33}
EXAMEN DIP\_2020 La reserva para contingencias de los costos debe ser:

\begin{enumerate}[label=\alph*)]
  \item Ocultada para prevenir que la dirección decida no permitir la reserva
  \item \textbf{Añadida a cada actividad para proveerle al cliente un camino crítico más corto.}
  \item Añadida a los costos del proyecto para que pueda afrontar los riesgos
  \item Mantenida por la dirección para cubrir sobrecostos
\end{enumerate}

\subsection*{Pregunta 34}
Ignacio es el director de un nuevo proyecto y el cliente le exige que acorte la duración del mismo. ¿Cuál de las siguientes opciones es la mejor a tener en cuenta?

\begin{enumerate}[label=\alph*)]
  \item \textbf{Evaluar las implicaciones asociadas a los requerimientos del cliente}
  \item Reunirse con el equipo del proyecto y revisar la programación
  \item Reasignar recursos de las actividades críticas
  \item Acortar la programación para tener al cliente contento
\end{enumerate}

\subsection*{Pregunta 35}
La EDP puede ser pensada para como una ayuda efectiva para las comunicaciones con \_\_\_\_\_\_\_\_\_\_\_\_ (responda la mejor opción)

\begin{enumerate}[label=\alph*)]
  \item Director del Proyecto Esta tiene trampa ya que grupo de interés
  \item Cliente
  \item \textbf{Equipo del Proyecto}
  \item Grupos de Interés
\end{enumerate}

\subsection*{Pregunta 36}
La herramienta que proporciona el trabajo que deben realizar los miembros del equipo en cada uno de sus paquetes de trabajo es: :

\begin{enumerate}[label=\alph*)]
  \item El alcance del producto
  \item El enunciado del alcance del proyecto
  \item El diccionario de la EDP
  \item El cronograma.
\end{enumerate}

\subsection*{Pregunta 37}
La principal diferencia entre un proceso y un proyecto es: :

\begin{enumerate}[label=\alph*)]
  \item El proyecto es temporal y singular
  \item El proyecto es singular y repetitivo.
  \item El proceso tiene una duración limitada aunque sea repetitiva.
  \item Los proyectos se componen de procesos.
\end{enumerate}

\subsection*{Pregunta 38}
Las relaciones lógicas o vínculos entre las actividades de un proyecto pueden ser del tipo:

\begin{enumerate}[label=\alph*)]
  \item Fin-Comienzo y Comienzo-Fin
  \item Predecesora-Sucesora
  \item \textbf{Comienzo-Comienzo, Fin-Fin, Fin-Comienzo y Comienzo-Fin}
  \item Salida-Entrada
\end{enumerate}

\subsection*{Pregunta 39}
María está dirigiendo un proyecto de ingeniería. Tiene que entregar el informe de proyecto definitivo dentro de dos semanas y ha repartido el trabajo entre su equipo. Éste recela de María porque antes tenían otra directora que les gustaba más. Por otra parte, María, aunque sepa que es un equipo competente, les ha solicitado que le entreguen a ella los informes para revisarlos antes de enviarlos al cliente. ¿Cree usted que ha actuado correctamente?

\begin{enumerate}[label=\alph*)]
  \item No, si María sabe que el equipo es competente tiene que fiarse de ellos
  \item No, si pedir los informes a los miembros del equipo le va a causar problemas, tendría que coger los informes directamente, a escondidas, sin comentar nada al equipo del proyecto y revisarlos
  \item \textbf{Sí, aunque el equipo sea competente, el trabajo de María también consiste en revisar esos informes, como entregables, a fin de determinar si el trabajo y los productos cumplen con los requisitos y los criterios de aceptación del producto}
  \item No, pero depende del volumen de trabajo que tenga María
\end{enumerate}

\subsection*{Pregunta 40}
Para dirigir un proyecto de manera efectiva se debe descomponer el trabajo en piezas pequeñas. ¿Cuál de las siguientes afirmaciones NO describe hasta dónde descomponer el trabajo?

\begin{enumerate}[label=\alph*)]
  \item \textbf{Hasta que lógicamente no se puede subdividir más}
  \item Hasta que puede estimarse de modo realista
  \item Hasta que puede hacerse por una persona
  \item Hasta que tiene una conclusión significativa
\end{enumerate}

\subsection*{Pregunta 41}
¿Podría identificar cuál de las siguientes definiciones hace referencia a la estimación paramétrica?

\begin{enumerate}[label=\alph*)]
  \item La técnica de la estimación paramétrica consiste en determinar valores asignados a parámetros cuantificables
  \item \textbf{La técnica de la estimación paramétrica utiliza una relación estadística entre los datos históricos y otras variables para calcular una estimación de parámetros de una actividad tales como costo, presupuesto y duración}
  \item La técnica de la estimación paramétrica consiste en determinar los parámetros adecuados para la evaluación de los proyectos
  \item La técnica de la estimación paramétrica es un método para estimar los componentes del trabajo. El costo de cada paquete de trabajo o de cada actividad se calcula con el mayor nivel de detalle
\end{enumerate}

\subsection*{Pregunta 42}
¿Qué es un programa?

\begin{enumerate}[label=\alph*)]
  \item Una regulación gubernamental
  \item \textbf{Conjunto de proyectos con un objetivo común que gestionados de forma conjunta se obtienen beneficios.}
  \item Un grupo de proyectos sin relación entre sí que se dirigen de forma coordinada
  \item Una iniciativa empezada por la dirección
\end{enumerate}

\subsection*{Pregunta 43}
Se ha autorizado un proyecto para construir un almacén de residuos peligrosos y Pepe ha sido asignado como director de proyecto. Ha desarrollado un plan para la gestión de las comunicaciones. Para asegurar que el plan satisface las necesidades de los interesados ¿qué proceso debería seguir Pepe?

\begin{enumerate}[label=\alph*)]
  \item Revisar las comunicaciones
  \item \textbf{Gestionar las comunicaciones}
  \item Controlar las comunicaciones
  \item Planificar la gestión de las comunicaciones
\end{enumerate}

\subsection*{Pregunta 44}
Si se le pide al director de proyecto un estimado rápido de los costos de un proyecto, una estimación \_\_\_\_\_\_\_\_\_\_ es adecuada. Si el mismo director de proyecto le pide un estimado bien elaborado y le da el tiempo necesario para prepararlo, entonces una estimación \_\_\_\_\_\_\_\_\_\_ es la más adecuada.

\begin{enumerate}[label=\alph*)]
  \item Análoga, descendente.
  \item \textbf{Descendente, ascendente.}
  \item Paramétrica, análoga.
  \item Ascendente, descendente.
\end{enumerate}

\subsection*{Pregunta 45}
¿Si un proyecto tiene un BAC (presupuesto hasta la conclusión) de 43.000€ y el EV (valor ganado) es de 13.200€, qué otro valor es necesario para calcular su variación hasta la conclusión (VAC) utilizando los datos proporcionados?

\begin{enumerate}[label=\alph*)]
  \item El PV (valor planificado)
  \item El SV (variación del cronograma)
  \item El AC (costo real)
  \item EL CPI (índice del desempeño del costo)
\end{enumerate}

\subsection*{Pregunta 46}
Un director de proyectos acaba de terminar el plan de respuesta a los riesgos de un proyecto de ingeniería con valor de US\$387.000. ¿Cuál debe ser su SIGUIENTE acción?

\begin{enumerate}[label=\alph*)]
  \item \textbf{Realizar una reevaluación de riesgos del proyecto}
  \item Determinar la calificación general del riesgo del proyecto.
  \item Comenzar a analizar los riesgos que aparecen en los planos del proyecto
  \item Añadir paquetes de trabajo a la estructura de descomposición del proyecto (EDP)
\end{enumerate}

\subsection*{Pregunta 47}
Un ejemplo sobre cómo el gobierno del proyecto puede añadir riesgo a los proyectos individuales es: Un ejemplo sobre cómo el gobierno del proyecto puede añadir riesgo a los proyectos individuales es: opción:

\begin{enumerate}[label=\alph*)]
  \item Soportar multiples metodologias que encajan con diferentes tipos de proyectos a. Soportar múltiples metodologías que encajan con diferentes tipos de proyectos
  \item Requerir casos de negocio previa aprobacidn del proyecto Ob. Requerir casos de negocio previa aprobación del proyecto cuestién
  \item Tomar los pasos para convertirse en una organizacion mas "dirigida por proyectos" €. Tomar los pasos para convertirse en una organización más "dirigida por proyectos"
  \item \textbf{Aprobar mas proyectos que los que los fondos, recursos y procesos de la organizacion pueden soportar d. Aprobar más proyectos que los que los fondos, recursos y procesos de la organización pueden soportar}
\end{enumerate}

\subsection*{Pregunta 48}
Un equipo del proyecto se encuentra elaborando un plan para la dirección del proyecto cuando de pronto la gerencia le solicita que identifique riesgos del proyecto y proporcione algún tipo de salida cualitativa lo antes posible . ¿Qué debe proporcionar el equipo del proyecto?

\begin{enumerate}[label=\alph*)]
  \item Lista de riesgos del proyecto en orden de prioridad
  \item Disparadores de riesgos
  \item \textbf{Reservas para contingencias}
  \item Probabilidad de cumplir con los objetivos de tiempo y costo 449
\end{enumerate}

\subsection*{Pregunta 49}
Un cronograma detallado del proyecto sólo puede crearse después de:

\begin{enumerate}[label=\alph*)]
  \item El plan para la dirección del proyecto
  \item \textbf{La estructura de descomposición del trabajo}
  \item El presupuesto del proyecto
  \item La evaluación detallada de riesgos
\end{enumerate}

\subsection*{Pregunta 50}
Usted está en las últimas semanas de un proyecto de 12 meses de duración. Su equipo está trabajando en el entregable final y las pruebas de integración. La relación de trabajo entre dos miembros clave del equipo se ha roto y ahora no se están hablando el uno al otro. Usted quiere asegurarse de que las pruebas de integración se den sin problemas para que su proyecto sea completado de acuerdo al hito programado. Necesita la cooperación de los dos miembros del equipo para hacerlo. El curso de acción preferido es:

\begin{enumerate}[label=\alph*)]
  \item Escalar el problema al patrocinador del proyecto y pedir que involucre al gerente funcional de los miembros del equipo para resolver el problema.
  \item Pedir a la alta dirección que amplíen el plazo del proyecto otras dos semanas para asegurar cualquier problema causado por la mala voluntad entre los dos miembros del equipo no tendrá un impacto negativo en el cumplimiento de la fecha de finalización del proyecto.
  \item No hacer nada. Quedan sólo dos semanas en el proyecto, y es poco probable que pueda hacer cualquier cosa dentro de ese tiempo para resolver el problema entre los dos miembros del equipo.
  \item \textbf{Hablar con ambos y explicarles que necesitan seguir adelante con su trabajo en el corto tiempo que queda en el proyecto.}
\end{enumerate}

\subsection*{Pregunta 51}
Usted está liderando un proyecto para una empresa que ha sufrido problemas de calidad en proyectos anteriores. Desea asegurarse de que el proyecto emplee todos los procesos de calidad necesarios para cumplir con los requisitos. Esto es:

\begin{enumerate}[label=\alph*)]
  \item Inspección de la calidad
  \item Control de calidad
  \item \textbf{Gestión de calidad}
  \item Costos de calidad
\end{enumerate}

\subsection*{Pregunta 52}
Al considerar a quién comunicar y qué información necesita comunicarse, el director de proyecto debe:

\begin{enumerate}[label=\alph*)]
  \item \textbf{Revisar el Plan de gestión de las comunicaciones del proyecto}
  \item Comunicar todos los detalles posibles a todos los interesados del proyecto
  \item Proporcionar únicamente información positiva sobre el proyecto
  \item Adaptar la información a las necesidades de los interesados del proyecto realizando un análisis de los requisitos de comunicación
\end{enumerate}

\subsection*{Pregunta 53}
Al principio de la vida de tu proyecto, mantienes una discusión con el patrocinador respecto de qué técnicas de estimación se deben usar. Tú deseas utilizar algún tipo de juicio de expertos, pero el patrocinador argumenta a favor de la estimación análoga. Lo MEJOR sería:

\begin{enumerate}[label=\alph*)]
  \item Convenir que se haga una estimación análoga, pues es parte del juicio de expertos .
  \item Sugerir el costeo del ciclo de vida a manera de conciliación.
  \item Determinar por qué el patrocinador desea una estimación tan precisa .
  \item Intentar convencer al patrocinador de que permita el juicio de expertos, ya que por lo general resulta más preciso . 283
\end{enumerate}

\subsection*{Pregunta 54}
Algunos factores clave que determinan las duraciones de las actividades son:

\begin{enumerate}[label=\alph*)]
  \item Motivación del personal, número de recursos y tecnología 1.00
  \item Ley de rendimientos decrecientes, éxito de proyectos anteriores y riesgos del proyecto
  \item Calendario laboral, número de recursos y experiencia del director de proyectos
  \item Reserva de gestión y contingencia del proyecto
\end{enumerate}

\subsection*{Pregunta 55}
Considere el siguiente histograma de recursos. Basado únicamente en la disponibilidad de recursos, ¿cuál de los siguientes es verdadero?

\begin{enumerate}[label=\alph*)]
  \item Existe un excedente de recursos durante partes de Abril y Mayo.
  \item Existe una falta de recursos durante partes de Abril y Mayo.
  \item Existe un excedente de recursos durante partes de Mayo y Junio.
  \item Existe una falta de recursos durante partes de Mayo y Junio.
\end{enumerate}

\subsection*{Pregunta 56}
¿Cuál de las siguientes definiciones describe mejor el camino crítico?

\begin{enumerate}[label=\alph*)]
  \item \textbf{El camino más corto}
  \item El camino más largo con el riesgo mayor
  \item El camino más largo pero no necesariamente el de mayor riesgo
  \item El camino más largo con el coste más bajo
\end{enumerate}

\subsection*{Pregunta 57}
Cuál de las siguientes no es una característica de un requisito:

\begin{enumerate}[label=\alph*)]
  \item No ambiguo
  \item Trazable 1.00
  \item Concreto
  \item Mensurable
\end{enumerate}

\subsection*{Pregunta 58}
¿Cuál de los siguientes proporciona un vínculo entre la definición del alcance del proyecto y la organización/ individuos ejecutores? :

\begin{enumerate}[label=\alph*)]
  \item Esquema de codificación único
  \item Estructura de Descomposición del Proyecto
  \item Diccionario de la EDP
  \item Matriz de asignación de responsabilidades
\end{enumerate}

\subsection*{Pregunta 59}
¿Cuál de los siguientes se incluye en el acta de constitución del proyecto?

\begin{enumerate}[label=\alph*)]
  \item Una estrategia de gestión de riesgos
  \item \textbf{Estimados de los paquetes de trabajo}
  \item Estimados detallados de los recursos
  \item El caso de negocio del proyecto 150
\end{enumerate}

\subsection*{Pregunta 60}
Cuando hablamos del costo total incurrido y registrado para llevar a cabo un trabajo realizado para una actividad o componente de la estructura de desglose del trabajo, nos estamos refiriendo a...

\begin{enumerate}[label=\alph*)]
  \item El CV (variación del costo)
  \item El AC (costo real)
  \item El EV (valor ganado)
  \item El PV (valor planificado).
\end{enumerate}

\subsection*{Pregunta 61}
Cuando un director de proyectos participa en negociaciones, las habilidades de comunicación no verbal son de: 00

\begin{enumerate}[label=\alph*)]
  \item Poca importancia
  \item Importantes sólo cuando se involucran los objetivos de los costes y el cronograma
  \item Importantes para garantizar que gane la negociación
  \item Gran importancia
\end{enumerate}

\subsection*{Pregunta 62}
El cierre del proyecto incluye todas las siguientes acciones, excepto:

\begin{enumerate}[label=\alph*)]
  \item Actualizar los activos de los procesos de la organización.
  \item Documentar hasta qué grado se cerró de forma correcta cada fase del proyecto después de su conclusión.
  \item Hacer entrega del producto del proyecto.
  \item \textbf{Determinar medidas del desempeño.}
\end{enumerate}

\subsection*{Pregunta 63}
El cliente principal de un proyecto ha solicitado un cambio en la aplicación durante la prueba del usuario . Como director del proyecto, ¿cómo resolverías MEJOR esta polémica?

\begin{enumerate}[label=\alph*)]
  \item Desarrollar un plan de mitigación de riesgos.
  \item \textbf{Crear una solicitud formal de cambio.}
  \item Informar al patrocinador del proyecto acerca de los cambios en el alcance, costo y cronograma.
  \item Asegurar que el cambio al alcance cumpla con todas las disposiciones relevantes del contrato.
\end{enumerate}

\subsection*{Pregunta 64}
El nivel de descomposición para los productos entregables varía...

\begin{enumerate}[label=\alph*)]
  \item \textbf{Según complejidad del proyecto}
  \item Según tamaño, complejidad y financiación del proyecto
  \item Según tamaño y complejidad del proyecto
  \item Según tamaño del proyecto
\end{enumerate}

\subsection*{Pregunta 65}
El proyecto es temporal y singular La conversión de las actividades de un proyecto en un diagrama hace referencia al proceso:

\begin{enumerate}[label=\alph*)]
  \item Desarrollar el cronograma
  \item Secuenciar las actividades
  \item Planificar la gestión del cronograma
  \item Definir las actividades
\end{enumerate}

\subsection*{Pregunta 66}
El rol de cada interesado es determinado por:

\begin{enumerate}[label=\alph*)]
  \item El interesado y el patrocinador.
  \item El director del proyecto y el patrocinador .
  \item \textbf{El equipo y el director del proyecto.}
  \item El director del proyecto y el interesado.
\end{enumerate}

\subsection*{Pregunta 67}
En el plan de gestión de costes de un proyecto se incluye normalmente:

\begin{enumerate}[label=\alph*)]
  \item Unidades de medida, cuentas de control y umbrales de variación permitidos
  \item Curva de la S de costes acumulados del proyecto
  \item \textbf{Reserva de gestión y reserva de contingencias}
  \item Resultados de la estimación del presupuesto del proyecto
\end{enumerate}

\subsection*{Pregunta 68}
En la planificación de la gestión de las comunicaciones:

\begin{enumerate}[label=\alph*)]
  \item Es cuando el Director de Proyecto transmite a su equipo la necesidad de cambios en los requisitos del proyecto 1.00
  \item Debe seleccionarse a aquellos grupos de interés con los que no será necesario tener una estrategia de comunicación
  \item Se han de tener en cuenta y documentar los métodos de almacenamiento, recuperación y disposición final de la información
  \item Se suele comenzar en fase avanzada del proyecto
\end{enumerate}

\subsection*{Pregunta 69}
¿En qué grupo de procesos de la dirección de proyectos se crea el presupuesto detallado del proyecto?

\begin{enumerate}[label=\alph*)]
  \item Ejecución
  \item \textbf{Planificación}
  \item Antes del proceso de dirección de proyectos
  \item Iniciación
\end{enumerate}

\subsection*{Pregunta 70}
Está realizando un análisis para mejorar su proceso de control de calidad. Tiene varias opciones y desea seleccionar la que proporciona los mejores resultados con la menor inversión. ¿Qué técnica de planificación de la calidad debe utilizar?

\begin{enumerate}[label=\alph*)]
  \item Estudios comparativos (benchmarking).
  \item Diagramas de control.
  \item Diseño de experimentos.
  \item Análisis costo-beneficio.
\end{enumerate}

\subsection*{Pregunta 71}
EXAMEN DIP\_2020 El estudio de programación del proyecto incluye, entre otras etapas:

\begin{enumerate}[label=\alph*)]
  \item Selección de todos los subcontratistas
  \item Establecimiento de relaciones contractuales
  \item Desarrollo completo del proceso tecnológico
  \item \textbf{Asignación de recursos y plazos de ejecución}
\end{enumerate}

\subsection*{Pregunta 72}
EXAMEN DIP\_2020 /EXAMEN DIP\_2020 Una forma de obtener la variación del costo (CV) es partir del valor ganado (EV) y...

\begin{enumerate}[label=\alph*)]
  \item Multiplicarlo por el costo real (AC)
  \item Restar el valor planificado (PV).
  \item Multiplicarlo por el valor planificado (PV).
  \item \textbf{Restar el costo real (AC)}
\end{enumerate}

\subsection*{Pregunta 73}
EXAMEN DIP\_2020 /EXAMEN DIP\_2020 La Directora de un proyecto desea informar sobre los resultados reales del mismo en comparación con los resultados planificados. Para ello debe utilizar:

\begin{enumerate}[label=\alph*)]
  \item \textbf{Informe de variación}
  \item Informe de estado
  \item Informe de tendencia
  \item Informe de proyección
\end{enumerate}

\subsection*{Pregunta 74}
Una forma de obtener la estimación a la conclusión (EAC) es partir del presupuesto hasta la conclusión y...

\begin{enumerate}[label=\alph*)]
  \item Multiplicarlo por el valor planificado (PV).
  \item \textbf{Dividirlo entre el índice de desempeño del costo (CPI)}
  \item Multiplicarlo por el costo real (AC)
  \item Dividirlo entre el coste real (AC)
\end{enumerate}

\subsection*{Pregunta 75}
Hay varias actividades de ejecución que están en marcha en tu proyecto . Comienzas a preocuparte sobre la precisión de los reportes de progreso que tu equipo está haciendo. ¿Cómo puedes verificar si hay un problema?

\begin{enumerate}[label=\alph*)]
  \item Auditorías de calidad
  \item Informes de cuantificación de riesgos
  \item Análisis de regresión
  \item Análisis Monte Carlo
\end{enumerate}

\subsection*{Pregunta 76}
La descomposición de todo el trabajo de proyecto generalmente implica...

\begin{enumerate}[label=\alph*)]
  \item \textbf{Todas las respuestas son correctas}
  \item Estructurar y organizar la EDT
  \item Descomponer los niveles superiores de la EDT en componentes detallados de nivel inferior
  \item Identificar los productos entregables y el trabajo relacionado
\end{enumerate}

\subsection*{Pregunta 77}
La diferencia entre reservas para contingencias y para gestión es que...

\begin{enumerate}[label=\alph*)]
  \item La reserva para contingencias sólo se disponen bajo acuerdo entre los miembros del equipo y la dirección mientras que la disponibilidad de la reserva de gestión es discrecional
  \item La reserva para contingencias se incluye en la línea base de costos y la de gestión en el presupuesto
  \item La reserva para contingencias se asocia a la dirección general de una empresa y la de gestión se asocia a los riesgos del proyecto
  \item \textbf{Ambas han de ir incluidas en la línea base de costos}
\end{enumerate}

\subsection*{Pregunta 78}
La estimación por analogía: :

\begin{enumerate}[label=\alph*)]
  \item Se utiliza con mayor frecuencia durante la ejecución del proyecto
  \item Es una técnica de estimación descendente
  \item Es una técnica de estimación ascendente
  \item Calcula las estimaciones utilizando los costes históricos reales detallados
\end{enumerate}

\subsection*{Pregunta 79}
Las estrategias alternativas de respuesta al riesgo para oportunidades o eventos de riesgo positivos incluyen: 00

\begin{enumerate}[label=\alph*)]
  \item Evitar, compartir y mejorar
  \item Explotar, compartir y mejorar
  \item Explotar, compartir y aceptar
  \item Explotar, transferir y mejorar
\end{enumerate}

\subsection*{Pregunta 80}
Las tolera ncias al riesgo se determinan con el fin de ayudar :

\begin{enumerate}[label=\alph*)]
  \item \textbf{A que el equipo califique los riesgos del proyecto.}
  \item A que el director del proyecto estime el proyecto.
  \item A que el equipo haga un cronograma del proyecto .
  \item A que la gerencia sepa de qué manera se comportarán los demás directores en el proyecto .
\end{enumerate}

\subsection*{Pregunta 81}
Los procesos de Dirección de Proyectos:

\begin{enumerate}[label=\alph*)]
  \item Pueden generar muchos proyectos
  \item No incorporan una metodología
  \item Son diferentes para cada industria
  \item \textbf{Son iguales para todos los proyectos}
\end{enumerate}

\subsection*{Pregunta 82}
Luis ha identificado un riesgo en donde se utiliza el antiguo hardware para ejecutar un software más actual. Luis ha recibido la aprobación para proseguir actualizando el hardware para asegurar que el riesgo no tenga lugar. ¿Qué método de respuesta al riesgo está aplicando Luis?

\begin{enumerate}[label=\alph*)]
  \item Evitar el riesgo
  \item Aceptar el riesgo
  \item \textbf{Mitigar el riesgo}
  \item Transferir el riesgo
\end{enumerate}

\subsection*{Pregunta 83}
Una manera de abordar la reducción de tiempo de un proyecto, manteniendo el alcance es:

\begin{enumerate}[label=\alph*)]
  \item \textbf{Incrementando el coste}
  \item Reduciendo los recursos clave del camino crítico
  \item Revisando la definición de alcance del proyecto
  \item Incrementando la reserva de gestión del proyecto
\end{enumerate}

\subsection*{Pregunta 84}
María considera que sólo existe una única estrategia de respuesta a los riesgos, evitarlos, de modo que éste es su planteamiento ante el cliente. ¿Está de acuerdo con el modo de proceder de María?

\begin{enumerate}[label=\alph*)]
  \item Depende de la priorización de los riesgos
  \item \textbf{No, ya que hay disponibles varias estrategias de respuesta a los riesgos (evitarlos, transferirlos, etc.).}
  \item Depende de la percepción que de los riesgos tenga su cliente
  \item Sí, la única estrategia válida es evitar cualquier riesgo
\end{enumerate}

\subsection*{Pregunta 85}
María está revisando los informes de estado de su equipo. En general parece que han logrado realizar trabajo por valor de 30.000€. Sin embargo, al añadir todas las tareas que su equipo debería haber realizado, su valor es de 40.000€. ¿Cuál sería la acción correcta en esta situación?

\begin{enumerate}[label=\alph*)]
  \item \textbf{Evaluar la ruta crítica para ver si el retraso en el crono grama afecta a la ruta crítica.}
  \item Añadir más recursos para compensar estar por detrás de cronograma.
  \item Ninguna. tienen una tendencia bajo presupuesto y están bien.
  \item Con un SPI de 0.75 se encuentra demasiado por detrás del cronograma para solucionarlo. Es hora de informar al cliente de que se retrasará.
\end{enumerate}

\subsection*{Pregunta 86}
Marta acaba de ser asignada directora de proyecto para la mejora de hardware. Hace tres años se realizó un proyecto similar. Para prepararse para este proyecto, Marta haría bien en realizar todo lo siguiente EXCEPTO:

\begin{enumerate}[label=\alph*)]
  \item Revisar la información de lecciones aprendidas de proyectos pasados
  \item Revisar la retroalimentación de los interesados en proyectos previos
  \item \textbf{Ignorar el proyecto previo porque cada proyecto es único}
  \item Revisar la documentación del proyecto previo
\end{enumerate}

\subsection*{Pregunta 87}
Para gestionar un proyecto de manera eficiente el trabajo debería ser descompuesto en componentes más pequeños. ¿Cuál de las siguientes afirmaciones NO describe cuánto debe de ser descompuesto del trabajo?

\begin{enumerate}[label=\alph*)]
  \item Hasta que pueda ser estimado de manera realista La clave está en "única". Si no sería verdadera. Ya que en las diapos dice persona
  \item \textbf{Hasta que pueda ser hecho por una única persona u organización}
  \item Hasta que no puede ser descompuesto más de manera lógica
  \item Hasta que pueda tener un resultado significativo
\end{enumerate}

\subsection*{Pregunta 88}
Qué aspecto es más importante en el cierre del proyecto:

\begin{enumerate}[label=\alph*)]
  \item \textbf{Confirmar que todos los requisitos del proyecto se hayan cumplido}
  \item Recolectar información histórica de proyectos anteriores
  \item Obtener la aprobación formal de los planes de gestión
  \item Trabajar con el cliente para determinar los criterios de aceptación
\end{enumerate}

\subsection*{Pregunta 89}
¿Qué debe hacerse con los riesgos incluidos en la lista de supérvisión? a. Documentarlos y regresar a ellos durante el seguimiento y control del proyecto

\begin{enumerate}[label=\alph*)]
  \item \textbf{Documentarlos y apartarlos, pues ya están considerados dentro de tus planes de contingencia}
  \item Documentarlos y entregárselos al cliente
\end{enumerate}

\subsection*{Pregunta 90}
¿Qué es un programa? Arona Ss: úbra 100

\begin{enumerate}[label=\alph*)]
  \item Una regulación gubernamental
  \item \textbf{Conjunto de proyectos con un objetivo común que gestionados de forma conjunta se obtienen beneficios.}
  \item Una iniciativa empezada por la dirección
  \item Un grupo de proyectos sin relación entre sí que se dirigen de forma coordinada
\end{enumerate}

\subsection*{Pregunta 91}
¿Qué método de estimación tiende a ser el más costoso para crear un estimado de costo del proyecto?

\begin{enumerate}[label=\alph*)]
  \item 50/50
  \item Paramétrico
  \item \textbf{Ascendente}
  \item Por analogía
\end{enumerate}

\subsection*{Pregunta 92}
Se acaba de designar un Director de Proyectos a un equipo que proviene de Brasil, Japón, Suecia y España. ¿Cuál es la mejor herramienta para tener éxito?:

\begin{enumerate}[label=\alph*)]
  \item Videoconferencias
  \item La matriz de asignación de responsabilidades RAM
  \item \textbf{La comunicación y habilidades personales bien desarrolladas}
  \item Comunicación a través de la EDP
\end{enumerate}

\subsection*{Pregunta 93}
Se consideran interesados en un proyecto...

\begin{enumerate}[label=\alph*)]
  \item La banca, la administración y el ejecutivo
  \item \textbf{Clientes, patrocinadores, la organización ejecutante y/o el público}
  \item Clientes, patrocinadores y la cultura de la organización ejecutante
  \item Las multinacionales, organismos internacionales y entidades sin ánimo de lucro
\end{enumerate}

\subsection*{Pregunta 94}
Se ha unido a la Oficina de Dirección de Proyecto después de cinco años trabajando en proyectos. Una de las herramientas que quiere introducir en su compañía es la necesidad de crear y utilizar la EDP. Alguno de los Directores de Proyecto están enfadados porque les está solicitando que hagan "trabajo extra". ¿Cuál de las siguientes afirmaciones podría decirles para convencerles del uso de esta herramienta? (mejor afirmación)

\begin{enumerate}[label=\alph*)]
  \item Decirles que no es necesaria su utilización
  \item Decirles que solo deben de hacer uso de la herramienta cuando el proyecto tenga un contrato
  \item Decirles que es la única forma de identificar los riesgos del proyecto
  \item Decirles que el uso de esta herramienta les prevendrá de que el trabajo sea mayor del necesario
\end{enumerate}

\subsection*{Pregunta 95}
Una señal de aviso indicando que un riesgo identificado puede estar aproximándose o ser inminente se le conoce como: 00

\begin{enumerate}[label=\alph*)]
  \item Reserva de contingencia (Contingency reserve)
  \item Indicador (Handle)
  \item Reserva de gestión (Management reserve)
  \item Disparador (trigger)
\end{enumerate}

\subsection*{Pregunta 96}
Si determina que uno de sus interesados del proyecto tiene un nivel de poder alto y poco interés o implicación en su proyecto, ¿qué estrategia debería implementar para gestionar a dicho interesado?

\begin{enumerate}[label=\alph*)]
  \item Monitorizar
  \item Mantener satisfecho
  \item Mantener informado
  \item Gestionar de cerca
\end{enumerate}

\subsection*{Pregunta 97}
Si el valor ganado (EV) es 1000, el costo real (AC) es 2000, el valor planificado 750. ¿Cuál es la variación del costo (CV)?

\begin{enumerate}[label=\alph*)]
  \item 250
  \item \textbf{-1000}
  \item 1250
  \item 1000
\end{enumerate}

\subsection*{Pregunta 98}
Si un director de proyecto está utilizando la estructura de desglose del trabajo como la mostrada en el gráfico, entonces la forma de la EDT es más probable que sea:

\begin{enumerate}[label=\alph*)]
  \item El enfoque de ensamblaje ascendente.
  \item El enfoque de entregables principales.
  \item \textbf{El enfoque de subproyectos.}
  \item El enfoque de fases de ciclo de vida.
\end{enumerate}

\subsection*{Pregunta 99}
Si usted fuese accionista de la empresa LLL y le presentasen los resultados de dos proyectos que se están ejecutando actualmente, ¿en cuál de los dos desearía haber invertido? Proyecto 1: Índice de desempeño del cronograma: 1,2; Índice de desempeño del costo: 1,2. Proyecto 2: 1: Índice de desempeño del cronograma: 0,1; Índice de desempeño del costo: 0,8

\begin{enumerate}[label=\alph*)]
  \item El proyecto 2
  \item El proyecto 1
  \item Ambos proyectos parecen rentables
  \item No me gustaría haber invertido en ninguno de los dos
\end{enumerate}

\subsection*{Pregunta 100}
Todo el trabajo técnico ha sido comp letado en el proyecto. ¿Qué falta por hacer?

\begin{enumerate}[label=\alph*)]
  \item Validar el Alcance
  \item Planificar la Respuesta a los Riesgos
  \item Crear un plan de gestión de personal
  \item \textbf{Completar las lecciones aprendidas 151}
\end{enumerate}

\subsection*{Pregunta 101}
Un director de proyecto está preocupado por la posibilidad de que un miembro clave del equipo abandone la organización y está pensando en preparar un plan de mitigación. ¿Cuál de las siguientes estrategias ayudará a mitigar el impacto de este evento?

\begin{enumerate}[label=\alph*)]
  \item Prometer al miembro del equipo un gran bono al final del proyecto relacionado con su rendimiento
  \item \textbf{Desarrollar un plan específico para involucrar activamente y mantener motivado al miembro del equipo}
  \item Desarrollar un plan de contingencia en caso de que el miembro del equipo se vaya
  \item Alinear las aspiraciones del miembro del equipo con los objetivos del proyecto
\end{enumerate}

\subsection*{Pregunta 102}
Un director de proyecto pierde parte de su equipo a la mitad de un proyecto. Se forma un nuevo equipo para sustituir al anterior. Como director de proyecto, ¿cuál es el primer tema a tratar en la reunión inicial con el nuevo equipo?

\begin{enumerate}[label=\alph*)]
  \item Identificar los roles y responsabilidades de los miembros del equipo.
  \item Discutir las estimaciones de costos.
  \item Revise el cronograma detallado.
  \item Enfatizar su autoridad.
\end{enumerate}

\subsection*{Pregunta 103}
Un miembro del equipo no está rindiendo bien porque carece de experiencia en el trabajo específico asignado. Na hay nadie mejor disponible. ¿Cuál es la mejor estrategia para el Director

\begin{enumerate}[label=\alph*)]
  \item de Proyecto? did
  \item \textbf{Obtener un nuevo recurso más cualificado (9) b. Proporcionar capacitación al miembro del equipo}
  \item Establecer incentivos al miembro del equipo al finalizar el proyecto
  \item Asignar parte de la reserva del cronograma del proyecto
\end{enumerate}

\subsection*{Pregunta 104}
Un proyecto ha enfrentado dificultades importantes en la calidad de sus entregables . La gerencia ahora afirma que la restricción más importante del proyecto es la calidad . Si sucediera otro prob lema con la calidad, ¿qué es lo MEJOR que podría hacer el director del proyecto?

\begin{enumerate}[label=\alph*)]
  \item Arreg lar el problema lo más pronto posible .
  \item \textbf{Permitir que el cronograma se modifique a través de una reducción de los costos .}
  \item Permitir que los costos incrementen arreglando la causa raíz del problem a.
  \item Permitir que el riesgo aumente recortando los costos.
\end{enumerate}

\subsection*{Pregunta 105}
Usted está a cargo de un equipo global de proyecto. Cada vez que se presenta un problema, su equipo decide enviar correos electrónicos a todos los que se incluyen en la lista de distribución del proyecto. Esta situación está fuera de control y hay mucha información en manos de personas que no necesitan conocerla o que no han mostrado interés en el proyecto. ¿Qué deberá hacer para corregir esta situación?

\begin{enumerate}[label=\alph*)]
  \item Pedirle a su equipo que investigue este problema y que le presente una recomendación.
  \item \textbf{Si ya se establecieron las reglas, hacer que el equipo les preste atención. De lo contrario, actualizar el plan de comunicación con los lineamientos para comunicación a través de correo electrónico.}
  \item Enviar un correo electrónico a los miembros del equipo y hacer que detengan esta práctica relacionada con la distribución de información.
  \item Publicar y distribuir los lineamientos y pedir a los miembros del equipo que lean la sección correspondiente a los correos electrónicos.
\end{enumerate}

\subsection*{Pregunta 106}
Usted está dirigiendo un proyecto para actualizar una red de cable existente, con el fin de evitar interrupciones. Parte de la responsabilidad de su equipo, es determinar la causa de estos cortes. Hay muchas opiniones y algunas buenas ideas, pero es difícil para su equipo para identificar los problemas más importantes. El equipo está empezando a sentirse frustrado y usted se siente fuera de control. ¿Qué puede hacer para ayudar al equipo a enfocarse?

\begin{enumerate}[label=\alph*)]
  \item Dirigir al equipo en el uso de un diagrama de causa y efecto (espina de pescado) para facilitar un enfoque estructurado para analizar las causas profundas de los cortes en la red.
  \item Desarrollar una sesión de tormenta de ideas para que en base a la experiencia y conocimientos del equipo, determinen por qué está fallando la red.
  \item Recopilar datos sobre las quejas del cliente para ayudar al equipo a enfocarse en los grandes problemas.
  \item Utilizar un histograma para medir y comprender mejor la distribución del tiempo promedio para reparar interrupciones del sistema.
\end{enumerate}

\subsection*{Pregunta 107}
Usted está preparando un acta para su proyecto de desarrollo de software. Su proyecto está organizado por fases. Las fases típicas para un proyecto de software incluyen Conceptualización, Planificación, Construcción, Pruebas, e Implementación. Los grupos de procesos de la dirección de proyectos:

\begin{enumerate}[label=\alph*)]
  \item Deben utilizarse en lugar de estas fases.
  \item No se necesitan debido a que las fases del proyecto son suficientes para la planificación.
  \item Sólo deben realizarse al término del proyecto para indicar la entrega.
  \item \textbf{Deberán realizarse dentro de cada fase al igual que al final del proyecto.}
\end{enumerate}

\subsection*{Pregunta 108}
1 n e s Procesos de Dirección de Proyectos La primera fase de tu proyecto ha llegado a su fin. ¿Qué es lo MÁS importante que debes asegurarte de hacer ANTES de pasar a la siguiente fase?

\begin{enumerate}[label=\alph*)]
  \item Verificar que los recursos estén disponibles para la siguiente fase.
  \item Revisar el progreso del proyecto en relación a sus líneas base.
  \item Confirmar que la fase haya cumplido con sus objetivos y lograr la aceptación formal de sus entregables.
  \item \textbf{Recomendar acción correctiva con el propósito de alinear los resultados del proyecto con las expectativas del proyecto .}
\end{enumerate}

\subsection*{Pregunta 109}
1 R E s Procesos de Dirección de Proyectos Las restricciones al cronograma del proyecto de alto nivel acaban de ser determinadas . ¿En qué grupo de procesos de la dirección de proyectos te encuentras?

\begin{enumerate}[label=\alph*)]
  \item Iniciación
  \item Planificación
  \item Ejecución
  \item \textbf{Seguimiento y control}
\end{enumerate}

\subsection*{Pregunta 110}
¿A cuál de los siguientes procesos es una entrada el registro de los interesados?

\begin{enumerate}[label=\alph*)]
  \item Planificar la Gestión de los Riesgos y Recopilar los Requisitos
  \item \textbf{Realizar el Control Integrado de Cambios y Planificar la Gestión de las Comunicaciones}
  \item Planificar la Gestión de la Calidad y Realizar el Aseguramiento de Calidad
  \item Identificar los Riesgos y Desarrollar el Acta de Constitución del Proyecto
\end{enumerate}

\subsection*{Pregunta 111}
A pesar de que los interesados pensaban que había suficiente dinero en el presupuesto, a la mitad del proyecto el índice de desempeño del costo (CPI) es de 0,7. Con el propósito de determinar la causa primordial, varios interesados hacen una auditoría del proyecto y descubren que el presupuesto de costo del proyecto se estimó de forma análoga. A pesar de que la suma de los estimados de las actividades da como resultado el estimado del proyecto, los interesados creen que algo faltó debido al modo en que se completó el estimado. ¿Cuál de las siguientes opciones describe lo que hizo falta?

\begin{enumerate}[label=\alph*)]
  \item Los costos estimados deben usarse para medir el CPI.
  \item Se debe usar el SPI, no el CPI.
  \item Debió usarse la estimación ascendente.
  \item \textbf{La historia pasada no se tuvo en cuenta. 284}
\end{enumerate}

\subsection*{Pregunta 112}
A qué personas no afecta el Código de Ética y Conducta Profesional del PMI®

\begin{enumerate}[label=\alph*)]
  \item a los que no siendo miembros han enviado una solicitud para iniciar un proceso de certificación.
  \item a los que poseen una certificación del PMI®
  \item a todas las partes interesadas de un proyecto liderado por un PMP®
  \item a los que no son miembros pero colaboran con carácter de voluntarios Respuestas razonadas
\end{enumerate}

\subsection*{Pregunta 113}
¿A qué se refiere la siguiente definición? "El punto en donde los beneficios o ganancias a ser recibidos por mejorar la calidad igualan el costo incremental para obtener esa calidad ".

\begin{enumerate}[label=\alph*)]
  \item Análisis de control de calidad
  \item Análisis marginal
  \item \textbf{Análisis de calidad estándar}
  \item Análisis de conformidad
\end{enumerate}

\subsection*{Pregunta 114}
Acabas de cerrar el grupo de procesos de inicio de un pequeño proyecto con lo que, entre otras cosas tienes el acta de constitución y la lista de interesados de los cuales uno te pide que le digas la fecha para la cual podrá lanzar el producto al mercado. ¿Qué le dirías?

\begin{enumerate}[label=\alph*)]
  \item Una primera estimación del plazo está incluida en el acta de constitución
  \item Es imposible dar una estimación antes de que se cree el plan de dirección del proyecto.
  \item El cronograma no estará finalizado hasta concluir el grupo de procesos de planificación
  \item Al ser un proyecto pequeño el plan de dirección del proyecto no incluirá el cronograma
\end{enumerate}

\subsection*{Pregunta 115}
Acabas de incorporarte a una empresa como Project manager donde trabajan con una estructura organizativa por proyectos. todos los miembros de tu equipo trabajan en la misma oficina donde tienes gran independencia y autoridad. Todas las siguientes son ventajas de una organización por proyectos excepto una ¿Cuál?:

\begin{enumerate}[label=\alph*)]
  \item Lealtad al proyecto
  \item Los miembros del equipo tiene su puesto fijo al acabar el proyecto
  \item La comunicación es más efectiva
  \item La organización del proyecto es más eficiente
\end{enumerate}

\subsection*{Pregunta 116}
Acabas de ser designado como director de proyectos de un proyecto de telecomunicaciones grande . Este proyecto de un año de duración está cerca de la mitad de su realización. El equipo del proyecto está compuesto por 5 vendedores y 20 empleados de tu compañía. Quieres entender quién es responsable de cada cosa en el proyecto. ¿Dónde encontrarías esta información?

\begin{enumerate}[label=\alph*)]
  \item Matriz de asignación de responsabilidades
  \item Histograma de recursos
  \item Diagrama de barras
  \item \textbf{Organigrama del proyecto}
\end{enumerate}

\subsection*{Pregunta 117}
Una actividad tiene un inicio temprano (ES) del día 3, un inicio tardío (LS) del día 13, una finalización temprana (EF) del día 9 y una finalización tardía (LF) del día 19. La actividad :

\begin{enumerate}[label=\alph*)]
  \item Está en la ruta crítica.
  \item Tiene un retraso .
  \item Está avanzando bien .
  \item No está en la ruta crítica.
\end{enumerate}

\subsection*{Pregunta 118}
Además del plan para la dirección del proyecto, el registro de cambios y los Activos de los procesos de la organización, ¿Qué otras tres Entradas son necesarias para gestionar adecuadamente las expectativas de los interesados?

\begin{enumerate}[label=\alph*)]
  \item La estrategia de gestión de los interesados, el registro de interesados y las solicitudes de cambios
  \item El plan de comunicación, la estrategia de gestión de interesados y las herramientas de distribución de la información
  \item El registro de interesados, la estrategia de gestión de las expectativas de los interesados y el registro de incidentes.
  \item Los métodos de comunicación, las plantillas e instrucciones de comunicación y el plan de gestión de las expectativas de los interesados
\end{enumerate}

\subsection*{Pregunta 119}
Al definir e implementar el proceso de planificación de riesgo, es importante reconocer la actitud de los principales stakeholders hacia el riesgo. Alguien que quiere mantener un enfoque equilibrado es una persona que:

\begin{enumerate}[label=\alph*)]
  \item Persigue el riesgo o busca el riesgo
  \item \textbf{Tolerar el riesgo o es neutral ante el riesgo}
  \item Prefiere evitar riesgos o es adverso al riesgo
  \item No entiende el riesgo o ignora el riesgo
\end{enumerate}

\subsection*{Pregunta 120}
Al Project manager se le ha asignado un proyecto complejo y completamente nuevo en la empresa en el que aparecen como patrocinadores tres unidades de negocio que están especialmente interesadas en el resultado del proyecto. Aunque el nuevo proyecto ha sido debatido y valorada su necesidad e importancia, ¿cuál será el principal reto en el proyecto?

\begin{enumerate}[label=\alph*)]
  \item Los objetivos no concluyentes entre patrocinadores
  \item El equipo de proyecto y puesta en marcha
  \item La estructura de desglose de tareas
  \item El control integrado de cambios
\end{enumerate}

\subsection*{Pregunta 121}
Al escribir este libro, estoy utilizando casos prácticos procedentes de las empresas en las que he trabajado y que cito en el libro. Si quien va a editar mi libro es un PMP® ¿Qué debiera hacer?

\begin{enumerate}[label=\alph*)]
  \item Solicitar aprobación explícita y escrita de que las empresas aceptan que sus nombres y sus casos aparezcan.
  \item Como editor responsabilizarse de que la información sea veraz
  \item Eliminar los casos de estudio del libro
  \item Aceptar unan vez corregidos los posibles errores que pudiera haber
\end{enumerate}

\subsection*{Pregunta 122}
Al final de un proyecto , un director de proyectos determina que el proyecto ha agregado cuatro áreas de funcionalidad y tres áreas de desempeño . El cliente ha expresado satisfacción con el proyecto . ¿Qué significa esto en términos del éxito del proyecto?

\begin{enumerate}[label=\alph*)]
  \item El proyecto tuvo un éxito rotundo .
  \item El proyecto no fue exitoso porque estaba "bañado en oro" (añadir funcionalidad extra).
  \item El proyecto no fue exitoso porque el hecho de que el cliente estaba contento significa que hubiera pagado más por el trabajo .
  \item \textbf{El proyecto fue exitoso porque el equipo tuvo una oportunidad de aprender nuevas áreas de funcionalidad y el cliente estuvo satisfecho.}
\end{enumerate}

\subsection*{Pregunta 123}
Al final del proyecto, el Project manager anuncia que el proyecto del nuevo juego informático a entregar al cliente ha añadido cuatro funcionalidades nuevas y tres mejoras en su funcionamiento operativo.

\begin{enumerate}[label=\alph*)]
  \item El proyecto es un éxito absoluto ya que ofrece más de lo que esperaba el cliente
  \item El proyecto no ha sido un éxito ya que se ha hecho trabajo que no estaba en el alcance acordado con el cliente
  \item El proyecto ha sido un fracaso ya que si el cliente está contento significa que se le podía haber vendido más caro
  \item El proyecto ha sido un éxito ya que ha permitido al equipo y la compañía mejorar sus productos con vistas a nuevos clientes
\end{enumerate}

\subsection*{Pregunta 124}
Al intentar concluir el proyecto más rápido, el director del proyecto examina los costos asociados con la compresión de cada actividad. El MEJOR enfoque a la compresión también incluiría examinar:

\begin{enumerate}[label=\alph*)]
  \item El impacto de los riesgos de la compresión de cada actividad.
  \item La opinión del cliente sobre qué actividades comprimir.
  \item La opinión del jefe sobre qué actividades comprimir y en qué orden.
  \item \textbf{Fase del ciclo de vida del proyecto en la que la actividad está programada para ocurrir.}
\end{enumerate}

\subsection*{Pregunta 125}
Al trabajar en tu proyecto tienes algunas actividades que pueden ir desarrollándose en paralelo pero deben finalizar a la vez. ¿Qué tipo de método de secuenciación de actividades se utiliza en estos casos?

\begin{enumerate}[label=\alph*)]
  \item Método de diagramación por precedencia, PDM
  \item Método del Camino crítico (CPM)
  \item IERT (Iterative Graphical Evaluation and Review Technique)
  \item Método de diagramación de Flechas (ADM)
\end{enumerate}

\subsection*{Pregunta 126}
Algunos directores de proyecto rechazan promociones para puestos funcionales o de linea, incluidos los puestos de alto nivel, debido a su deseo de permanecer como directores de proyectos. En este caso, los responsables del proyecto creen que la gestión del proyecto satisface sus necesidades de \_\_\_\_\_\_\_\_\_\_

\begin{enumerate}[label=\alph*)]
  \item Autoestima.
  \item Seguridad.
  \item \textbf{Auto-realización.}
  \item Pertenencia o sociales.
\end{enumerate}

\subsection*{Pregunta 127}
Calcula el Valor Monetario Esperado (Expected Monetary Value, EVM) para el riesgo en el siguiente caso en el que la probabilidad de que el evento de riesgo ocurra es del 60\%. Si ocurre, existe una probabilidad del 60\% de tener un impacto positivo de 1.000.000 euros en la cuenta de resultados pero también una probabilidad del 40\% de que tenga un impacto negativo de -500.000. 60\% +1.000.000 euros 60\% OCURRE 40\% --500.000 euros Evento de riesgo 40\% NO OCURRE

\begin{enumerate}[label=\alph*)]
  \item 300.000
  \item 600.000
  \item -200.000
  \item 400.000
\end{enumerate}

\subsection*{Pregunta 128}
Cambiando de ejemplo, si el valor esperado de una actividad fuera de 100 días con una desviación estándar de 5 días, cuál de las siguientes respuestas es correcta:

\begin{enumerate}[label=\alph*)]
  \item Hay una probabilidad del 95\% de que la duración esté entre 115 y 85
  \item Hay una probabilidad del 68\% de que esté entre 105 y 95
  \item Hay una probabilidad del 99\% de que esté entre 120 y 80
  \item Ninguna de las anteriores
\end{enumerate}

\subsection*{Pregunta 129}
Una compañía está haciendo un esfuerzo por mejorar el rendimiento de su proyecto y crear registros históricos de proyectos pasados . ¿Cuál es la MEJOR manera de lograrlo?

\begin{enumerate}[label=\alph*)]
  \item Crear planes para la dirección del proyecto.
  \item Crear lecciones aprendidas .
  \item Crear diagramas de red.
  \item Crear informes del estado .
\end{enumerate}

\subsection*{Pregunta 130}
Una compañía organizada jerárquicamente en relación a una compañía organizada por proyectos tiende a:

\begin{enumerate}[label=\alph*)]
  \item Ser más productiva en sus reuniones
  \item Dedicar menos tiempo a reuniones
  \item Tener una mayor productividad
  \item Ser más eficiente en la gestión del tiempo de trabajo
\end{enumerate}

\subsection*{Pregunta 131}
Con el objetivo de que puedas optar a un puesto muy atractivo y para el que te crees plenamente capaz, preparas un buen curriculum, bien maquetado, con una presentación muy cuidada y aumentas en un par de años tu formación, cambias ligeramente el nombre de los cargos ostentados para que suene más importante e incluyes proyectos en los que no has participado directamente. ¿En qué aspecto de la integridad individual del Project manager estás fallando?

\begin{enumerate}[label=\alph*)]
  \item En ninguno ya que realmente adornando y haciendo más atractivo mi curriculum.
  \item No estoy siendo honesto
  \item No estoy siendo responsable
  \item No estoy siendo justo con la compañía que aspiro a que me contrate
\end{enumerate}

\subsection*{Pregunta 132}
¿Cuál de la siguientes herramientas o técnicas corresponde al proceso de cierre de las adquisiciones

\begin{enumerate}[label=\alph*)]
  \item Documentación del contrato
  \item Negociaciones del contrato
  \item Auditorías de la adquisición
  \item Sistemas de pago
\end{enumerate}

\subsection*{Pregunta 133}
Cuál de las afirmaciones siguientes no es correcta en relación al análisis de riesgos?

\begin{enumerate}[label=\alph*)]
  \item El análisis cuantitativo se realiza si y solo si se ha realizado ya el cualitativo
  \item el análisis cualitativo normalmente se realiza antes del cuantitativo
  \item el objetivo de los análisis cualitativo y cuantitativo es priorizar los riesgos identificados
  \item el registro de riesgos es una entrada para los análisis cuantitativo y cualitativo.
\end{enumerate}

\subsection*{Pregunta 134}
¿Cuál de las afirmaciones siguientes relativas a la Holgura no es correcta?

\begin{enumerate}[label=\alph*)]
  \item La Holgura de las tareas que están en el camino crítico es por definición cero
  \item La Holgura es el tiempo que marcamos de adelanto en una relación Final a inicio para no tener problemas en el plazo total.
  \item La Holgura es el tiempo máximo que se puede retrasar una tarea sin afectar al camino crítico
  \item La Holgura puede ser negativa
\end{enumerate}

\subsection*{Pregunta 135}
Cuál de las respuestas presenta el orden cronológico normalmente utilizado para la planificación de riesgos

\begin{enumerate}[label=\alph*)]
  \item identificación de riesgos, análisis cualitativo, análisis cuantitativo, plan de respuesta a riesgos, plan de gestión de riesgos .
  \item plan de gestión de riesgos, análisis cuantitativo, análisis cualitativo, identificación de riesgos, plan de respuesta a riesgos.
  \item plan de gestión de riesgos , identificación de riesgos, análisis cuantitativo, análisis cualitativo, plan de respuesta a riesgos.
  \item plan de gestión de riesgos, identificación de riesgos, análisis cualitativo, análisis cuantitativo, plan de respuesta a riesgos.
\end{enumerate}

\subsection*{Pregunta 136}
¿Cuál de las siguientes acciones es la MÁS adecuada durante el grupo de procesos de iniciación?

\begin{enumerate}[label=\alph*)]
  \item Crear una descripción detallada de los entregables del proyecto.
  \item Familiarizarse con la cultura de la compañía y su estructura en relación con el proyecto.
  \item \textbf{Identificar la causa raíz de los problemas .}
  \item Asegurarse de que todos los procesos de dirección de proyectos hayan sido completados .
\end{enumerate}

\subsection*{Pregunta 137}
¿Cuál de las siguientes actividades de realiza primero?

\begin{enumerate}[label=\alph*)]
  \item Crear el acta de constitución
  \item Crear la Estructura de desglose de trabajo, EDT
  \item Crear el Enunciado del Alcance del Proyecto
  \item Crear los requerimientos del proyecto.
\end{enumerate}

\subsection*{Pregunta 138}
¿Cuál de las siguientes actividades no forma parte de la definición de los requisitos de compra?

\begin{enumerate}[label=\alph*)]
  \item Autorización o aprobación de los requisitos
  \item Decisión de contratar o realizar el trabajo internamente
  \item Selección del proveedor que garantice el cumplimiento de los requisitos
  \item Elaboración del Enunciado del Trabajo y Especificaciones
\end{enumerate}

\subsection*{Pregunta 139}
¿Cuál de las siguientes afirmaciones define mejor el papel del Project manager en relación a la gestión integrada de cambios en un proyecto?

\begin{enumerate}[label=\alph*)]
  \item Hacer caso al cliente y poner en práctica lo que éste indique
  \item Monitorizar y controlar los cambios para asegurar que benefician al proyecto
  \item Actuar sobre los factores que afectan a los cambios
  \item Evitar los cambios en el proyecto
\end{enumerate}

\subsection*{Pregunta 140}
¿Cuál de las siguientes afirmaciones en relación al control del alcance es correcta? :

\begin{enumerate}[label=\alph*)]
  \item El control del alcance es una actividad realizada por el cliente para confirmar que el Entregable está completo y correcto y procede la aprobación formal.
  \item En este proceso se monitoriza el estado del alcance del producto y el producto y se gestionan cambios a la línea de base del producto.
  \item Como resultado se obtiene una matriz de trazabilidad de requisitos que evite que se dejen de verificar requisitos por parte del cliente
  \item El análisis de variación es una herramienta utilizada en este proceso para evaluar el retraso en plazo y coste y sus implicaciones sobre el alcance final al cliente.
\end{enumerate}

\subsection*{Pregunta 141}
Cuál de las siguientes afirmaciones en relación al sistema de gestión de registros no es correcta?

\begin{enumerate}[label=\alph*)]
  \item Es una salida del proceso de administrar las adquisiciones
  \item Permite mantener un registro de todas las acciones, documentación, incidencias, cambios y correspondencia relativa al contrato
  \item Es un documento muy importante desde el punto de vista legal y contractual
  \item Debe mantenerse disponibles para ambas partes del contrato.
\end{enumerate}

\subsection*{Pregunta 142}
¿Cuál de las siguientes afirmaciones es claramente correcta?

\begin{enumerate}[label=\alph*)]
  \item No debe aceptar un proyecto que claramente está mal planificado y no va a entregarse en plazo si no se aceptan los cambios que proponga.
  \item Puede aceptar un proyecto sin una acta de constitución que deje claros desde el principio el alcance y los requisitos del proyecto,
  \item Por razones de superior criterio, debe aceptar que la dirección le impida planificar el trabajo antes de iniciar la fase de ejecución,
  \item En ocasiones no tendrá más remedio que comprometer al equipo del proyecto al dar concesiones que vayan en contra de la formación y desarrollo profesional del equipo,
\end{enumerate}

\subsection*{Pregunta 143}
¿Cuál de las siguientes afirmaciones es incorrecta en relación a las Conferencias de Licitadores?

\begin{enumerate}[label=\alph*)]
  \item Trata de que todos los licitantes compiten en una manera justa y equilibrada
  \item Los licitadores hacen preguntas sobre el contrato
  \item A cada licitador se le manda un acta con la respuesta a las preguntas que ha realizado
  \item Se convocan una vez que la compañía ha preparado una oferta o propuesta
\end{enumerate}

\subsection*{Pregunta 144}
¿Cuál de las siguientes afirmaciones no es cierta en relación a una reserva de contingencia?

\begin{enumerate}[label=\alph*)]
  \item la suma total depende de la tolerancia al riesgo definida por la dirección
  \item son cantidades planificadas para atender situaciones difíciles de predecir
  \item se aplican únicamente a los costes no a los plazos del proyecto
  \item son costes previstos en el presupuesto para riesgos identificados y con plan de respuesta que se activará en cuanto el evento de riesgo ocurra
\end{enumerate}

\subsection*{Pregunta 145}
¿Cuál de las siguientes afirmaciones no es correcta en relación a la calidad:

\begin{enumerate}[label=\alph*)]
  \item es clave saber cuáles son los niveles que las partes interesadas y el usuario final van a considerar aceptables.
  \item El Project manager gestiona la calidad estableciendo unos objetivos y midiendo unos indicadores que indican si está cumpliéndolos.
  \item El proyecto debe tener como objetivo productos de alto grado, de alto valor, que impacten al cliente
  \item tiene mucho más que ver con la prevención que con la inspección,
\end{enumerate}

\subsection*{Pregunta 146}
¿Cuál de las siguientes afirmaciones no es correcta en relación al Método de la ruta crítica?

\begin{enumerate}[label=\alph*)]
  \item Permite identificar el plazo final del proyecto
  \item Utiliza datos de disponibilidad de recursos para afinar el camino crítico
  \item Identifica las Holguras que tienen las diferentes actividades
  \item Permite calcular el día más pronto que puede empezar una tarea
\end{enumerate}

\subsection*{Pregunta 147}
¿Cuál de las siguientes afirmaciones no es correcta en relación al Plan de comunicaciones que, como Project manager, vas a obtener como resultado de este proceso?

\begin{enumerate}[label=\alph*)]
  \item Trata de que todas las partes comprendan la razón y objetivos de las reuniones, videoconferencias, informes que el Project manager envíe
  \item Indica qué información enviará el Project manager a quién y con qué frecuencia
  \item Este plan debe incluir la lista de interesados , la matriz de asignación de responsabilidades de cada uno y los modelos de comunicación empleados
  \item Debe incluir el responsable de comunicar la información, los formatos específicos y medios de transmisión y una instrucción para actualizar el plan
\end{enumerate}

\subsection*{Pregunta 148}
¿Cuál de las siguientes afirmaciones no es correcta sobre el solapamiento de fases?

\begin{enumerate}[label=\alph*)]
  \item Cada fase debe desarrollarse en los cinco grupos de procesos
  \item Hay un riesgo añadido al proyecto debido a los errores a corregir
  \item Retrasará el proyecto si la fase posterior no puede empezar si no acaba la anterior
  \item Cada fase la hace un equipo diferente
\end{enumerate}

\subsection*{Pregunta 149}
¿Cuál de las siguientes afirmaciones refleja mejor la diferencia entre planificar las adquisiciones y planificar el contrato?

\begin{enumerate}[label=\alph*)]
  \item El plan de gestión de las adquisiciones del proyecto implica documentar los requerimientos del producto y los potenciales proveedores mientras que el planificar el contrato implica obtener cotizaciones, ofertas, propuestas
  \item El plan de gestión de las adquisiciones del proyecto implica documentar los requerimientos del producto e identificar potenciales proveedores mientras que el planificar el contrato implica definir qué se va a comprar o subcontratar
  \item El plan de gestión de las adquisiciones del proyecto implica identificar qué es necesario comprar o subcontratar mientras que la planificación del contrato implica actividades como documentar los requisitos del producto, identificar potenciales vendedores,…
  \item Ninguna de las anteriores es correcta
\end{enumerate}

\subsection*{Pregunta 150}
¿Cuál de las siguientes afirmaciones refleja mejor la teoría de los factores higiénicos?

\begin{enumerate}[label=\alph*)]
  \item hay aspectos del ambiente de trabajo que si están no se aprecian y si faltan son si están, motivan fuertemente.
  \item El Project manager necesita dar a su equipo unas expectativas de reconocimiento y recompensa para motivarlos pero solamente es efectivo si el logro es alcanzable.
  \item las personas se motivan por los logros conseguidos, el poder adquirido y la cercanía y buena relación con el equipo y la dirección r
  \item la motivación de las personas tiene que ver con su nivel de satisfacción de sus expectativas vitales
\end{enumerate}

\subsection*{Pregunta 151}
¿Cuál de las siguientes afirmaciones sobre el Earned Value Management no es cierta?

\begin{enumerate}[label=\alph*)]
  \item facilitar el reporte mensual del desempeño del proyecto
  \item valora el avance del programa y el cumplimiento de presupuestos basándose en los datos hasta la fecha
  \item permite conclusiones como que el proyecto va adelantado pero con un sobrecoste importante
  \item facilita el seguimiento al no necesitar realizar un control exhaustivo de horas, porcentajes de avance ni facturas
\end{enumerate}

\subsection*{Pregunta 152}
¿Cuál de las siguientes afirmaciones sobre la responsabilidad del Project manager en la gestión de riesgos no es cierta?

\begin{enumerate}[label=\alph*)]
  \item Inicia y lidera el proceso de gestión de riesgos
  \item Identifica y elabora el registro de riesgos con el equipo de proyecto
  \item Da directrices al equipo en procesos y herramientas a utilizar
  \item Elaborar el plan de respuestas a los riesgos
\end{enumerate}

\subsection*{Pregunta 153}
¿Cuál de las siguientes afirmaciones sobre los diagramas de comportamiento no es correcta?:

\begin{enumerate}[label=\alph*)]
  \item Son similares a los de control pero sin mostrar límites de rechazo.
  \item Tratan de identificar tendencias ya cíclicas, ya constantes, ya monótonas crecientes.
  \item Son diagramas que ayudan a ver cómo se relacionan las diferentes fases entre sí y las muestras bajadas en las prestaciones –valores inusualmente bajos durante un período-
  \item que lleven a estudiar que variables pudieron cambiar esos días en los que la calidad fue sensiblemente menor.
\end{enumerate}

\subsection*{Pregunta 154}
¿Cuál de las siguientes afirmaciones sobre los riesgos es correcta?

\begin{enumerate}[label=\alph*)]
  \item Un riesgo es un evento con una probabilidad alta de que ocurra
  \item Un riesgo siempre afectará negativamente al proyecto
  \item Se pueden categorizar en una Estructura de desglose de riesgos
  \item las contingencias son riesgos con los que se cuenta que ocurran
\end{enumerate}

\subsection*{Pregunta 155}
¿Cuál de las siguientes afirmaciones son ciertas en relación al CPI (índice de desempeño del coste)?

\begin{enumerate}[label=\alph*)]
  \item Un CPI menor que uno indica que el desempeño es mejor que el esperado
  \item Un CPI mayor que uno indica que el desempeño es mejor que el esperado
  \item Un CPI negativo indica que el desempeño del coste es mejor que el esperado
  \item Un CPI positivo indica que el desempeño en costes es mejor que el esperado
\end{enumerate}

\subsection*{Pregunta 156}
Cuál de las siguientes afirmaciones son falsas en relación al sistema de Control de cambios en el coste?

\begin{enumerate}[label=\alph*)]
  \item Informa sobre el desempeño en el coste: qué tareas han variado su presupuesto
  \item Define el nivel de aprobación necesario para autorizar los cambios
  \item Es un sistema integrado de control de cambio que se debe obligatoriamente cumplir
  \item Define procedimientos por los que la línea de base del coste puede cambiarse.
\end{enumerate}

\subsection*{Pregunta 157}
¿Cuál de las siguientes condiciones expresa mejor un ejemplo de contrato de precio fijo con incentivos por desempeño?

\begin{enumerate}[label=\alph*)]
  \item El Project manager paga 2.500.000 euros
  \item El Project manager paga 2.500.000 euros en los que están incluidos los costes más un tarifa porcentual a los mismos.
  \item El proveedor recibe los 2.000.000 correspondientes a los costes más 500.000 euros si se entrega el producto antes de final de año.
  \item El proveedor recibe 2.500.000 euros más 500.000 euros si el producto cumple unos parámetros de calidad complementarios definidos en el contrato.
\end{enumerate}

\subsection*{Pregunta 158}
¿Cuál de las siguientes describe MEJOR el análisis del producto?

\begin{enumerate}[label=\alph*)]
  \item Trabajar con el cliente para determinar la descripción del producto
  \item \textbf{Analizar matemáticamente la calidad deseada para el proyecto}
  \item Adquirir una mejor comprensión del producto del proyecto para crear el enunciado del alcance del proyecto
  \item Determinar si pueden alcanzarse los estándares de calidad del proyecto 191
\end{enumerate}

\subsection*{Pregunta 159}
¿Cuál de las siguientes describe MEJOR la relación entre la desviación estándar y los riesgos?

\begin{enumerate}[label=\alph*)]
  \item No hay relación.
  \item La desviación estándar te dice si el estimado es preciso .
  \item \textbf{La desviación estándar te dice qué tan incierto es el estimado.}
  \item La desviación estándar te dice si el estimado incluye un relleno-:- ·
\end{enumerate}

\subsection*{Pregunta 160}
Cuál de las siguientes elementos no forma parte de los resultados del proceso de realización del control de calidad:

\begin{enumerate}[label=\alph*)]
  \item Las mediciones de los controles de calidad
  \item Los cambios validados
  \item Los informes de desempeño
  \item Los Entregables validados
\end{enumerate}

\subsection*{Pregunta 161}
Cuál de las siguientes Entradas no es necesaria para iniciar al proceso de secuenciar las actividades

\begin{enumerate}[label=\alph*)]
  \item Lista de actividades y atributos de la actividad
  \item Lista de hitos que deben cumplirse
  \item Línea base del alcance del producto)
  \item Los Adelantos y retrasos entre actividades
\end{enumerate}

\subsection*{Pregunta 162}
¿Cuál de las siguientes Entradas no forma parte del proceso de Estimar los costos?

\begin{enumerate}[label=\alph*)]
  \item Línea base del alcance
  \item Plan de recursos humanos
  \item Registro de riesgos
  \item Calendario de recursos
\end{enumerate}

\subsection*{Pregunta 163}
¿Cuál de las siguientes es una característica de los procesos de dirección de proyectos?

\begin{enumerate}[label=\alph*)]
  \item Iterativo
  \item Único
  \item Innecesario
  \item \textbf{Estandarizado}
\end{enumerate}

\subsection*{Pregunta 164}
¿Cuál de las siguientes es una de las características más importantes de la técnica Delphi?

\begin{enumerate}[label=\alph*)]
  \item Extrapolación de registros históricos de proyectos anteriores
  \item Opinión experta
  \item Proceso de jerarquía analítica
  \item Enfoque ascendente
\end{enumerate}

\subsection*{Pregunta 165}
¿Cuál de las siguientes es lo MEJOR que puedes hacer cuando te solicitan que concluyas un proyecto dos días antes de lo planificado?

\begin{enumerate}[label=\alph*)]
  \item Decirle a la gerencia senior que la ruta crítica del proyecto no permite que el proyecto se finalice más temprano .
  \item Decirle a tu jefe.
  \item Reunirte con el equipo para buscar opciones para realizar una compresión o ejecución rápida de la ruta crítica.
  \item Trabajar mucho y ver cuál es el estado del proyecto al mes siguiente.
\end{enumerate}

\subsection*{Pregunta 166}
¿Cuál de las siguientes es una salida CLAVE del proceso Validar el Alcance?

\begin{enumerate}[label=\alph*)]
  \item Un plan de gestión del alcance más completo
  \item \textbf{La aceptación del cliente de los entregab les del proyecto}
  \item Estimaciones del cronograma mejoradas
  \item Un sistema de información para la dirección de proyectos mejorado
\end{enumerate}

\subsection*{Pregunta 167}
¿Cuál de las siguientes es una ventaja de una estructura organizativa enfocada a proyectos?

\begin{enumerate}[label=\alph*)]
  \item La mayor competencia de la unidad de negocio
  \item La optimización por que todo está alineado entorno al proyecto
  \item Necesitar la aprobación del gerente funcional
  \item Los miembros del equipo tienen un lugar donde ir cuando el proyecto está finalizado
\end{enumerate}

\subsection*{Pregunta 168}
¿Cuál de las siguientes explica por qué la calidad se planifica y no se inspecciona?

\begin{enumerate}[label=\alph*)]
  \item Reduce la calidad y es menos costoso.
  \item Mejora la calidad y es más costoso.
  \item \textbf{Reduce la calidad y es más costoso.}
  \item Mejora la calidad y es menos costoso.
\end{enumerate}

\subsection*{Pregunta 169}
¿Cuál de las siguientes formas de abordar un conflicto debería utilizar un Project manager?

\begin{enumerate}[label=\alph*)]
  \item Eludir el problema
  \item Forzar la solución más adecuada desde el punto de vista del proyecto
  \item Buscar una solución que aporte cierto grado de satisfacción a cada una de las partes
  \item Recoger datos, evidencias, puntos de vista y buscar una solución parcialmente aceptada.
\end{enumerate}

\subsection*{Pregunta 170}
Cuál de las siguientes funciones no realiza el Project manager en la fase de control del cronograma.

\begin{enumerate}[label=\alph*)]
  \item Consultar el plan de dirección del proyecto
  \item Aplicar el análisis de variación para identificar las variaciones en fechas
  \item Recomendar cambios para solucionar problemas de retrasos
  \item Analizar los efectos en el coste de los retrasos
\end{enumerate}

\subsection*{Pregunta 171}
¿Cuál de las siguientes GENERALMENTE se ilustra MEJOR mediante diagramas de barras que por diagramas de red?

\begin{enumerate}[label=\alph*)]
  \item Relaciones lógicas
  \item Rutas críticas
  \item \textbf{Compensaciones de recursos}
  \item Avance o estado
\end{enumerate}

\subsection*{Pregunta 172}
¿Cuál de las siguientes herramientas/ técnicas incluye recopilar y analizar sistemáticamente la información para determinar los intereses de quien deben tenerse en cuenta durante el proyecto?

\begin{enumerate}[label=\alph*)]
  \item \textbf{Análisis de los interesados}
  \item Habilidades de gestión
  \item Planificación de recursos humanos
  \item Análisis de riesgos
\end{enumerate}

\subsection*{Pregunta 173}
Cuál de las siguientes matrices relaciona la EDT con las funciones y responsabilidades de la organización:

\begin{enumerate}[label=\alph*)]
  \item Estructura de desglose de la organización, OBS
  \item Matriz de asignación de responsabilidades, RAM)
  \item Matriz RACI
  \item Matriz RBS (Estructura de desglose de recursos).
\end{enumerate}

\subsection*{Pregunta 174}
¿Cuál de las siguientes no es una forma de modificación del contrato ?

\begin{enumerate}[label=\alph*)]
  \item Término del contrato
  \item Solicitud de cambio
  \item Cambio constructivo
  \item Adenda al contrato
\end{enumerate}

\subsection*{Pregunta 175}
¿Cuál de las siguientes no es responsabilidad del Project manager en la planificación y gestión de compras?

\begin{enumerate}[label=\alph*)]
  \item Seleccionar el tipo de contrato y el documento de compras
  \item Identificar los riesgos derivados de la subcontratación del trabajo
  \item Revisar el alcance del contrato para verificar que está completo
  \item Comprender los términos del contrato y sus condiciones
\end{enumerate}

\subsection*{Pregunta 176}
Cuál de las siguientes no es responsabilidad del equipo de proyecto en relación a la gestión de riesgos?

\begin{enumerate}[label=\alph*)]
  \item Informar periódicamente al patrocinador sobre nuevos riesgos
  \item Entender y seguir el procedimiento de gestión de riesgos
  \item Ejecutar las estrategias definidas para la gestión de riesgos
  \item Registrar el estado del proceso de gestión de riesgos
\end{enumerate}

\subsection*{Pregunta 177}
Cuál de las siguientes no es un resultado del proceso de Control del Alcance

\begin{enumerate}[label=\alph*)]
  \item Medición del desempeño del trabajo
  \item Solicitudes de cambio
  \item Documentos rechazados por el cliente
  \item Actualizaciones del Plan de Dirección del proyecto
\end{enumerate}

\subsection*{Pregunta 178}
¿Cuál de las siguientes opciones DEBE ser un punto de la agenda en las reuniones del equipo?

\begin{enumerate}[label=\alph*)]
  \item Discusión de los riesgos del proyecto
  \item Estado de las actividades actuales
  \item \textbf{Identificación de nuevas actividades}
  \item Revisión de los problemas del proyecto 450
\end{enumerate}

\subsection*{Pregunta 179}
¿Cuál de las siguientes opciones describe MEJOR el rol del director de proyectos como integrador?

\begin{enumerate}[label=\alph*)]
  \item Ayudar a los miembros del equipo a familiarizarse con el proyecto.
  \item Unir todas las piezas de un proyecto en un todo cohesivo.
  \item Unir todas las piezas del proyecto en un programa .
  \item Unir a todos los miembros del equipo en un todo cohesivo.
\end{enumerate}

\subsection*{Pregunta 180}
¿Cuál de las siguientes opciones describe MEJOR para qué se puede utilizar un acta de constitución del proyecto cuando se está completando el trabajo?

\begin{enumerate}[label=\alph*)]
  \item Para asegurarse de que todos los miembros del equipo sean recompensados
  \item Para ayudar a determinar si debe aprobarse un cambio en el alcance
  \item Para evaluar la efectividad del sistema de control de cambios
  \item Para garantizar que toda la documentación sobre el proyecto esté completa
\end{enumerate}

\subsection*{Pregunta 181}
¿Cuál de las siguientes opciones describe MEJOR un plan para la dirección del proyecto?

\begin{enumerate}[label=\alph*)]
  \item Una impresión de un software de dirección de proyectos
  \item Un diagrama de barras
  \item \textbf{El alcance del proyecto 145}
\end{enumerate}

\subsection*{Pregunta 182}
¿Cuál de las siguientes opciones NO es una entrada del grupo de procesos de iniciación?

\begin{enumerate}[label=\alph*)]
  \item Procesos de la compañía
  \item Cultura de la compañía
  \item EDT históricas
  \item Enunciado del alcance del proyecto
\end{enumerate}

\subsection*{Pregunta 183}
¿Cuál de las siguientes opciones representa el valor estimado del trabajo que realmente se ha llevado a cabo?

\begin{enumerate}[label=\alph*)]
  \item Valor ganado (EV)
  \item Valor planificado (PV)
  \item Costo real (AC)
  \item \textbf{Variación del costo ( CV) 282}
\end{enumerate}

\subsection*{Pregunta 184}
¿Cuál de las siguientes pueden crear el MAYOR malentendido del enunciado del alcance del proyecto?

\begin{enumerate}[label=\alph*)]
  \item Lenguaje impreciso
  \item Patrón, estructura y orden cronológico deficientes
  \item Pequeñas variaciones en el tamaño de los paquetes de trabajo o detalle de trabajo
  \item \textbf{Mucho detalle}
\end{enumerate}

\subsection*{Pregunta 185}
¿Cuál de las siguientes secuencias representa una depreciación en línea recta?

\begin{enumerate}[label=\alph*)]
  \item \$100, \$100, \$100
  \item \$100,\$120,\$140
  \item \$100,\$120,\$160
  \item \textbf{\$160, \$140, \$120 149}
\end{enumerate}

\subsection*{Pregunta 186}
¿Cuál de las siguientes técnicas de resolución de conflictos generará la solución MÁS duradera?

\begin{enumerate}[label=\alph*)]
  \item Forzar
  \item Suavizar
  \item Consentir
  \item Solución de problemas
\end{enumerate}

\subsection*{Pregunta 187}
Cuál de las siguientes técnicas o herramientas no forma parte del proceso de adquisición del equipo

\begin{enumerate}[label=\alph*)]
  \item La asignación previa de recursos
  \item La adquisición
  \item La negociación
  \item La asignación de personal al proyecto
\end{enumerate}

\subsection*{Pregunta 188}
¿Cuál de las siguientes técnicas o herramientas se utiliza para crear la Estructura de Desglose de Tareas?

\begin{enumerate}[label=\alph*)]
  \item El Juicio de expertos
  \item Los Activos de los procesos de la organización
  \item La Descomposición
  \item La planificación gradual
\end{enumerate}

\subsection*{Pregunta 189}
¿Cuál de las siguientes técnicas o herramientas tiene como objetivo identificar los puntos fuertes y débiles del proyecto y los riesgos y oportunidades externas que pueden afectarle?

\begin{enumerate}[label=\alph*)]
  \item Análisis DAFO
  \item Análisis PDCA
  \item Análisis de asunciones e hipótesis
  \item Diagramas de influencias
\end{enumerate}

\subsection*{Pregunta 190}
¿Cuál de las siguientes Técnicas y Herramientas del proceso de realizar el control de calidad implica medir o testar los resultados para determinar si cumplen con los requisitos y estándares de calidad?

\begin{enumerate}[label=\alph*)]
  \item Muestreo compuesto
  \item Inspección
  \item Cuadros de Control
  \item Supervisión
\end{enumerate}

\subsection*{Pregunta 191}
¿Cuál de los ejemplos siguientes no es un ejemplo de operación?

\begin{enumerate}[label=\alph*)]
  \item Realizar evaluaciones del desempeño anuales
  \item Crear un sistema de compra on-line de un supermercado
  \item Aplicar procesos de mejora continua en una cadena de producción
  \item Supervisar y controlar el avance de una epidemia
\end{enumerate}

\subsection*{Pregunta 192}
¿Cuál de los siguientes afirmaciones es correcta en relación a los elementos que pueden formar parte de la propiedad intelectual de la empresa:

\begin{enumerate}[label=\alph*)]
  \item elementos desarrollados por una organización que tienen valor comercial pero no son tangibles (ideas, estrategias,…),
  \item información sujeta a acuerdos de confidencialidad
  \item procesos patentados o procesos industriales, de negocio o de fabricación desarrollados por una organización para un uso específico
  \item los tres anteriores.
\end{enumerate}

\subsection*{Pregunta 193}
¿Cuál de los siguientes aspectos debe tratarse siempre en las reuniones de seguimiento del equipo de proyecto?

\begin{enumerate}[label=\alph*)]
  \item Los Incidentes registrados en el periodo a evaluar
  \item El estado de las actividades en marcha
  \item Las nuevas actividades a iniciar en el periodo siguiente
  \item El seguimiento de los riesgos del proyecto
\end{enumerate}

\subsection*{Pregunta 194}
¿Cuál de los siguientes considerarías un conflicto de intereses?

\begin{enumerate}[label=\alph*)]
  \item El gerente de la principal empresa proveedora es antiguo compañero de colegio y vecino tuyo.
  \item Tu familia es accionista minoritaria del principal proveedor de cable de tu proyecto
  \item Recibes un regalo del adjudicatario del contrato de obra civil
  \item Un suministrador te entrega información confidencial que te permite seleccionar a otro proveedor.
\end{enumerate}

\subsection*{Pregunta 195}
Cuál de los siguientes contenidos no suele contener el registro de riesgos?

\begin{enumerate}[label=\alph*)]
  \item Causas de los riesgos
  \item Lista de riesgos identificados
  \item Indicadores del riesgo
  \item Presupuesto de la respuesta
\end{enumerate}

\subsection*{Pregunta 196}
¿Cuál de los siguientes documentos de los activos de los procesos de la organización no es necesario actualizar en el proceso de Administrar las adquisiciones?

\begin{enumerate}[label=\alph*)]
  \item Auditoría de compras
  \item Correspondencia
  \item Solicitudes y fechas de pago
\end{enumerate}

\subsection*{Pregunta 197}
¿Cuál de los siguientes documentos se utiliza como base para evaluar las solicitudes de cambio y para que sirva de guía en el proceso de ejecución del proyecto?

\begin{enumerate}[label=\alph*)]
  \item La línea de base de coste y la línea de base de tiempo
  \item La línea base del alcance del producto
  \item La Estructura de Desglose de Tareas
  \item El Enunciado del Trabajo
\end{enumerate}

\subsection*{Pregunta 198}
¿Cuál de los siguientes elementos es una entrada del proceso de gestión del alcance?

\begin{enumerate}[label=\alph*)]
  \item La Estructura de Desglose de Tareas
  \item Los Cambios solicitados
  \item Los Entregables aceptados
  \item El registro de interesados
\end{enumerate}

\subsection*{Pregunta 199}
Cuál de los siguientes elementos facilita el proceso de monitorización y control de riesgos:

\begin{enumerate}[label=\alph*)]
  \item Las reuniones de estado del proyecto
  \item Las reuniones de planificación y análisis
  \item Los análisis SWOT
  \item El análisis de hipótesis
\end{enumerate}

\subsection*{Pregunta 200}
¿Cuál de los siguientes elementos no es una restricción al proyecto?

\begin{enumerate}[label=\alph*)]
  \item Calidad
  \item Escala
  \item Tiempo
  \item Riesgo
\end{enumerate}

\subsection*{Pregunta 201}
Cuál de los siguientes elementos no es una salida del proceso de estimación de recursos de las actividades

\begin{enumerate}[label=\alph*)]
  \item Requisitos de recursos de la actividad
  \item Estructura de desglose de recursos
  \item Asignación del coste unitario de cada recurso identificado
  \item Actualizaciones a los documentos del proyecto
\end{enumerate}

\subsection*{Pregunta 202}
¿Cuál de los siguientes elementos no forma parte de las Entradas al proceso de planificar la gestión de riesgos?

\begin{enumerate}[label=\alph*)]
  \item Enunciado del Trabajo
  \item Plan de Dirección del Proyecto
  \item lecciones aprendidas, como parte de los activos de la organización
  \item Cambios solicitados
\end{enumerate}

\subsection*{Pregunta 203}
¿Cuál de los siguientes elementos no forma parte del proceso de cierre del proyecto o cierre de fase?

\begin{enumerate}[label=\alph*)]
  \item El plan de Gestión de proyectos
  \item La metodología de gestión del proyecto
  \item Los entregables aceptados
  \item Las activos de procesos de la organización
\end{enumerate}

\subsection*{Pregunta 204}
¿Cuál de los siguientes elementos se utiliza para identificar las personas que pueden verse afectadas por el proyecto?

\begin{enumerate}[label=\alph*)]
  \item La lista de recursos
  \item El registro de partes interesadas
  \item Los factores ambientales de la empresa
  \item El plan de gestión del proyecto
\end{enumerate}

\subsection*{Pregunta 205}
¿Cuál de los siguientes enunciados describe mejor el proceso de Definir el Alcance?

\begin{enumerate}[label=\alph*)]
  \item Es el proceso por el cual se garantiza que los objetivos de todas las partes interesadas se van a cumplir
  \item Es el proceso que consiste en desarrollar una descripción detallada del proyecto y el producto.
  \item Crear un esquema gráfico en el que se muestran las actividades que deben estar incluidas y las que no en la EDT.
  \item En este proceso se definen todos los requerimientos del producto
\end{enumerate}

\subsection*{Pregunta 206}
¿Cuál de los siguientes enunciados describe MEJOR cómo los interesados participan en un proyecto?

\begin{enumerate}[label=\alph*)]
  \item Ayudan a determinar el cronograma , los entregables y los requisitos del proyecto .
  \item Ayudan a determinar las restricciones del proyecto y los entregables del producto.
  \item Ayudan a determinar las necesidades de los recursos y las restricciones de los recursos en el proyecto .
  \item \textbf{Aprueban el acta de constitución del proyecto, ayudan a proporcionar los supuestos y crear los planes de gestión .}
\end{enumerate}

\subsection*{Pregunta 207}
¿Cuál de los siguientes enunciados es incorrecto en relación al plan de gestión del alcance?

\begin{enumerate}[label=\alph*)]
  \item Define como verificar y controlar el alcance
  \item Describe cómo elaborar la Estructura de Desglose del Trabajo
  \item Describe cómo identificar los requisitos del trabajo
  \item Es la línea base del alcance para comparar el alcance del producto
\end{enumerate}

\subsection*{Pregunta 208}
¿cuál de los siguientes enunciados relativos al caso de negocio y a la necesidad del negocio de un proyecto son correctas?

\begin{enumerate}[label=\alph*)]
  \item El caso de negocio se incluye en el enunciado del alcance del trabajo
  \item El caso de negocio suele incluir el análisis coste-beneficio para el proyecto
  \item La necesidad del negocio es una salida del proceso de desarrollar el acta de constitución
  \item La necesidad de negocio aparece o por demandas del mercado o cambios tecnológicos Respuestas razonadas
\end{enumerate}

\subsection*{Pregunta 209}
¿cuál de los siguientes enunciados relativos al control de calidad y la verificación del alcance es correcto?

\begin{enumerate}[label=\alph*)]
  \item El control de calidad se encarga de la aceptación de productos
  \item La verificación del alcance puede realizarse paralelamente al control de calidad
  \item La verificación del alcance no debiera realizarse si el proyecto se ha anulado.
  \item La verificación del alcance y el control de calidad son procesos muy parecidos.
\end{enumerate}

\subsection*{Pregunta 210}
¿Cuál de los siguientes es lo MÁS apropiado que debe hacerse en el cierre del proyecto?

\begin{enumerate}[label=\alph*)]
  \item \textbf{Trabajar con el cliente para determinar los criterios de aceptación.}
  \item Recolectar información histórica de proyectos anteriores .
  \item Confirmar que todos los requisitos del proyecto se hayan cumplido .
  \item Obtener la aprobación formal de los planes de gestión.
\end{enumerate}

\subsection*{Pregunta 211}
¿Cuál de los siguientes es un beneficio de un estimado análogo del proyecto?

\begin{enumerate}[label=\alph*)]
  \item Estará más cerca de lo que el trabajo en realidad va a necesitar.
  \item Está basada en una comprensión detallada de lo que el trabajo necesita.
  \item \textbf{Proporciona al equipo del proyecto una comprensión de las expectativas de la gerencia.}
  \item Ayuda al director del proyecto a determinar si el proyecto cumplirá con el cronograma.
\end{enumerate}

\subsection*{Pregunta 212}
¿Cuál de los siguientes es un ejemplo de una estimación paramétrica?

\begin{enumerate}[label=\alph*)]
  \item Dólares por módulo
  \item Curva de aprendizaje
  \item Ascendente
  \item \textbf{CPM}
\end{enumerate}

\subsection*{Pregunta 213}
¿Cuál de los siguientes es un ejemplo de un evento de riesgo positivo u oportunidad? :

\begin{enumerate}[label=\alph*)]
  \item Si asignamos el desarrollo inicial de casos de prueba al mismo recurso que realizará las pruebas podemos reducir el tiempo de prueba en un 30 \%.
  \item Si mejoramos los beneficios de nuestra compañía podemos reducir la probabilidad de perder los recursos clave del proyecto.
  \item Si construimos una red redundante podemos reducir el impacto de pérdida de información del tiempo improductivo de la red.
  \item Si utilizamos una tecnología probada podemos evitar los problemas asociados con trabajar con tecnología de vanguardia.
\end{enumerate}

\subsection*{Pregunta 214}
Cuál de los siguientes factores no es un factor higiénico según Herzberg

\begin{enumerate}[label=\alph*)]
  \item salario
  \item reconocimiento del trabajo realizado
  \item buen clima de trabajo
  \item oficinas luminosas y bien diseñadas
\end{enumerate}

\subsection*{Pregunta 215}
¿Cuál de los siguientes herramientas no se utiliza para hacer informes de avance del proyecto?

\begin{enumerate}[label=\alph*)]
  \item Gráficos de barras
  \item Curvas S
  \item Las curvas de EV y PV
  \item Diagrama de espina de pez
\end{enumerate}

\subsection*{Pregunta 216}
¿Cuál de los siguientes métodos es el más adecuado para identificar los "pocos vitales"? :

\begin{enumerate}[label=\alph*)]
  \item Análisis de control de procesos.
  \item Análisis de causa y efecto.
  \item Análisis de Pareto.
  \item Análisis de tendencias.
\end{enumerate}

\subsection*{Pregunta 217}
¿Cuál de los siguientes no es una entrada al proceso de Verificar el alcance?

\begin{enumerate}[label=\alph*)]
  \item Plan para la dirección del proyecto
  \item Documentación de Requisitos
  \item Solicitudes de cambio aprobadas
  \item Matriz de trazabilidad de requerimientos
\end{enumerate}

\subsection*{Pregunta 218}
¿Cuál de los siguientes no es una entrada para el proceso de identificación de interesados?

\begin{enumerate}[label=\alph*)]
  \item El acta de constitución
  \item Los documentos de la adquisiciones y compras
  \item Los Factores ambientales de la empresa
  \item El caso de negocio
\end{enumerate}

\subsection*{Pregunta 219}
Cuál de los siguientes no es un proceso de gestión del alcance

\begin{enumerate}[label=\alph*)]
  \item Recopilar requisitos
  \item Controlar el alcance
  \item Definir actividades
  \item Crear la EDT
\end{enumerate}

\subsection*{Pregunta 220}
¿Cuál de los siguientes NO es un tipo de dependencia entre actividades?

\begin{enumerate}[label=\alph*)]
  \item Obligatoria (mandatory)
  \item Interna
  \item Externa
  \item \textbf{Discrecional (discretionary)}
\end{enumerate}

\subsection*{Pregunta 221}
¿Cuál de los siguientes NO se necesita para generar un estimado del proyecto?

\begin{enumerate}[label=\alph*)]
  \item \textbf{UnaEDT}
  \item Un diagrama de red
  \item Riesgos
  \item Solicitud de cambios
\end{enumerate}

\subsection*{Pregunta 222}
Cuál de los siguientes procesos de calidad del área de conocimiento del gestión de calidad forma parte del grupo de procesos de ejecución del proyecto?

\begin{enumerate}[label=\alph*)]
  \item Realizar el control de calidad
  \item Identificación de requisitos del producto
  \item Realizar el aseguramiento de la calidad)
  \item Planificar la calidad
\end{enumerate}

\subsection*{Pregunta 223}
¿Cuál de los siguientes procesos incluye preguntarles a los miembros del equipo acerca de los estimados de tiempo para sus actividades y llegar a acuerdos sobre la fecha de calendario para cada actividad?

\begin{enumerate}[label=\alph*)]
  \item Secuenciar las Actividades
  \item \textbf{Desarrollar el Cronograma}
  \item Definir el Alcance
  \item Desarrollar el Acta de Constitución del Proyecto
\end{enumerate}

\subsection*{Pregunta 224}
¿Cuál de los siguientes procesos realizará en primer lugar un Project manager?

\begin{enumerate}[label=\alph*)]
  \item Elaborar el Presupuesto
  \item Desarrollar el cronograma
  \item Estimar los costos de cada actividad
  \item Planificar las adquisiciones
\end{enumerate}

\subsection*{Pregunta 225}
¿cuál de los siguientes proyectos elegirías?

\begin{enumerate}[label=\alph*)]
  \item Proyecto A con 4.000 euros de costes hundidos y un coste de oportunidad de 3.000
  \item Proyecto B con costes hundidos de 5.000 euros y un coste de oportunidad de 3.000
  \item Proyecto C con unos costes hundidos de 5000
  \item Faltan datos para tomar la decisión
\end{enumerate}

\subsection*{Pregunta 226}
¿Cuál de los siguientes riesgos puede tener un impacto positivo?

\begin{enumerate}[label=\alph*)]
  \item La subida de precios de una materia prima relevante
  \item La aparición de una nueva tecnología de lanzamiento de puentes
  \item La dimisión del Project manager al cambiar de empresa
  \item La rotura de stock de un componente
\end{enumerate}

\subsection*{Pregunta 227}
Cuál de los siguientes será una salida del proceso de controlar los costes del proyecto

\begin{enumerate}[label=\alph*)]
  \item La base para estimaciones
  \item Las proyecciones del presupuesto
  \item El sistema de control de cambios
  \item La información de desempeño del trabajo
\end{enumerate}

\subsection*{Pregunta 228}
Cuál de los siguientes términos y condiciones no tiene porqué aparecer en el contrato ?

\begin{enumerate}[label=\alph*)]
  \item Criterios de aceptación de Entregables
  \item Condiciones y términos de pago
  \item Valor de la fianza a depositar
  \item Equipo del proyecto que desarrollará el contrato
\end{enumerate}

\subsection*{Pregunta 229}
¿Cuál de los siguientes tiene un orden cronológico adecuado?

\begin{enumerate}[label=\alph*)]
  \item Elaborar la acta de constitución, crear la EDT, Verificar el Alcance
  \item Definir el alcance, definir los interesados, controlar el alcance
  \item Recopilar requisitos, elaborar acta de constitución y crear la EDT
  \item Recopilar requisitos, crear la EDT y Definir el Alcance
\end{enumerate}

\subsection*{Pregunta 230}
¿cuál debe ser tu principal preocupación a la hora de gestionar los cambios?

\begin{enumerate}[label=\alph*)]
  \item Mantener un registro de problemas para las solicitudes de cambios
  \item Ignorar las solicitudes de cambios de las partes interesadas externas
  \item Prevenir cambios innecesarios
  \item Mantener informados a los patrocinadores de todas las solicitudes de cambios
\end{enumerate}

\subsection*{Pregunta 231}
Cuál es el grupo de procesos en el que no hay ningún proceso relativo a la gestión de la comunicación?

\begin{enumerate}[label=\alph*)]
  \item Inicio
  \item Planificación
  \item Seguimiento y control
  \item Cierre del proyecto
\end{enumerate}

\subsection*{Pregunta 232}
¿Cuál es la definición de Información confidencial?

\begin{enumerate}[label=\alph*)]
  \item Cualquier material considerado confidencial, sensible, marca registrada implícita o explícitamente y que no sea de Domicio público Información como patentes, marcas registradas,… b.
  \item El material que la compañía registre como Confidencial
  \item Material que no es de dominio público y que se utiliza dentro de la compañía
\end{enumerate}

\subsection*{Pregunta 233}
¿Cuál es la diferencia entre la línea base de coste y el presupuesto de coste?

\begin{enumerate}[label=\alph*)]
  \item La reserva de gestión
  \item Los costes indirectos
  \item La reserva de contingencia
  \item La diferencia entre el presupuesto previsto y el estimado a su finalización
\end{enumerate}

\subsection*{Pregunta 234}
¿Cuál es la estimación PERT por tres valores para una única actividad teniendo una visión optimista de 17 horas, una visión más probable de 23 horas y una visión pesimista de 41 horas?

\begin{enumerate}[label=\alph*)]
  \item \textbf{25 (17+4*23+41)/6}
  \item 23 d 27
\end{enumerate}

\subsection*{Pregunta 235}
Cuál es la mejor herramienta para mostrar al comité de seguimiento el estado del avance del proyecto durante la fase de ejecución?

\begin{enumerate}[label=\alph*)]
  \item El diagrama de redes
  \item La tabla de hitos
  \item La Estructura de Desglose de Tareas
  \item El diagrama de Gantt
\end{enumerate}

\subsection*{Pregunta 236}
Cuál es la principal función de un Gerente funcional

\begin{enumerate}[label=\alph*)]
  \item Controlar los recursos humanos
  \item Tomar las decisiones del proyecto cuando el Project manager no está disponible
  \item Definir los procesos del negocio
  \item Gestionar a los Project managers
\end{enumerate}

\subsection*{Pregunta 237}
Cuál es la relación de precedencia más habitual en el método de diagrama de de precedencias (PDM)

\begin{enumerate}[label=\alph*)]
  \item Inicio a Final (Start to finish, SF)
  \item Final a final (Finish-to-Finish, FF)
  \item Inicio a inicio
  \item Final a inicio (Finish-to-Start, FS)
\end{enumerate}

\subsection*{Pregunta 238}
¿Cuál no es un tipo de predecesor en una secuencia de actividades?

\begin{enumerate}[label=\alph*)]
  \item Externos: cuando el proyecto depende de Entregables que debe facilitar, por ejemplo, un proveedor
  \item Legales: son los que se deben cumplir de acuerdo a la legislación vigente.
  \item Discrecionales: que se estiman por la experiencia del Project manager y el equipo
  \item Obligatorios: vienen definidos por la propia naturaleza del trabajo
\end{enumerate}

\subsection*{Pregunta 239}
¿Cuáles de las siguientes no es una entrada para el proceso de llevar a cabo el control de calidad?

\begin{enumerate}[label=\alph*)]
  \item Registro de riesgos
  \item Activos de los procesos de la organización
  \item Plan de gestión de la calidad
  \item Medidas del desempeño del trabajo
\end{enumerate}

\subsection*{Pregunta 240}
¿Cuáles son las cinco fases de la vida de un proyecto y su orden correcto de acuerdo a las teorías de Tuckman:

\begin{enumerate}[label=\alph*)]
  \item Formación, Hacer equipo, Ejecución, Ajuste y Ejecución
  \item Formación, Turbulencias, Normalización, Desarrollo, Disolución
  \item Adquisición, Conocimiento, Desarrollo, Normalización, Disolución
  \item Formación, Turbulencias, Desarrollo, Normalización, Disolución
\end{enumerate}

\subsection*{Pregunta 241}
¿Cuáles son las principales habilidades interpersonales que debe tener el Project manager durante la fase de dirección y gestión de la ejecución del proyecto?:

\begin{enumerate}[label=\alph*)]
  \item habilidades de liderazgo, comunicación y negociación
  \item habilidades presupuestarias, de conocimiento general del negocio y de comunicación
  \item habilidades de comunicación, negociación y resolución de problemas
  \item liderazgo, comunicación y conocimientos generales del negocio
\end{enumerate}

\subsection*{Pregunta 242}
¿Cuáles son las principales ventajas de una estructura de tipo funcional?

\begin{enumerate}[label=\alph*)]
  \item El proceso organizativo está enfocado al cliente
  \item La estructura organizativa es estable
  \item Hay flexibilidad para formar y deshacer equipos de proyectos
  \item El proceso está enfocado a la gestión por proyectos
\end{enumerate}

\subsection*{Pregunta 243}
¿Cuáles son los cinco grupos de procesos en orden correcto?

\begin{enumerate}[label=\alph*)]
  \item Inicio, ejecución, planificación, supervisión y control y cierre
  \item Inicio, planificación, ejecución, supervisión y control y cierre.
  \item Inicio, supervisión y control, planificación, ejecución y cierre
  \item Inicio, planificación, supervisión y control, ejecución y cierre.
\end{enumerate}

\subsection*{Pregunta 244}
Cuando coges un proyecto ya comenzado, y estás revisando la documentación disponible, te encuentras con un documento que contiene una lista de habilidades y competencias y el momento en que serán necesarias en el proyecto. ¿De qué documento estamos hablando?

\begin{enumerate}[label=\alph*)]
  \item Los requisitos de recursos
  \item Los requisitos de plantilla
  \item Las evaluaciones de equipos
  \item El plan preliminar del proyecto
\end{enumerate}

\subsection*{Pregunta 245}
¿Cuándo debe realizarse el proceso de Verificar el Alcance?

\begin{enumerate}[label=\alph*)]
  \item Al final del proyecto
  \item Durante el proceso de supervisión y control del proyecto
  \item Al inicio de cada Fase del proyecto
  \item Cada vez que lo solicite el cliente
\end{enumerate}

\subsection*{Pregunta 246}
¿Cuándo debería hacerse el proceso Validar el Alcance?

\begin{enumerate}[label=\alph*)]
  \item Al final del proyecto
  \item Al inicio del proyecto
  \item \textbf{Al final de cada fase del proyecto}
  \item Durante los procesos de planificación
\end{enumerate}

\subsection*{Pregunta 247}
¿Cuándo es mejor tener reuniones de inicio?

\begin{enumerate}[label=\alph*)]
  \item Al comienzo del proyecto
  \item Cada vez que un entregable es finalizado
  \item Al inicio de cada fase
  \item Cuando el plan de comunicación ha sido aprobado
\end{enumerate}

\subsection*{Pregunta 248}
Cuando estaba a punto de cerrarse el proyecto de construcción de una nueva estación de autobuses, aparece un riesgo no identificado con anterioridad que, además, puede afectar a la funcionalidad y, por tanto, a la aceptación del proyecto por parte de la administración contratante. ¿Qué se debería hacer a continuación?

\begin{enumerate}[label=\alph*)]
  \item Avisar al patrocinador de los impactos potenciales en el coste, plazo y alcance
  \item Analizar cualitativamente el riesgo para poder valor su impacto
  \item Mitigar desarrollando un nuevo plan de respuestas
  \item Desarrollar un modificado que solucione el problema.
\end{enumerate}

\subsection*{Pregunta 249}
Cuando la estructura del edificio de 10 plantas que estás dirigiendo estaba ya medio levantada, el equipo de proyecto ha identificado un riesgo no previsto que te lo ha comunicado rápidamente. Ello te ha permitido evaluar el efecto potencial del cambio y has creado una solicitud de cambio. Cuál es el siguiente paso a dar?

\begin{enumerate}[label=\alph*)]
  \item Informar inmediatamente de la gravedad del hecho al patrocinador del proyecto
  \item evaluar la idoneidad del cambio informando del posible cambio al comité de cambios,
  \item proceder a la aprobación o rechazo de la propuesta de cambios o de actuación.
  \item preparar un plan de implantación de la solución adoptada
\end{enumerate}

\subsection*{Pregunta 250}
¿cuando se monitoriza y controla el trabajo del proyecto, qué es lo que no se lleva a cabo en él?

\begin{enumerate}[label=\alph*)]
  \item Se compara el desempeño previsto con el esperado
  \item Se monitoriza la puesta en marcha de los cambios aprobados
  \item Dar información para el reporte del estado de los proyectos
  \item Dar proyecciones del coste y plazo final estimado
\end{enumerate}

\subsection*{Pregunta 251}
¿Cuándo se realiza el análisis de las partes interesadas?

\begin{enumerate}[label=\alph*)]
  \item Cuando se elabora el plan de gestión de las comunicaciones del proyecto
  \item Cuando se crea el plan de dirección del proyecto
  \item Cuando se realiza la identificación de riesgos
  \item Cuando se elabora el Acta de constitución
\end{enumerate}

\subsection*{Pregunta 252}
Cuando trabajas en un entorno muy colaborador, ¿cuál de las siguientes técnicas será más efectiva para hacer equipo?

\begin{enumerate}[label=\alph*)]
  \item Implicar al equipo en el proceso de planificación
  \item Prometer una remuneración extra si el proyecto cumple objetivos Realizar actividades conjuntas (comidas, encuentros,…) c.
  \item Dar libertad para que el equipo se organice y actúe libremente
\end{enumerate}

\subsection*{Pregunta 253}
Cuando un Project manager informa de los objetivos de un proyecto a su equipo, éstos deben cumplir preferentemente todas estas características menos una. ¿cuál?

\begin{enumerate}[label=\alph*)]
  \item tener una fecha concreta de cumplimiento
  \item ser específicos y medibles
  \item realistas pero retadores
  \item ser definidos por el Project manager y trasladados al equipo
\end{enumerate}

\subsection*{Pregunta 254}
¿Cuánto dura un hito?

\begin{enumerate}[label=\alph*)]
  \item Es más corto que la duración de la actividad más larga.
  \item Es más corto que la actividad que representa .
  \item No tiene duración.
  \item Tiene la misma duración que la actividad que representa .
\end{enumerate}

\subsection*{Pregunta 255}
De acuerdo a la teoría de McGregor y sus conceptos de Teoría X y teoría Y, cuál de las siguientes afirmaciones es cierta?

\begin{enumerate}[label=\alph*)]
  \item En la teoría X los gestores confían en sus trabajadores y consideran que pueden delegar generando equipos autónomos
  \item En la teoría Y los gestores desconfían den su equipo tratando por tanto de no delegar nada y hacer micro-Management..
  \item En la teoría X se busca que el equipo tenga un sentimiento de pertenencia a la empresa, como si fuera una familia).
  \item En la teoría X los gestores no confían en sus trabajadores y consideran que deben tener una actitud autoritaria, no delegando.
\end{enumerate}

\subsection*{Pregunta 256}
De acuerdo al PMBOK® ¿Cuál de las siguientes no es razón para cerrar el contrato?

\begin{enumerate}[label=\alph*)]
  \item Por cambios importantes en el alcance del contrato
  \item Incumplimiento del contrato por parte del comprador
  \item Incumplimiento del contrato por vendedor
  \item Por conveniencia del comprador
\end{enumerate}

\subsection*{Pregunta 257}
De acuerdo al PMBOK®, ¿cual de las siguientes respuestas describe el objetivo de los grupos de procesos y los procesos individuales realizados dentro de cada grupo?:

\begin{enumerate}[label=\alph*)]
  \item Son utilizados para guiar el desarrollo de las actividades y procesos del proyecto, y más específicamente el proceso de ejecución.
  \item Se utilizan para monitorizar el cumplimiento de los entregables del proyecto y asegurar que las partes interesadas están satisfechas
  \item Son utilizados como guías para aplicar los conocimientos y habilidades de la Gestión de proyectos durante el proyecto.
  \item Se utilizan para medir, monitorizar y controlar el trabajo del proyecto para que se cumplan las expectativas de las partes interesadas.
\end{enumerate}

\subsection*{Pregunta 258}
De la siguiente descripción -los eventos que sabemos que van a pasar, los eventos que conocemos pero no sabemos si van a ocurrir y los eventos que desconocemos que puedan ocurrir- ¿cuántos se pueden considerar riesgos?

\begin{enumerate}[label=\alph*)]
  \item Ninguno de ellos
  \item Todos
\end{enumerate}

\subsection*{Pregunta 259}
De la siguiente información, ¿cuál es la que mejor describe el contenido completo de la documentación indexada del contrato , incluido el contrato cerrado, para su inclusión en al archivo final del proyecto

\begin{enumerate}[label=\alph*)]
  \item El plan de gestión del contrato
  \item Las auditorías del contrato
  \item El archivo del contrato
  \item Las lecciones aprendidas
\end{enumerate}

\subsection*{Pregunta 260}
Una de las muchas funciones del Project manager es la de avisar a los Project managers y jefes funcionales cuándo necesitará determinadas personas y cuándo quedarán liberadas. Esta información la refleja en:

\begin{enumerate}[label=\alph*)]
  \item El plan de comunicación
  \item En la Estructura de desglose de recursos
  \item En el Plan de recursos humanos
  \item En el Cronograma del proyecto
\end{enumerate}

\subsection*{Pregunta 261}
De las siguientes técnicas de identificación de riesgos, identifica cuál de ellas tiene como característica diferencial, mantener el anonimato de los participantes:

\begin{enumerate}[label=\alph*)]
  \item Entrevistas con expertos
  \item Tormenta de ideas
  \item Técnica Delphi
  \item Análisis DAFO
\end{enumerate}

\subsection*{Pregunta 262}
¿De qué depende el grado de éxito en la comunicación entre el Project manager y los miembros de su equipo?

\begin{enumerate}[label=\alph*)]
  \item De la relación que establezca el Project manager con ellos
  \item Del nivel jerárquico del Project manager en la compañía
  \item Del tamaño y complejidad del proyecto
  \item Del grado de conocimiento técnico que tenga el Project manager
\end{enumerate}

\subsection*{Pregunta 263}
De todos los ejemplos siguientes, ¿cuál no es un proyecto?

\begin{enumerate}[label=\alph*)]
  \item Desarrollar un nuevo servicio que va a ofrecer la empresa
  \item Cambiar la contraseña de usuario para acceder a una cuenta de red social
  \item Preparar una campaña de lanzamiento de un nuevo cosmético
  \item Diseñar e implantar un nuevo sistema de documentación de compras Respuestas razonadas
\end{enumerate}

\subsection*{Pregunta 264}
Debes planificar la elaboración de un informe de auditoría sobre una pequeña compañía. Sabes que dos de tus compañeros realizaron ese trabajo en cuatro semanas para tres empresas algo más pequeñas. Por lo tanto, decides estimar que te llevará 5 semanas. Este es un ejemplo de estimación:

\begin{enumerate}[label=\alph*)]
  \item estimación PERT por tres valores
  \item estimación por analogía
  \item estimación paramétrica
  \item Juicio de expertos
\end{enumerate}

\subsection*{Pregunta 265}
Demanda de mercado, una necesidad de negocios y/o una disposición legal son ejemplos de:

\begin{enumerate}[label=\alph*)]
  \item Razones para contratar a un director del proyecto .
  \item Razones por las que se inician los proyectos .
  \item Razones por las que las personas o las empresas se convierten en interesados .
  \item \textbf{Razones para patrocinar un proyecto. 103}
\end{enumerate}

\subsection*{Pregunta 266}
Dentro del proceso de establecer la línea base de coste para luego poder medir el desempeño del proyecto en relación a los costes inicialmente previstos, todos excepto uno de los siguientes elementos son salidas del proceso

\begin{enumerate}[label=\alph*)]
  \item actualización del plan de gestión del coste
  \item la línea de base del coste
  \item las estimaciones de costes de las actividades
  \item los cambios requeridos
\end{enumerate}

\subsection*{Pregunta 267}
Una dependencia que necesita que el diseño se finalice antes de que se pueda iniciar la manufactura es un ejemplo de:

\begin{enumerate}[label=\alph*)]
  \item Dependencia discrecional.
  \item Dependencia externa.
  \item Dependencia obligatoria.
  \item Dependencia del alcance.
\end{enumerate}

\subsection*{Pregunta 268}
Desarrollar el espíritu de equipo en la Gestión de proyectos es fundamental y por muchas razones. ¿Cuál de las siguientes no sería una de las razones?

\begin{enumerate}[label=\alph*)]
  \item Por la propia naturaleza finita y temporal del proyecto
  \item porque el equipo puede no tener experiencia de trabajo conjunta
  \item que todos los miembros estén alineados con los objetivos del proyecto
  \item para facilitar la asignación de remuneraciones y recompensas
\end{enumerate}

\subsection*{Pregunta 269}
Desarrollar opciones y acciones para las oportunidades del proyecto y riesgos es parte de cuál de los siguientes procesos:

\begin{enumerate}[label=\alph*)]
  \item Lecciones Aprendidas.
  \item \textbf{Planificar la Respuesta a los Riesgos.}
  \item Analizar el Riesgo del Proyecto.
  \item Evaluar el Riesgo del Proyecto.
\end{enumerate}

\subsection*{Pregunta 270}
¿Cómo describirías mejor el ``análisis del producto= entre el Project manager y el cliente?

\begin{enumerate}[label=\alph*)]
  \item Es la búsqueda conjunta de una buena descripción del producto
  \item Es definir los procedimientos de calidad que debe pasar el producto
  \item Conocer mejor el producto deseado por el cliente para poder elaborar mejor el Enunciado del Trabajo
  \item Clarificar cuál va ser la relación entre los procesos de gestión del producto y su coordinación con los de gestión del proyecto
\end{enumerate}

\subsection*{Pregunta 271}
Desde el equipo de marketing te solicitan que cambies la pantalla de entrada para la intranet de la compañía para hacerla más sencilla y que, además, te mantenga tu anterior contraseña. Esto es considerado:

\begin{enumerate}[label=\alph*)]
  \item el inicio de un proyecto
  \item operaciones en marcha
  \item un proyecto
  \item la ejecución de un proyecto
\end{enumerate}

\subsection*{Pregunta 272}
Deseas conocer qué objetivos de negocio van a tratar de cumplir un grupo de proyectos y programas. ¿Cuál es el mejor lugar para encontrar esta información?

\begin{enumerate}[label=\alph*)]
  \item En el Plan de Proyecto
  \item En el Portfolio de la Compañía
  \item En el Acta de constitución (Project charter) del Proyecto
  \item En la base de programas de la compañía
\end{enumerate}

\subsection*{Pregunta 273}
Diagramas de Causa y Efecto (Ishikawa o Espina de Pescado), Diagramas Pareto, e Histogramas son herramientas que se usan como parte de:

\begin{enumerate}[label=\alph*)]
  \item El aseguramiento de calidad.
  \item El control de la calidad.
  \item El análisis cualitativo de riesgos.
  \item \textbf{La gestión de la calidad.}
\end{enumerate}

\subsection*{Pregunta 274}
Como Project manager cuál de las siguientes técnicas o herramientas podrían utilizar para estimar costes?

\begin{enumerate}[label=\alph*)]
  \item Estimación por analogía
  \item Análisis de ofertas de proveedores
  \item Relaciones históricas
  \item Coste de la calidad
\end{enumerate}

\subsection*{Pregunta 275}
Como Project manager, ¿cuál de los siguientes es un resultado que deberás tener preparado en el proceso de cierre de las compras?

\begin{enumerate}[label=\alph*)]
  \item Adquisiciones cerradas
  \item Aceptación verbal de tu cliente
  \item Auditorías de la adquisición
  \item Entregables aceptados
\end{enumerate}

\subsection*{Pregunta 276}
Como Project manager de un proyecto de lanzamiento de un nuevo producto informático, sabes que uno de los factores clave será el coste del producto en su fase de desarrollo. A la hora de preparar el plan de gestión de calidad, deberás:

\begin{enumerate}[label=\alph*)]
  \item Dejar que el gestor del programa decida el nivel de calidad más acorde con la empresa.
  \item Imputarle al cliente todos los costes de calidad
  \item Poner en marcha controles de calidad cuyos beneficios en el Ciclo de vida del proyecto sean superiores a los costes
  \item Poner en marcha procesos que, como indicaba Juran, logren cero errores
\end{enumerate}

\subsection*{Pregunta 277}
Como Project manager dentro de una organización enfocada a proyectos, tienes que valorar los objetivos de desempeño de tu equipo para cuantificar el variable que reciben por objetivos. ¿Qué documento deberás consultar?

\begin{enumerate}[label=\alph*)]
  \item El plan de reconocimientos y recompensas
  \item El plan de Dirección del personal
  \item El plan de gestión de los recursos humanos
  \item El presupuesto del proyecto
\end{enumerate}

\subsection*{Pregunta 278}
Como Project manager estás monitorizando y registrando los resultados del desempeño y recomendando los cambios necesarios. ¿Qué obtienes tras ello?

\begin{enumerate}[label=\alph*)]
  \item Medidas de control de calidad
  \item El plan de gestión de calidad
  \item Métricas del control de calidad
  \item Listas de control de calidad
\end{enumerate}

\subsection*{Pregunta 279}
Como Project manager formas parte de una evaluación de ofertas entre las que se encuentra una compañía en la que tu familia tiene una pequeña participación. De las siguientes opciones de actuación, ¿Cuál no es adecuada?

\begin{enumerate}[label=\alph*)]
  \item Ser honesto y ecuánime en las evaluaciones y atenerte a los criterios de adjudicación
  \item Informar al cliente y al consejo de administración de la compañía
  \item No participar en la mesa de adjudicación aun siendo tu responsabilidad,
  \item Delegar temporalmente la responsabilidad en ese caso aunque tu seas el mejor preparado
\end{enumerate}

\subsection*{Pregunta 280}
Como Project manager recibes los diagramas de control en que se anotan los resultados de las mediciones y se estudia si las variaciones se encuentran en valores aceptables. Si, nos encontramos con varios valores que muestran una variación respecto a la media y el intervalo aceptable muy altos o muy bajos podremos pensar que el proceso está fuera de control. Si por el contrario, lo que se ve, por el contrario es un ajuste al intervalo de aceptación con alguna fluctuación de algún valor de vez en cuando pero que no preocupa las llamaremos:.

\begin{enumerate}[label=\alph*)]
  \item Variaciones normales
  \item Variaciones sistemáticas
  \item Variaciones aleatorias
  \item Variaciones puntuales.
\end{enumerate}

\subsection*{Pregunta 281}
Como Project manager te encargan la responsabilidad de desarrollar un nuevo parque eólico en el extranjero. Ese mercado requiere diseñar nuevos componentes de la turbina para garantizar su óptimo funcionamiento y cumplir con los estándares locales. Para crear una buena Estructura de Desglose de Tareas decides:

\begin{enumerate}[label=\alph*)]
  \item Presentar a las partes interesadas, la EDT, la Matriz RACI y la RBS.
  \item registrar el código de identificación, una descripción del trabajo, el responsable de la organización y un cronograma de los principales hitos en el diccionario de la EDT para todos los componentes de la EDT.
  \item Elaborar un nuevo Enunciado del Trabajo que incluya las especificaciones de los nuevos componentes
  \item Al ser un país nuevo, preparar un plan de gestión de riesgos que nos complete las actividades a realizar con las correspondientes al plan de Contingencia
\end{enumerate}

\subsection*{Pregunta 282}
Como Project manager te encuentras con un diagrama de red como el siguiente en el que tienes nueva actividades a realizar para concluir el proyecto, con las precedencias que en él se indican. C: 3 dias D: 4 dias B: 7 dias F: 2 dias E: 10 dias I: 2 dias L: 4 dias G: 4 dias H: 8 dias ¿Cuál sería el camino crítico?

\begin{enumerate}[label=\alph*)]
  \item I-B-C-D-L
  \item I-E-C-D-L
  \item I-E-F-L
  \item I-G-H-L
\end{enumerate}

\subsection*{Pregunta 283}
Como Project manager tienes una lista de posibles proveedores a los que puedes acudir en caso de no tener tiempo para hacer todo el proceso de selección y adjudicación. ¿Cómo se denomina esta lista?

\begin{enumerate}[label=\alph*)]
  \item Lista de vendedores cualificados
  \item Lista de vendedores preseleccionados
  \item Lista de licitadores invitados
  \item Lista de vendedores seleccionados
\end{enumerate}

\subsection*{Pregunta 284}
Como Project manager vas a proceder a solicitar ofertas a posibles proveedores, para lo cual vas a preparar una solicitud de cotización y para ello debes preparar los Documentos de la adquisición. ¿Cuál de los siguientes contenidos no forma parte de este documento?

\begin{enumerate}[label=\alph*)]
  \item Información general del proyecto
  \item Precios unitarios y precio total
  \item Enunciado del alcance del contrato
  \item Las condiciones y términos del contrato.
\end{enumerate}

\subsection*{Pregunta 285}
Como Project manager vas a trabajar en un país en el que es práctica habitual hacer regalos a las personas que deben agilizar los trámites de licencias de obras de la nueva planta de criterios. Mañana tienes una reunión con el funcionario que debe tramitar la licencia. ¿Qué harías?

\begin{enumerate}[label=\alph*)]
  \item Averiguar con discreción cuál es el importe requerido por el funcionario y si no afecta a la rentabilidad del proyecto, entregárselo ya que si no, el proyecto puede detenerse indefinidamente.
  \item Puesto que el obsequio no es para cambiar su opinión o percepción, ni el propio Project manager va a obtener beneficios económicos a costa de la empresa puede pagar la cantidad establecida por el funcionario.
  \item Consultar a la compañía las políticas propias en este sentido y ver si el regalo excede lo establecido en la empresa.
  \item Se asegurará de que es una cantidad que va a ser utilizada para mejorar las condiciones del entorno social del proyecto y sus condiciones ambiéntales y culturales.
\end{enumerate}

\subsection*{Pregunta 286}
Como director de proyectos, no puedes dedicar el tiempo que desearías para interactuar con los interesados. ¿Conocer a cuál de los siguientes interesados sería tu prioridad?

\begin{enumerate}[label=\alph*)]
  \item \textbf{El interesado que es experto en el producto del proyecto, pero que no está interesado en implementarlo en su departamento.}
  \item El gerente del departamento que utilizará el producto del proyecto. Se sabe que se resiste a los cambios.
  \item El patrocinador del proyecto, con quien has trabajado con éxito en muchos proyectos.
  \item El empleado del departamento que no está familiarizado con el producto del proyecto, pero que está abierto a los impactos positivos que cree que el producto tendrá en su entorno de trabajo .
\end{enumerate}

\subsection*{Pregunta 287}
Diriges una obra civil en la que el representante del promotor viene todos los días a hacer su visita de Inspección semanal a la hora de comer y te pide que le invites a comer tras cada visita. ¿Qué es lo mejor que puedes hacer?

\begin{enumerate}[label=\alph*)]
  \item Decirle de forma correcta que no es ético desde tu punto de vista como Project manager
  \item Invitarle a comer pero siendo consciente que es por la empresa
  \item Invitarle a comer pero poniéndolo en conocimiento del superior del representante
  \item Invitarle a comer ya que no te queda más remedio
\end{enumerate}

\subsection*{Pregunta 288}
Diseño de experimentos:

\begin{enumerate}[label=\alph*)]
  \item Identifica qué variables tendrán una mayor influencia en un resultado de calidad .
  \item Identifica qué variables tendrán una menor influencia en un resultado de calidad.
  \item Determina qué es un resultado de calidad.
  \item \textbf{Determina los métodos a utilizar para la investigación y el desarrollo .}
\end{enumerate}

\subsection*{Pregunta 289}
Dos Project managers trabajan en una organización que tiene una estructura de matriz débil. En estos casos, es difícil considerarles Project managers. Podrían ser Coordinadores de Proyecto o Expedidores de Proyecto. ¿Cuál es la diferencia entre estas dos figuras?

\begin{enumerate}[label=\alph*)]
  \item El expedidor de proyectos no puede tomar decisiones
  \item El expedidor de proyectos puede tomar más decisiones
  \item El expedidor de proyectos reporta a un nivel más alto
  \item El expedidor de proyectos tiene cierta autoridad
\end{enumerate}

\subsection*{Pregunta 290}
Durante cinco años has trabajado como Project manager y. al ser una persona proactiva has participado activamente en la elaboración de todas las plantillas y procedimientos utilizados por la empresa dedicándole tiempo en fines de semana en tu casa. Ahora has recibido una oferta de otra compañía y dejas tu anterior empresa ¿Cuál de las siguientes opciones es válida?

\begin{enumerate}[label=\alph*)]
  \item Puesto que las plantillas son el resultado de tu trabajo, te las podrás llevar para ajustarlas a las necesidades de la nueva empresa
  \item Te las llevarás siempre que tengas el permiso y consentimiento de la compañía en la que has trabajado
  \item El Project manager deberá comprobar primeramente si está protegida
  \item El Project manager deberá comprobar si en el registro aparece que la información es privada
\end{enumerate}

\subsection*{Pregunta 291}
¿Durante cuál parte del proceso de dirección de proyectos se crea el enunciado del alcance del proyecto?

\begin{enumerate}[label=\alph*)]
  \item Iniciación
  \item Planificación
  \item \textbf{Ejecución}
  \item Seguimiento y control
\end{enumerate}

\subsection*{Pregunta 292}
Durante el proceso Identificar los Riesgos, un director de proyectos y los interesados utilizaron diversos métodos para identificar los riesgos, y luego crearon una larga lista de esos riesgos. Luego, el director del proyecto se aseguró de que se entendieran todos los riesgos y de que se identificaran todos los disparadores. Posteriormente, en el proceso Planificar la Respuesta a los Riesgos, tomó todos los riesgos identificados por los interesados y determinó las maneras en que habrían de mitigarse . ¿Qué hizo mal?

\begin{enumerate}[label=\alph*)]
  \item El director del proyecto debió haber esperado hasta el proceso Realizar el Análisis Cualitativo de Riesgos para involucrar a los interesados.
  \item Más personas debieron estar involucradas en el proceso Planificar la Respuesta a los Riesgos.
  \item El director del proyecto debió haber creado soluciones temporales.
  \item Los disparadores no deben identificarse hasta el proceso Identificar los Riesgos.
\end{enumerate}

\subsection*{Pregunta 293}
Durante el proyecto, estás recogiendo datos sobre cómo va el proyecto en sus diferentes líneas de base para distribuirlas a las partes interesadas. la información enviada incluye cómo se están utilizando los recursos para alcanzar los objetivos del proyecto. Esto se realiza mediante el:

\begin{enumerate}[label=\alph*)]
  \item Proceso del plan de comunicación
  \item Proceso de informe del desempeño
  \item Proceso de distribución de información
  \item Gestión del valor ganado
\end{enumerate}

\subsection*{Pregunta 294}
Durante la ejecución del proyecto, se realizan un gran número de cambios al proyecto. El director del proyecto debe:

\begin{enumerate}[label=\alph*)]
  \item Esperar hasta que se conozcan todos los cambios e imprimir un nuevo cronograma.
  \item Realizar cambios aprobados como sea necesario, pero mantener la línea base del cronograma .
  \item Realizar sólo los cambios aprobados por la gerencia.
  \item \textbf{Hablar con la gerencia antes de realizar cualquier cambio. 246}
\end{enumerate}

\subsection*{Pregunta 295}
Durante la ejecución del proyecto, un miembro del equipo identifica un riesgo que no está en el registro de riesgos. ¿Qué debes hacer?

\begin{enumerate}[label=\alph*)]
  \item Obtener más información sobre la manera en la que el miembro del equipo identificó el riesgo, debido a que tú llevaste a cabo un análisis detallado de riesgos y no identificaste el riesgo.
  \item \textbf{Ignorar el riesgo, pues los riesgos ya fueron identificados durante la planificación del proyecto.}
  \item Informar al cliente sobre el riesgo.
  \item Analizar el riesgo. 448
\end{enumerate}

\subsection*{Pregunta 296}
Durante la ejecución del proyecto, un miembro del equipo se acerca al director del proyecto porque no está seguro sobre qué trabajo debe realizar en el proyecto. ¿Cuál de los siguientes documentos contiene descripciones detalladas de los paquetes de trabajo?

\begin{enumerate}[label=\alph*)]
  \item Diccionario de la EDT
  \item Lista de actividades
  \item Enunciado del alcance del proyecto
  \item \textbf{El plan para la gestión del alcance}
\end{enumerate}

\subsection*{Pregunta 297}
Durante la fase de ejecución del proyecto, un miembro del equipo solicita aclaraciones al Project manager puesto que no tiene claro qué trabajo le corresponde realizar. ¿Cuál de los siguientes documentos tienen una descripción detallada de los paquetes de trabajo?

\begin{enumerate}[label=\alph*)]
  \item El Diccionario de la EDT
  \item La Lista de actividades
  \item El Enunciado del Trabajo
  \item El plan de gestión del alcance del producto
\end{enumerate}

\subsection*{Pregunta 298}
Durante la fase de planificación consideraste que, teniendo en cuenta los riesgos que tenía el proyecto, había que incluir una partida en el presupuesto para afrontar los riesgos previstos e imprevistos. Durante la ejecución del proyecto vas controlando si esa cantidad es suficiente o no y, en la medida de lo posible tomas medidas al respecto. ¿Qué técnica o herramienta de seguimiento y control de riesgos estás utilizando?

\begin{enumerate}[label=\alph*)]
  \item Re-evaluación de riesgos
  \item Auditoría de riesgos
  \item Análisis de reservas
  \item Análisis de variación y tendencias
\end{enumerate}

\subsection*{Pregunta 299}
Durante la fase de planificación del proyecto, algunas de las partes interesadas que se llevan muy bien con el Project manager le solicitan aclaración por lo que consideran cambios en la forma de proceder de la compañía en la Gestión de proyectos y no están seguros de que este nuevo sistema pueda funcionar mejor. ¿Qué debería hacer el Project manager?

\begin{enumerate}[label=\alph*)]
  \item Tranquilizar a los interesados indicándoles que en todo momento estarán informados de acuerdo al plan de comunicación establecido
  \item Convocar una reunión con la Dirección de Programas para aclarar las cosas
  \item Notificárselo a la Oficina de Gestión de proyectos (Project Management Office, PMO)
  \item Indicarles que como responsable del proyecto, cada proyecto tiene sus peculiaridades y éste en concreto, se gestionará en algunos aspectos concretos de forma diferente.
\end{enumerate}

\subsection*{Pregunta 300}
Durante la finalización del trabajo del proyecto, el patrocinador pide al director de proyectos que le informe sobre cómo va el proyecto . Para preparar el informe , el director de proyectos pregunta a todos los miembros del equipo qué porcentaje de su trabajo está finalizado . Hay un miembro del equipo que desde el principio ha sido difícil de dirigir. En respuesta a la pregunta sobre el porcentaje que lleva finalizado, el miembro del equipo pregunta, "¿Porcentaje finalizado de qué?" Cansado de tales comentarios, el director de proyectos informa al jefe del miembro del equipo que el miembro del equipo no está cooperando. ¿Cuál de los siguientes es MÁS probable que sea el verdadero problema?

\begin{enumerate}[label=\alph*)]
  \item El director del proyecto no obtuvo el compromiso del gerente para los recursos del proyecto.
  \item El director de proyectos no creó un sistema de recompensa adecuado para los miembros del equipo para mejorar su cooperación .
  \item El director de proyectos debió haber tenido una reunión con el jefe del miembro del equipo la primera vez que el miembro del equipo causó problemas.
  \item El director del proyecto no asignó paquetes de trabajo.
\end{enumerate}

\subsection*{Pregunta 301}
Durante la planificación del proyecto en una organización matricial, el director de proyectos determina que se necesitan más recursos humanos. ¿A quién le solicitaría estos recursos?

\begin{enumerate}[label=\alph*)]
  \item Al director de la oficina de dirección de proyectos
  \item Al gerente funcional
  \item \textbf{Al equipo}
  \item Al patrocinador del proyecto
\end{enumerate}

\subsection*{Pregunta 302}
Durante las reuniones del equipo del proyecto, el director de proyectos le pide a cada miembro del equipo que describa el trabajo que está llevando a cabo y asigna nuevas actividades a los miembros del equipo. La duración de estas reuniones ha aumentado debido a que hay muchas actividades diferentes que deben ser asignadas. Esto podría estar ocurriendo por todas las siguientes razones, EXCEPTO:

\begin{enumerate}[label=\alph*)]
  \item Ausencia de una EDT.
  \item \textbf{Ausencia de una matriz de asignación de responsabilidades.}
  \item Ausencia de nivelación de recursos.
  \item El equipo no se involucró en la planificación del proyecto.
\end{enumerate}

\subsection*{Pregunta 303}
¿Durante qué grupo de procesos de gestión del proyecto tienen más influencia las partes interesadas y el riesgo de no cumplir los objetivos es mayor?

\begin{enumerate}[label=\alph*)]
  \item Planificación
  \item Ejecución
  \item Inicio
  \item Supervisión y Control
\end{enumerate}

\subsection*{Pregunta 304}
¿Durante qué grupo de procesos de la dirección de proyectos se crean las proyecciones del presupuesto?

\begin{enumerate}[label=\alph*)]
  \item Seguimiento y control
  \item Planificación
  \item \textbf{Iniciación}
  \item Ejecución
\end{enumerate}

\subsection*{Pregunta 305}
¿Durante qué grupo de procesos el equipo analiza y mide el trabajo que está siendo realizado en el proyecto?

\begin{enumerate}[label=\alph*)]
  \item Iniciación
  \item Ejecución
  \item Seguimiento y control
  \item Cierre
\end{enumerate}

\subsection*{Pregunta 306}
¿Durante qué proceso de la gestión de los riesgos se determinan las soluciones temporales?

\begin{enumerate}[label=\alph*)]
  \item Identificar los Riesgos
  \item Realizar el Análisis Cuantitativo de Riesgos
  \item Planificar la Respuesta a los Riesgos
  \item Controlar los Riesgos
\end{enumerate}

\subsection*{Pregunta 307}
Durante una reunión del equipo del proyecto, un miembro del equipo sugiere una mejora para el alcance que va más allá del alcance del acta de constitución del proyecto. El director del proyecto señala que el equipo necesita concentrarse en finalizar todo el trabajo y sólo el trabajo solicitado. Esto es un ejemplo de:

\begin{enumerate}[label=\alph*)]
  \item \textbf{Proceso de la gestión de los cambios .}
  \item Gestión del alcance .
  \item Análisis de calidad .
  \item Desglose del alcance .
\end{enumerate}

\subsection*{Pregunta 308}
Durante una reunión del equipo, el equipo agrega un área específica de trabajo extra al proyecto porque han determinado que beneficiará al cliente. ¿Cuál es el error de esta situación?

\begin{enumerate}[label=\alph*)]
  \item El equipo está "bañando en oro" (añadiendo funcionalidad extra).
  \item Esto no debería realizarse en las reunione s.
  \item Nada. Es la manera de alcanzar y superar las expectativas de los clientes .
  \item \textbf{Nada . El director de proyectos está en control de la situación . 325}
\end{enumerate}

\subsection*{Pregunta 309}
Durante una reunión del equipo, un miembro del equipo sugiere una ampliación de alcance que va más allá del alcance establecido en el Acta de Constitución. El Director del Proyecto les recuerda que se tienen que centrar en completar el trabajo requerido y solo el trabajo requerido. Esto es un ejemplo de: :

\begin{enumerate}[label=\alph*)]
  \item Gestión del alcance del proyecto
  \item Proceso de gestión de cambios
  \item Descomposición del alcance
  \item Análisis de calidad.
\end{enumerate}

\subsection*{Pregunta 310}
Durante una reunión del equipo, un miembro hace una pregunta respecto de las medidas que se utilizarán en el proyecto para evaluar el desempeño . El miembro del equipo siente que algunas de las medidas relacionadas con las actividades que le han sido asignadas no son medidas válidas . ¿Es MEJOR considerar el proyecto en qué parte de los grupos de procesos de dirección de proyectos?

\begin{enumerate}[label=\alph*)]
  \item Cierre
  \item Seguimiento y control
  \item Ejecución
  \item Iniciación
\end{enumerate}

\subsection*{Pregunta 311}
El acta de constitución suele aprobarla el patrocinador del proyecto, aunque en algunos casos pueden ser aprobados por las partes interesadas. ¿Cómo podemos mejor definir el papel del patrocinador en un proyecto?

\begin{enumerate}[label=\alph*)]
  \item El patrocinador gestiona el proyecto
  \item El patrocinador provee o aprueba la provisión de fondos al proyecto
  \item El patrocinador verifica que todos los entregables están completos
  \item El patrocinador negocia los contratos
\end{enumerate}

\subsection*{Pregunta 312}
El alcance del producto La principal diferencia entre un proceso y un proyecto es:

\begin{enumerate}[label=\alph*)]
  \item El proceso tiene una duración limitada aunque repetitiva
  \item Los proyectos se componen de procesos
  \item \textbf{El proyecto es singulary repetitivo}
\end{enumerate}

\subsection*{Pregunta 313}
El análisis de Montecarlo es un tipo de análisis?

\begin{enumerate}[label=\alph*)]
  \item Método del camino crítico De escenarios ¿Qué pasa si…?
  \item Equilibrio de recursos
  \item Método de la cadena crítica
\end{enumerate}

\subsection*{Pregunta 314}
El análisis de reservas es una herramienta o técnica de todos los procesos siguientes menos:

\begin{enumerate}[label=\alph*)]
  \item elaboración del presupuesto
  \item estimación de los recursos de cada actividad
\end{enumerate}

\subsection*{Pregunta 315}
El análisis de valor se lleva a cabo para obtener:

\begin{enumerate}[label=\alph*)]
  \item Mayor valor de análisis de costo.
  \item Que la dirección crea en el proyecto.
  \item \textbf{Que el equipo crea en el proyecto .}
  \item Una forma menos costosa de realizar el mismo trabajo.
\end{enumerate}

\subsection*{Pregunta 316}
El análisis Monte Cado se utiliza para :

\begin{enumerate}[label=\alph*)]
  \item Recibir una indicación del riesgo que involucra el proyecto.
  \item Estimar la duración de una actividad.
  \item Simular posibles polémicas de calidad en el proyecto .
  \item \textbf{Comprobarle a la gerencia que se necesita personal adicional.}
\end{enumerate}

\subsection*{Pregunta 317}
El balance de doble amortización decreciente es una forma de:

\begin{enumerate}[label=\alph*)]
  \item Depreciación desacelerada.
  \item Depreciación de línea recta.
  \item Depreciación acelerada.
  \item Costeo del ciclo de vida .
\end{enumerate}

\subsection*{Pregunta 318}
El ciclo de vida de un proyecto difiere del ciclo de vida de un producto en que el ciclo de vida del proyecto:

\begin{enumerate}[label=\alph*)]
  \item No tiene ninguna metodología concreta
  \item Difiere según sectores industriales y mercados
  \item Puede abarcar varios proyectos
  \item Describe las actividades para la gestión del proyecto
\end{enumerate}

\subsection*{Pregunta 319}
El ciclo de vida del proyecto difiere del proceso de la dirección de proyectos en que el proceso de la dirección de proyectos:

\begin{enumerate}[label=\alph*)]
  \item Es igual para todos los proyectos.
  \item No incorpora una metodología.
  \item Es diferente para cada industria.
  \item \textbf{Puede generar muchos proyectos.}
\end{enumerate}

\subsection*{Pregunta 320}
El cierre incluye todas las siguientes acciones, EXCEPTO :

\begin{enumerate}[label=\alph*)]
  \item Determinar medidas del desempeño .
  \item Hacer entrega del producto del proyecto.
  \item Documentar hasta qué grado se cerró de forma correcta cada fase del proyecto después de su conclusión.
  \item \textbf{Actualizar los activos de los procesos de la organización . 102}
\end{enumerate}

\subsection*{Pregunta 321}
El cliente, cuando el proyecto de lanzamiento de un juego de ordenador lleva ya desarrollado un 70\% de su alcance, considera que el enfoque inicialmente dado al producto debe cambiar y hacer algunos cambios relevantes. ¿Qué es lo primero que debe hacer el Project manager?

\begin{enumerate}[label=\alph*)]
  \item Informar al cliente que el proyecto está ya demasiado avanzado como para integrar el cambio en el proyecto.
  \item Evaluar el impacto del cambio del cambio en el proyecto e informarle de las opciones posibles y su influencia en el coste, plazo, riesgos,…
  \item Dejar pasar la solicitud confiando en que el cliente dejará de ver como necesario dicho cambio
  \item Hacer caso al cliente y aplicar los cambios indicados por él y lanzar el producto al mercado.
\end{enumerate}

\subsection*{Pregunta 322}
El cliente de un proyecto le dice al director de proyectos que se quedó sin dinero para pagar el proyecto. ¿Qué debe realizar PRIMERO el director del proyecto?

\begin{enumerate}[label=\alph*)]
  \item \textbf{Trasladar más trabajo para que sea realizado más adelante en el cronograma para que el cliente tenga tiempo de obtener los fondos.}
  \item Cerrar el Proyecto o Fase.
  \item Detener el trabajo.
  \item Liberar parte del equipo del proyecto.
\end{enumerate}

\subsection*{Pregunta 323}
El cliente exige cambios a las especificaciones del producto que agregarán sólo dos semanas a la ruta crítica. ¿Cuál de las siguientes opciones indica lo MEJOR que puede realizar el director del proyecto?

\begin{enumerate}[label=\alph*)]
  \item Comprimir el cronograma para recuperar las dos semanas.
  \item Cortar el alcance para recuperar las dos semanas.
  \item Consultar al patrocinador sobre opciones.
  \item \textbf{Aconsejar al cliente acerca del impacto del cambio. 148}
\end{enumerate}

\subsection*{Pregunta 324}
El cliente ha aceptado el alcance del proyecto completado. No obstante, las lecciones aprendidas requeridas por la oficina de dirección de proyectos no se han completado. ¿Cuál es el estado del proyecto?

\begin{enumerate}[label=\alph*)]
  \item El proyecto está incompleto porque necesita volver a planificarse.
  \item El proyecto no está completo hasta que todos los entregables del producto y del proyecto estén completos y aceptados .
  \item El proyecto está completo porque el cliente ha aceptado los entregables.
  \item \textbf{El proyecto está completo porque ha alcanzado su fecha de entrega.}
\end{enumerate}

\subsection*{Pregunta 325}
El cliente, que tiene dudas sobre la calidad de los millones de etiquetas de yogures que han impreso en la compañía solicita que se revise su adecuación a los requisitos. ¿Qué es lo que harías, como Project manager?

\begin{enumerate}[label=\alph*)]
  \item Revisar todas y cada una de las etiquetas
  \item Tomar una muestra representativa y realizar la revisión sobre ella
  \item Asegurarle al cliente que el sistema de calidad evita estos errores
  \item Inspeccionar las treinta últimas en presencia del cliente
\end{enumerate}

\subsection*{Pregunta 326}
El cliente responsable de supervisar tu proyecto te pide que le entregues una estimación del costo por escrito que sea 30 por ciento superior a tu estimado del costo del proyecto . Te explica que el proceso de elaboración del presupuesto exige que los directores hagan estimados pesimistas para asegurarse de que se asigne dinero suficiente para los proyectos. ¿Cuál es la MEJOR manera de enfrentarte a esto?

\begin{enumerate}[label=\alph*)]
  \item \textbf{Agregar el 30 por ciento como un monto total de fondo para contingencias para enfrentarse a los riesgos del proyecto.}
  \item Agregar el 30 por ciento a tu estimado de costo repartiéndolo uniformemente a lo largo de todas las actividades del proyecto .
  \item Crear una primera línea base de costos para la asignación del presupuesto y una segunda para el plan para la dirección del proyecto real.
  \item Pedir información sobre los riesgos que podrían conducir a que tu estimado sea demasiado bajo.
\end{enumerate}

\subsection*{Pregunta 327}
El conocimiento que el equipo y partes interesadas va adquiriendo durante la vida del proyecto se llama:

\begin{enumerate}[label=\alph*)]
  \item Mejores prácticas
  \item Benchmarking
  \item Factores culturales de la empresa
  \item lecciones aprendidas
\end{enumerate}

\subsection*{Pregunta 328}
El desarrollo de la línea base del alcance puede describirse MEJOR como la responsabilidad de:

\begin{enumerate}[label=\alph*)]
  \item \textbf{Los gerentes funcionales .}
  \item El equipo del proyecto .
  \item Todos los interesados.
  \item El expedidor del proyecto. 189
\end{enumerate}

\subsection*{Pregunta 329}
El Project manager está tratando de definir el alcance del producto. ¿Cuál de las tareas siguientes es la más importante a abordar?

\begin{enumerate}[label=\alph*)]
  \item El caso de negocio
  \item La estructura de desglose de tareas
  \item El registro de riesgos e incidencias
  \item Verificar que todas las partes interesadas han aportado su visión
\end{enumerate}

\subsection*{Pregunta 330}
El Project manager está tratando de reunirse con diferentes personas que pueden verse afectadas por el proyecto tratando de identificar sus necesidades así como posibles limitaciones e hipótesis a considerar para poder tener claras las expectativas del proyecto y poder avanzar en la elaboración del acta de constitución. ¿cuál de las siguientes respuestas define mejor el trabajo que está realizando?

\begin{enumerate}[label=\alph*)]
  \item Gestión de las expectativas de las partes interesadas
  \item Comité de Control de cambios, CCB
  \item Análisis del equipo de proyecto
  \item Identificación de partes interesadas
\end{enumerate}

\subsection*{Pregunta 331}
El Project manager se hace cargo de un nuevo proyecto y de un equipo de trabajo que llevan trabajando juntos varios años. Como tal, es responsable de recabar toda la información histórica disponible sobre riesgos. ¿En cuál de estas fuentes debiera confiar más?

\begin{enumerate}[label=\alph*)]
  \item Ficheros de proyectos
  \item Informes comparativos entre proyectos publicados
  \item El conocimiento documentado por el equipo del proyecto
  \item Las bases de datos de lecciones aprendidas
\end{enumerate}

\subsection*{Pregunta 332}
El director de proyectos está asignando los estimados de costos generales a las actividades individuales, con el fin de establecer una línea base para la medición del desempeño del proyecto. ¿En qué proceso estamos?

\begin{enumerate}[label=\alph*)]
  \item Gestión de los Costos
  \item Estimar los Costos
  \item \textbf{Determinar el Presupuesto}
  \item Controlar los Costos
\end{enumerate}

\subsection*{Pregunta 333}
El director de proyectos está evaluando las necesidades en cuanto a recursos del proyecto y las lecciones aprendidas de proyectos anteriores . Esta información ocasiona que el director de proyectos se preocupe de si será posible adquirir suficientes recursos para el proyecto en seis meses. ¿Cuál de las siguientes constituye la acción preventiva MENOS efectiva?

\begin{enumerate}[label=\alph*)]
  \item Asegurarse de que los gerentes funcionales tengan una copia del histograma de recursos .
  \item Mostrarle los datos al patrocinador y explicarle la preocupación del director de proyectos .
  \item \textbf{Determinar las métricas que se usarán como alerta temprana de que los recursos no estarán disponibles.}
  \item Pedirle su opinión a los gerentes funcionales.
\end{enumerate}

\subsection*{Pregunta 334}
El director de proyectos está intentando recordar el método de comunicación preferido de un Grupo de Interés. ¿Dónde puede encontrar esa información? : — (5) a. Plan de gestión de los Grupos de Interés

\begin{enumerate}[label=\alph*)]
  \item \textbf{Matriz de evaluación del compromiso de los Grupos de Interés Cc. Diagrama RACI}
  \item Plan de gestión de los recursos humanos
\end{enumerate}

\subsection*{Pregunta 335}
El director del equipo de Project managers recibe un informe en el que el se detecta un retraso en plazo de un 10\% y un sobrecoste de 15\% respecto a lo previsto. El director deberá:

\begin{enumerate}[label=\alph*)]
  \item Cambiar al Project manager del proyecto
  \item Convocar una reunión urgente del equipo de trabajo y partes interesadas directamente implicadas en la ejecución del proyecto
  \item Preguntar al Project manager si estos datos indican una tendencia o son producto de una causa puntual
  \item Verificar por su cuenta que los datos de base usados para el informe son correctos
\end{enumerate}

\subsection*{Pregunta 336}
El director del proyecto anterior de tu proyecto, lo dirigió sin mucha organización. Existe una falta de control de la gestión y no hay entregables claramente definidos en el proyecto. ¿Cuál de las siguientes sería la MEJOR opción para mejorar la organización de tu proyecto?

\begin{enumerate}[label=\alph*)]
  \item Adoptar un enfoque del ciclo de vida para el proyecto.
  \item \textbf{Desarrollar las lecciones aprendidas para cada fase.}
  \item Desarrollar los planes de trabajo específicos para cada paquete de trabajo .
  \item Desarrollar una descripción del producto del proyecto.
\end{enumerate}

\subsection*{Pregunta 337}
El director del proyecto debe ser un comunicador efectivo para asegurar que los interesados del proyecto reciben y entienden la información referente al proyecto y su estado. Antes de distribuir la información a los interesados del proyecto el director del proyecto debe intentar:

\begin{enumerate}[label=\alph*)]
  \item Entender las características de los interesados que afectan a sus necesidades de comunicación, elaborando un registro de interesados
  \item Identificar únicamente aquellos interesados que cuenten con la misma experiencia previa como director de proyecto
  \item Restringir la información para incluir únicamente detalles técnicos específicos
  \item \textbf{Filtrar la información para eliminar detalles innecesarios}
\end{enumerate}

\subsection*{Pregunta 338}
El director del proyecto está trabajando para describir claramente el nivel de participación esperado de todos los involucrados en el proyecto, a fin de prevenir el reproceso, el conflicto y los problemas de coordinación. ¿Cuál de los siguientes describe MEJOR los esfuerzos del director del proyecto?

\begin{enumerate}[label=\alph*)]
  \item Desarrollar el Plan para la Dirección del Proyecto y Planificar la Gestión de la Calidad
  \item Gestionar el Compromiso de los Interesados y Dirigir y Gestionar el Trabajo del Proyecto
  \item Validar el Alcance y Controlar la Calidad
  \item \textbf{Identificar los Riesgos y Desarrollar el Equipo del Proyecto}
\end{enumerate}

\subsection*{Pregunta 339}
Como el director del proyecto, estás preparando tu plan de gestión de la calidad. Estás buscando una herramienta que pueda demostrar la relación entre los eventos y sus efectos resultantes. Quieres utilizar esta herramienta para ilustrar los eventos que causan un efecto negativo en la calidad. ¿Cuál de las siguientes es la MEJOR opción para lograr tu objetivo?

\begin{enumerate}[label=\alph*)]
  \item Histograma
  \item \textbf{Diagrama de Pareto}
  \item Diagrama de Ishikawa
  \item Diagrama de control
\end{enumerate}

\subsection*{Pregunta 340}
El Documento del PMI® sobre Ética profesional e individual de un Project Manager describe las obligaciones que, en la gestión de un proyecto, se deben contraer. Estas obligaciones se refieren a las siguientes categorías: Las responsabilidades para con la profesión, el cliente,… a.

\begin{enumerate}[label=\alph*)]
  \item Informar con veracidad
  \item Informar de las afiliaciones y/o asociaciones
  \item Respeto a otras a otras realidades socioculturales
\end{enumerate}

\subsection*{Pregunta 341}
El documento más completo para que el equipo de proyecto realice una buena integración de las actividades es:

\begin{enumerate}[label=\alph*)]
  \item el acta de constitución
  \item el plan para la dirección del proyecto
  \item el plan de gestión de fases del proyecto
  \item el procedimiento de Control de cambios del alcance
\end{enumerate}

\subsection*{Pregunta 342}
El efecto aureola se refiere a la tendencia a:

\begin{enumerate}[label=\alph*)]
  \item \textbf{Promover desde adentro .}
  \item Contratar a los mejores.
  \item Colocar a las personas en la dirección de proyectos porque son buenas en sus áreas técnicas.
  \item Colocar a las personas en la dirección de proyectos porque han tenido capacitación en dirección de proyectos.
\end{enumerate}

\subsection*{Pregunta 343}
El equipo de proyecto, tras mucho trabajo, ha definido finalmente los requerimientos para un cambio relevante en el alcance del trabajo. Teniendo en cuenta lo dicho, ¿qué es lo primero que debiera hacer el Project manager?

\begin{enumerate}[label=\alph*)]
  \item Notificar a las partes interesadas sobre el nuevo alcance
  \item Buscar el refrendo y apoyo del cambio propuesto
  \item Calcular los riesgos asociados con este cambio
\end{enumerate}

\subsection*{Pregunta 344}
El equipo del proyecto acaba de completar el cronograma inicial del proyecto y el presupuesto. Lo SIGUIENTE que debes hacer es:

\begin{enumerate}[label=\alph*)]
  \item \textbf{Identificar los riesgos.}
  \item Comenzar iteraciones.
  \item Determinar los requisitos de comunicaciones.
  \item Crear un diagrama de barras (de Gantt).
\end{enumerate}

\subsection*{Pregunta 345}
El índice de desempeño del costo (CPI) de un proyecto es 0,6 y el índice de desempeño del cronograma (SPI) es 0,71. El proyecto tiene 625 paquetes de trabajo y se está completando a lo largo de un período de cuatro años. Los miembros del equipo tienen muy poca experiencia y el proyecto recibió poco apoyo en cuanto a su adecuada planificación. ¿Cuál de las siguientes es la MEJOR opción a realizar?

\begin{enumerate}[label=\alph*)]
  \item Actualizar el análisis e identificación del riesgo.
  \item Dedicar más tiempo a mejorar los estimados de costo.
  \item Eliminar la mayor cantidad de paquetes de trabajo que sea posible.
  \item \textbf{Reorganizar la matriz de asignación de responsabilidades .}
\end{enumerate}

\subsection*{Pregunta 346}
El índice de desempeño del costo (CPI) en el proyecto es 1,13, y la relación costo beneficio es 1,2. El alcance del proyecto lo crearon el equipo y los interesados. A lo largo del proyecto, los requisitos del proyecto han estado cambiando. Independientemente de lo que el director de proyectos ha tratado de lograr dirigiendo el proyecto, ¿qué es lo MÁS probable que él tendrá que enfrentar en el futuro?

\begin{enumerate}[label=\alph*)]
  \item Tener que reducir los costos en el proyecto e incrementar los beneficios
  \item Asegurarse de que el cliente apruebe el alcance del proyecto
  \item No ser capaz de medir la finalización del producto del proyecto
  \item Tener que añadir recursos al proyecto
\end{enumerate}

\subsection*{Pregunta 347}
El índice de desempeño del costo del proyecto (CPI) es de 1,02, la relación costo beneficio es de 1,7 y la más reciente ronda de revisiones del desempeño descubrió que se necesitan pocos ajustes. El equipo del proyecto fue reubicado (colocated) en un nuevo edificio cuando se inició el proyecto. Todos comentaron lo emocionados que estaban de tener acceso a espacios completamente nuevos. El patrocinador está brindando apoyo adecuado para el proyecto y son pocos los riesgos no identificados que se han presentado. En un intento por mejorar el desempeño, el director del proyecto gasta parte del presupuesto del proyecto en sillas nuevas para los miembros del equipo y agrega el término "senior" al nombre del puesto de cada miembro del equipo. ¿Qué es lo MÁS correcto que puede decirse acerca de este proyecto o director de proyectos?

\begin{enumerate}[label=\alph*)]
  \item Que el director ha entendido mal la teoría de Herzberg .
  \item El proyecto lentamente está gastando más dinero del que debería . El director de proyectos debería comenzar a tener más cuidado con los costos.
  \item \textbf{La revisión del desempeño debe manejarse con más cuidado para encontrar más ajustes.}
  \item El director de proyectos debe usar el buen juicio para determinar qué variaciones son importantes. 373
\end{enumerate}

\subsection*{Pregunta 348}
El MEJOR momento para asignar un proyecto al director del proyecto es durante:

\begin{enumerate}[label=\alph*)]
  \item Integración .
  \item Selección de proyectos.
  \item Iniciación .
  \item Planificación .
\end{enumerate}

\subsection*{Pregunta 349}
¿El nivel de autoridad de un director de proyecto es mayor cuando gestiona un proyecto en cuál de la siguientes estructuras organizacionales?

\begin{enumerate}[label=\alph*)]
  \item Matriz fuerte
  \item Matriz equilibrada
  \item \textbf{Orientada a proyectos}
  \item Funcional
\end{enumerate}

\subsection*{Pregunta 350}
El nuevo proyecto de instalación de software está en marcha. El director de proyectos está trabajando con el departamento de aseguramiento de la calidad para mejorar la confianza de los interesados de que el proyecto satisfará los estándares de calidad. ¿Cuál de los siguientes DEBEN tener antes de que inicien este proceso?

\begin{enumerate}[label=\alph*)]
  \item Problemas de calidad
  \item Mejora de calidad
  \item Mediciones de control de calidad
  \item \textbf{Reproceso 326}
\end{enumerate}

\subsection*{Pregunta 351}
El nuevo proyecto resulta emocionante tanto para el director de proyectos como para el equipo. Se trata de la primera función del director de proyectos como tal. El equipo siente que serán capaces de completar trabajo que nunca se ha intentado. Hay 29 personas que contribuyen a la descripción del producto y el equipo está conformado por 9 expertos con experiencia en su campo. disienten en cuanto al alcance de dos de los entregab les. Uno de ellos tiene problemas con la versión preliminar de la EDT, pues dice que hace falta agregarle dos paquetes adiciona les de trabajo . Otro de ellos dice que uno de los paquetes de trabajo ni siquiera debería llevarse a cabo. El tercero está de acuerdo con ambos. ¿Cuál sería la MEJOR manera en la que el director de proyectos debería enfrentar este conflicto?

\begin{enumerate}[label=\alph*)]
  \item Debería escuchar las diferentes opiniones , determinar la mejor opción e implementar esa opción.
  \item Debería posponer las discusiones, reunirse con cada individuo y determinar el mejor enfoque .
  \item Debería escuchar las diferencias de opinión, fomentar las discusiones lógicas y facilitar un acuerdo.
  \item \textbf{Debería ayudar al equipo a enfocarse en los aspectos de sus opiniones en los que concuerdan y construir unidad por medio de técnicas de relajación y de formación de equipo con un enfoque común.}
\end{enumerate}

\subsection*{Pregunta 352}
El patrocinador del proyecto acaba de firmar el acta de constitución del proyecto. ¿Qué es lo SIGUIENTE que debemos hacer?

\begin{enumerate}[label=\alph*)]
  \item \textbf{Comenzar a completar los paquetes de trabajo.}
  \item Validar el alcance.
  \item Comenzar el control integrado de cambios.
  \item Empezar a crear los planes de gestión. 100
\end{enumerate}

\subsection*{Pregunta 353}
El plan de gestión de las adquisiciones del proyecto contiene procesos en los siguientes grupos:

\begin{enumerate}[label=\alph*)]
  \item Inicio, planificación, seguimiento y control y cierre
  \item planificación, ejecución, seguimiento y control y cierre
  \item inicio, planificación, ejecución y control y cierre
  \item inicio, planificación, ejecución, supervisión y control y cierre.
\end{enumerate}

\subsection*{Pregunta 354}
El Plan para la dirección del proyecto, las métricas de calidad, la Información sobre el Desempeño del Trabajo y las mediciones del control de calidad son Entradas para qué proceso:

\begin{enumerate}[label=\alph*)]
  \item Dirigir y Gestionar la ejecución del proyecto
  \item Realizar el control de calidad
  \item Planificar la calidad
  \item Realizar el aseguramiento de la calidad
\end{enumerate}

\subsection*{Pregunta 355}
El presupuesto Estimación a la conclusión (Estimate to complete, ETC) permite predecir el coste total del proyecto más probable teniendo en cuenta el desempeño hasta la fecha de análisis. ¿cuál de las siguientes fórmulas no permiten calcularlo?

\begin{enumerate}[label=\alph*)]
  \item EAC=AC+BAC+PV
  \item EAC=AC+ETC
  \item EAC=AC+BAC-EV
  \item EAC=(AC+(BAC-EV)/CPI)
\end{enumerate}

\subsection*{Pregunta 356}
El presupuesto estimado a su finalización (Estimate at Completion, EAC) es una evaluación periódica de:

\begin{enumerate}[label=\alph*)]
  \item el presupuesto actualizado con las desviaciones hasta la fecha.
  \item el presupuesto final aprobado inicialmente
  \item el valor del trabajo realizado a final de obra
  \item el coste del trabajo realizado hasta la fecha
\end{enumerate}

\subsection*{Pregunta 357}
El PRINCIPAL objetivo de escribir una historia de usuarios es:

\begin{enumerate}[label=\alph*)]
  \item Documentar las características o funciones requeridas por los interesados.
  \item \textbf{Crear un registro de las polémicas enfrentadas en el proyecto .}
  \item Realizar el análisis "qué pasaría si".
  \item Comunicar el avance.
\end{enumerate}

\subsection*{Pregunta 358}
El proceso de Adquirir el equipo del Proyecto forma parte del grupo de procesos de:

\begin{enumerate}[label=\alph*)]
  \item Inicio
  \item Planificación
  \item Ejecución
  \item Supervisión y Control
\end{enumerate}

\subsection*{Pregunta 359}
El proceso de desarrollar el equipo, dentro del grupo de proceso de gestión de los recursos humanos consiste en:

\begin{enumerate}[label=\alph*)]
  \item Descripción de información sobre aspectos como responsabilidades, autoridad, competencias y requisitos.
  \item \textbf{Mejorar las competencias, interacción de los miembros del equipo y el ambiente para mejorar el desempeño}
  \item Adquirir el conjunto de recursos necesarios para constituir el equipo de proyecto y la consecución de los objetivos comunes
  \item Asegurar que los recursos asignados están disponibles tal y como se planificó, monitorizando utilización planificada frente a la real y tomar acciones correctivas cuando sea necesario
\end{enumerate}

\subsection*{Pregunta 360}
El proceso de Gestión de las Comunicaciones está relacionado principalmente con:

\begin{enumerate}[label=\alph*)]
  \item Las estimaciones necesarias para recopilar la información
  \item El diseño del informe
  \item \textbf{Hacer que la información necesaria esté disponible para los interesados del proyecto}
  \item A qué fase del proyecto se refiere la información
\end{enumerate}

\subsection*{Pregunta 361}
El proceso "Identificación de los Grupos de Interés":

\begin{enumerate}[label=\alph*)]
  \item Debe ser monitorizado regularmente
  \item Corresponde al grupo de procesos de Planificación
  \item Es muy relevante pero no es un proceso de Dirección de Proyectos
  \item \textbf{Prácticamente debe comenzar al principio de un proyecto}
\end{enumerate}

\subsection*{Pregunta 362}
El programa se planificó años atrás, antes de que hubiera una introducción masiva de nueva tecnología. Mientras se planifica el siguiente proyecto en este programa, el director del proyecto ha ampliado el plan de gestión del alcance del proyecto porque, conforme un proyecto se vuelve más complejo, el nivel de incertidumbre en el alcance:

\begin{enumerate}[label=\alph*)]
  \item Permanece igual.
  \item Disminuye.
  \item Disminuye y luego aumenta.
  \item \textbf{Aumenta. 186}
\end{enumerate}

\subsection*{Pregunta 363}
El proyecto A tiene una tasa interna de retorno (IRR) de 21 por ciento. El proyecto B tiene una tasa interna de retorno de 7 por ciento. El proyecto C tiene una tasa interna de retorno de 31 por ciento. El proyecto D tiene una tasa interna de retorno de 19 por ciento. ¿Cuál de estos sería el MEJOR proyecto?

\begin{enumerate}[label=\alph*)]
  \item El proyecto A
  \item El proyecto B
  \item El proyecto C
  \item El proyecto D
\end{enumerate}

\subsection*{Pregunta 364}
El proyecto apenas empieza y está conformado por personas de 14 departamentos diferentes. El acta de constitución del proyecto fue firmada por una persona y contiene más de 30 requisitos importantes que deben cumplirse en el proyecto. El patrocinador le ha informado al director de proyectos que el índice de desempeño del cronograma (SPI) debe mantenerse entre 0,95 y 1,1. Unos cuantos minutos de investigación llevaron a la identificación de 34 interesados, y los objetivos del cronograma del proyecto han sido restringidos. El director de proyectos acaba de ser contratado . ¿Cuál de los siguientes tipos de poderes de dirección de proyectos será el que MEJOR ayude al director de proyectos a obtener la cooperación de los demás?

\begin{enumerate}[label=\alph*)]
  \item Formal
  \item Referente
  \item Penalidad
  \item \textbf{Experto 371}
\end{enumerate}

\subsection*{Pregunta 365}
Como el proyecto en el que trabajas esta fuera de costes, tu director te indica que cargues las horas de proyecto ejecutadas a los códigos de otro proyecto. Como Project manager deberías:

\begin{enumerate}[label=\alph*)]
  \item Informar a los auditores de la compañía
  \item Hacer caso a tu superior
  \item Solicitar una orden escrita de la acción propuesta por parte del director y notificarlo al Project manager del otro proyecto.
  \item Cerrar el Proyecto si es posible ya que no se puede perder dinero
\end{enumerate}

\subsection*{Pregunta 366}
El proyecto es temporal y singular La identificación efectiva de los riesgos depende en gran medida de la habilidad del director de proyecto para emplear varias técnicas diferentes. ¿Cuál de las siguientes técnicas es útil en la identificación de riesgos?

\begin{enumerate}[label=\alph*)]
  \item Diagramas de dispersión
  \item \textbf{Análisis de causa raíz}
  \item Matriz de Probabilidad-Impacto
  \item Estructura de desglose de los recursos
\end{enumerate}

\subsection*{Pregunta 367}
El proyecto está casi finalizado. El proyecto tiene una variación del cronograma de 300 y una variación de costos de -900. Todas las inspecciones de control de calidad se han finalizado, salvo una, y todas han alcanzado los requisitos de calidad. Todas las acciones en el registro de polémicas se han resuelto. Muchos de los recursos se han liberado. El patrocinador está a punto de llamar a una reunión para obtener la validación del producto cuando el cliente notifica al director del proyecto que quieren hacer un cambio importante en el alcance. El director del proyecto debe:

\begin{enumerate}[label=\alph*)]
  \item Reunirse con el equipo del proyecto para determinar si se puede hacer este cambio.
  \item Pedir al cliente una descripción del cambio.
  \item Explicar que no se puede hacer el cambio en ese momento del proceso.
  \item \textbf{Informar a la gerencia. 187}
\end{enumerate}

\subsection*{Pregunta 368}
El proyecto ha ido bien, excepto por el número de cambios que se han realizado. El producto del proyecto está instalándose dentro de siete departamentos diferentes dentro de la compañía y cuando esté operando mejorará enormemente el desempeño departamental. El equipo ha seleccionado los procesos adecuados para utilizar en el proyecto . El director del proyecto es un experto técnico y también ha sido capacitado en comunicaciones y dirección de personal. ¿Cuál de las siguientes es la causa MÁS probable de los problemas del proyecto?

\begin{enumerate}[label=\alph*)]
  \item El director del proyecto no estaba capacitado para entender el entorno de la compañía.
  \item El proyecto debió tener mayor supervisión de dirección, puesto que se traducirá en grandes beneficios para la compañía.
  \item El proyecto debió utilizar más de los procesos de la dirección de proyectos .
  \item No se identificó a algunos de los interesados .
\end{enumerate}

\subsection*{Pregunta 369}
El proyecto iba bien y, de repente, hubo cambios al proyecto que vinieron de múltiples interesados . Después de que todos los cambios fueron determinado s, el director del proyecto se reunió con todos los interesados para investigar por qué hubo cambios y para descubrir cualquier otro cambio. El trabajo del proyecto se ha estabilizado pero un miembro del equipo le menciona al pasar al director del proyecto que agregó funcionalidad a un producto del proyecto. "No te preocupes", dice, "¡No afecté el tiempo, el costo ni la calidad!" ¿Qué debe hacer PRIMERO el director del proyecto?

\begin{enumerate}[label=\alph*)]
  \item Preguntar al miembro del equipo cómo se determinó la necesidad de la funcionalidad.
  \item Realizar una reunión para revisar el trabajo completo del miembro del equipo .
  \item Buscar otra funcionalidad agregada.
  \item Preguntar al miembro del equipo cómo sabe que no hay impacto sobre el tiempo, los costos o la calidad.
\end{enumerate}

\subsection*{Pregunta 370}
El proyecto no finaliza hasta que

\begin{enumerate}[label=\alph*)]
  \item El alcance del producto está completo, el cierre administrativo está completo y se ha recibido el pago acordado.
  \item Se recibe la aceptación formal y cualquier otro requerimiento para el cierre del proyecto establecido en el contrato se ha cumplido
  \item El cliente está satisfecho con el resultado y se ha recibido la certificación final.
  \item Las lecciones aprendidas se han documentado.
\end{enumerate}

\subsection*{Pregunta 371}
El proyecto presentó algunos problemas , aunque ahora parece estar bajo control. En los últimos meses se ha utilizado casi toda la reserva y la mayor parte de los impactos negativos de los eventos que se predijeron han ocurrido . Sólo faltan cuatro actividades, y dos de ellas aparecen en la ruta crítica. La gerencia le ha informado al director del proyecto que lo más conveniente para la organización ejecutante sería terminar el proyecto dos semanas antes de lo estipulado, a fin de obtener una ganancia adicional. Como respuesta, el director del proyecto envía una solicitud de propuesta para una parte del trabajo que el equipo tenía previsto hacer, con la esperanza de encontrar otra compañía que sea capaz de realizar el trabajo más rápido. La opción que MEJOR representa aquello con lo que el director del proyecto intenta trabajar es:

\begin{enumerate}[label=\alph*)]
  \item Reserva.
  \item Oportunidades.
  \item Validación del alcance .
  \item Amenazas.
\end{enumerate}

\subsection*{Pregunta 372}
El proyecto que estás gestionando tiene un TCPI de 1,25. qué es lo mejor que puedes hacer?

\begin{enumerate}[label=\alph*)]
  \item Puedes ir tranquilo, tu gestión de costos del proyecto está siendo muy buena
  \item Gestionar los contratos restantes y las horas imputadas al proyecto de forma estricta
  \item Rehacer el cronograma con una nueva línea de base de plazos
  \item Crear un nuevo presupuesto que represente la nueva realidad
\end{enumerate}

\subsection*{Pregunta 373}
El proyecto que usted dirige parece haber entrado en una fase difícil. El equipo de 7 personas que forman ustedes necesita ahora de apoyo por parte de 3 consultores externos que clarifiquen hacia dónde va el proyecto. Con esta nueva incorporación ¿cuántos nuevos canales de comunicación se añaden el proyecto? :

\begin{enumerate}[label=\alph*)]
  \item 21 ¿Qué es un programa? :
  \item Una regulación gubernamental beneficios.
  \item Una iniciativa empezada por la dirección
  \item Un grupo de proyectos sin relación entre sí que se dirigen de forma coordinada
\end{enumerate}

\subsection*{Pregunta 374}
El proyecto se calculó para finalizar cuatro días después de la fecha de conclusión deseada. No tienes acceso a recursos adicionales. El proyecto es de bajo riesgo, se espera que la relación costo beneficio sea de 1,6, y las dependencias son preferenciales . Bajo estas circunstancias, ¿cuál es la MEJOR cosa por hacer?

\begin{enumerate}[label=\alph*)]
  \item Reducir recursos de una actividad.
  \item Hacer más actividades simultáneas .
  \item Mover recursos de las dependencias preferenciales a las dependencias externas .
  \item Quitar una actividad del proyecto.
\end{enumerate}

\subsection*{Pregunta 375}
El proyecto va retrasado en casi un mes y el sobrecoste estimado a fecha de hoy es del 15\%. El Project manager necesita abordar estos problemas y uno de los primeros elementos que utiliza es la Estructura de Desglose de Tareas. ¿Para qué puede utilizarla?

\begin{enumerate}[label=\alph*)]
  \item Para comunicarse e informar al cliente
  \item Para mostrar las fechas previstas y reales del proyecto
  \item Para identificar las causas que han producido el retraso y el sobrecoste
  \item Para identificar a los jefes funcionales de las tareas que están retrasadas.
\end{enumerate}

\subsection*{Pregunta 376}
El punto más alto de la jerarquía de necesidades de Maslow es:

\begin{enumerate}[label=\alph*)]
  \item \textbf{Satisfacción fisiológica.}
  \item Obtención de la supervivencia .
  \item Necesidad de asociación .
  \item Estima .
\end{enumerate}

\subsection*{Pregunta 377}
¿el registro de riesgos no es una entrada de cuál de los siguientes procesos?

\begin{enumerate}[label=\alph*)]
  \item Identificar riesgos
  \item Evaluar cuantitativamente los riesgos
  \item Evaluar cualitativamente los riesgos
  \item Control de los riesgos
\end{enumerate}

\subsection*{Pregunta 378}
El riesgo de costos significa:

\begin{enumerate}[label=\alph*)]
  \item Hay riesgos que le costarán dinero al proyecto.
  \item El proyecto es demasiado riesgoso desde una perspectiva de costos.
  \item \textbf{Existe un riesgo de que los costos del proyecto puedan elevarse más de lo planificado.}
  \item Existe un riesgo de que el costo del proyecto sea menor de lo planificado .
\end{enumerate}

\subsection*{Pregunta 379}
El Servicio de atención al cliente ha mostrado su preocupación sobre un proyecto en el que estás implicado. Parece ser que si el proyecto se ejecuta tal y como estaba planificado, deberán comprar equipos adicionales para equipar el centro de trabajo del cliente. Este coste no estaba previsto en el presupuesto del proyecto. El patrocinador del proyecto insiste en que el proyecto debe seguir adelante o el cliente se molestará. ¿Qué respuesta es la adecuada?

\begin{enumerate}[label=\alph*)]
  \item El responsable del Servicio al Cliente tiene razón. Si el coste no había sido considerado al inicio del proyecto, el proyecto no puede ejecutarse tal y como había sido planificado.
  \item El conflicto debe resolverse a favor del cliente mediante procesos de cambio establecidos
  \item Se debe actuar conforme a lo establecido por el patrocinador del proyecto
  \item El conflicto debe resolverse a favor del departamento de Servicio al Cliente
\end{enumerate}

\subsection*{Pregunta 380}
El trabajo de un proyecto está llevándose a cabo cuando el director del proyecto escucha a dos trabajadores discutiendo sobre lo que significa un conjunto de instrucciones . El director de proyectos investiga y descubre que las instrucciones para la construcción de los cimientos de concreto que se están colocando fueron mal traducidas entre los diferentes idiomas en uso en el proyecto. ¿Cuál de los siguientes es lo MEJOR que el director de proyectos puede hacer PRIMERO?

\begin{enumerate}[label=\alph*)]
  \item Hacer que las instrucciones sean traducidas por un grupo más experimentado.
  \item Buscar los impactos en la calidad de la mala traducción de las instrucciones para los cimientos.
  \item Llevar el incidente a la atención del equipo y pedirles que busquen más problemas de traducción .
  \item \textbf{Informar al patrocinador del problema en el siguiente informe del proyecto. 328}
\end{enumerate}

\subsection*{Pregunta 381}
El trabajo operativo se distingue del trabajo de proyecto debido a que el trabajo operativo es:

\begin{enumerate}[label=\alph*)]
  \item Único.
  \item Temporal.
  \item Continuo y repetitivo .
  \item Parte de todas las actividades del proyecto.
\end{enumerate}

\subsection*{Pregunta 382}
Elaborar los informes con toda la información recogida sobre el avance del proyecto para que las partes interesadas estén informadas corresponde al proceso de

\begin{enumerate}[label=\alph*)]
  \item Planificar las comunicaciones
  \item Distribuir la información
  \item Gestionar las expectativas de los interesados
  \item Informar del desempeño
\end{enumerate}

\subsection*{Pregunta 383}
En una conversación con otro director de proyectos de su organización, le comenta que algunos recursos asignados a su proyecto han sido movilizados a otro proyecto con mayor prioridad para la organización. Como resultado, un camino casi crítico se ha convertido en crítico y varios de los recursos restantes del equipo del proyecto se encuentran sobre asignados y, para solucionar esta situación, está valorando utilizar la nivelación de recursos. ¿Cuál sería la principal advertencia que le podría hacer el director del proyectos a su colega sobre la nivelación de recursos?

\begin{enumerate}[label=\alph*)]
  \item \textbf{Que la nivelación de recursos a menudo provoca cambios en el camino crítico del proyecto, por lo que podrían aparecer uno o varios nuevos caminos críticos}
  \item Que los recursos que se nivelan tengan la misma capacitación y experiencia que los recursos que han sido movilizados porque, de lo contrario, se modificará el camino crítico
  \item Que en lugar de utilizar la nivelación de recursos debería utilizar adelantos y retrasos entre actividades donde sea posible
  \item Que previamente es necesario analizar la relación coste/recurso para determinar si es posible utilizar la nivelación de recursos
\end{enumerate}

\subsection*{Pregunta 384}
¿En cuál de los siguientes casos subcontratarías paquetes de trabajo en vez de realizarlos con el personal propio?

\begin{enumerate}[label=\alph*)]
  \item La compañía tiene recursos suficientes y es capaz de realizarlos
  \item Los productos son confidenciales en lo que tiene que ver con diseño o tecnología aplicada.
  \item Cuando hay otras empresas que pueden realizarlo a un precio más barato y en menor tiempo pero con peor calidad
  \item La compañía no tiene recursos suficientes y es incapaz de realizarlos
\end{enumerate}

\subsection*{Pregunta 385}
¿En cuál de los siguientes costes no se centrará el control de costos de un Project manager?

\begin{enumerate}[label=\alph*)]
  \item el control de las horas dedicadas a cada actividad
  \item los gastos generales imputados al proyecto
  \item los gastos incurridos en materiales y equipos
  \item las adquisiciones realizadas a terceros
\end{enumerate}

\subsection*{Pregunta 386}
¿En cuántos grupos de procesos encontramos procesos relativos al área de conocimiento de gestión de la calidad?

\begin{enumerate}[label=\alph*)]
  \item En los 5 grupos
  \item En 4 grupos
  \item En 3 grupos
  \item En 2 grupos
\end{enumerate}

\subsection*{Pregunta 387}
En una discusión con un proveedor relativa a la certificación final de la obra, éste pierde los papeles y empieza a gritar e insultar al Project manager. ¿Qué es lo mejor que puede hacer el Project manager?

\begin{enumerate}[label=\alph*)]
  \item De manera educada dar por concluida la reunión
  \item Calmarle al proveedor indicándole que sus exigencias serán aceptadas
  \item Indicarle al proveedor de forma correcta que su actitud no es adecuada y que, una vez se haya calmado, se reanude la reunión para tratar de buscar una solución de consenso.
  \item Registrar en el libro de incidencias que no se vuelva a contratar a esta empresa
\end{enumerate}

\subsection*{Pregunta 388}
En el ciclo conocido como Ciclo de Deming cuál es la secuencia correcta

\begin{enumerate}[label=\alph*)]
  \item Planificar, Desarrollar, Actuar, Controlar
  \item Iniciar, Planificar, Ejecutar y Supervisar
  \item Planificar, Desarrollar, Controlar, Actuar
  \item Planificar, Ejecutar, Controlar y Actuar
\end{enumerate}

\subsection*{Pregunta 389}
En el proceso de cierre de un contrato, el Project manager es responsable de que se cierren todos los contratos antes de cerrar administrativamente el proyecto. En este contexto, la negociación para la resolución de conflictos pendientes, reclamaciones y disputas se hace mediante

\begin{enumerate}[label=\alph*)]
  \item El cierre del contrato
  \item Las auditorías de compras
  \item La documentación de reclamaciones
  \item Acuerdos negociados
\end{enumerate}

\subsection*{Pregunta 390}
En el proceso de planificación de la calidad estás utilizando el benchmarking que consiste todo menos en:

\begin{enumerate}[label=\alph*)]
  \item Se hace únicamente dentro de la propia organización
  \item Genera ideas para mejorar y da una base para la medida del desempeño
  \item Puede compararse con prácticas de la misma o de otra área
\end{enumerate}

\subsection*{Pregunta 391}
En el proyecto de elaboración de manuales de formación de un conocido programa informático habéis tenido que rehacer partes del manual ya escrito y reescribir dos capítulos completos ya que no cumplían con los parámetros de diseño, enfoque y edición solicitados por el cliente. ¿Cuál de estas afirmaciones consideras correcta?

\begin{enumerate}[label=\alph*)]
  \item El plan de gestión de riesgos ayuda a resolver estos problemas.
  \item El plan de aseguramiento de la calidad funcionó correctamente lo que hizo que se detectaran esos errores.
  \item La identificación de actividades no se hizo de forma correcta
  \item Todos estos problemas surgieron por un mal Enunciado del Trabajo
\end{enumerate}

\subsection*{Pregunta 392}
En la estimación de la duración de las actividades, ¿ cuándo es más aconsejable utilizar la comparación análoga?

\begin{enumerate}[label=\alph*)]
  \item \textbf{Cuando se dispone de poco tiempo de planificación del Proyecto}
  \item Cuando el número de actividades está entre 50 - 100
  \item Cuando el Proyecto tiene pocos Paquetes de Trabajo
  \item Cuando no hay una necesidad de que el grado de precisión sea muy estricto
\end{enumerate}

\subsection*{Pregunta 393}
En ocasiones, el alcance y contenido de un paquete de trabajo puede solaparse con el de otro en una EDP. Cuando ocurre esto: :

\begin{enumerate}[label=\alph*)]
  \item Esto no puede ocurrir, ya que los paquetes de trabajo deben tener un alcance y contenido único
  \item Un mismo recurso debería ocuparse de realizar el trabajo requerido en ambos
  \item Hay que coordinar adecuadamente ambos paquetes de trabajo para no repetir actividades
  \item Hay que definir con detalle los entregables de ambos paquetes de trabajo
\end{enumerate}

\subsection*{Pregunta 394}
En una organización orientada a proyectos, el equipo del proyecto:

\begin{enumerate}[label=\alph*)]
  \item Reporta a muchos jefes.
  \item No tiene lealtad al proyecto.
  \item Reporta al gerente funcional.
  \item No siempre tendrá una "sede''.
\end{enumerate}

\subsection*{Pregunta 395}
¿En qué consiste la línea base del alcance?

\begin{enumerate}[label=\alph*)]
  \item El acta de constitución, el Enunciado del Alcance del producto y la EDT
  \item El Enunciado del Alcance del Proyecto, el Documento de Requisitos y la EDT
  \item La EDT, el diccionario de la EDT y el Enunciado del Trabajo
  \item El plan de Incidentes del producto
\end{enumerate}

\subsection*{Pregunta 396}
¿En qué consiste la pirámide jerárquica de Maslow?

\begin{enumerate}[label=\alph*)]
  \item Una disposición jerárquica de misión, visión objetivos estratégicos y portfolio de proyectos de una compañía
  \item Una pirámide de expectativas y necesidades humanas
  \item Una representación piramidal del crecimiento profesional del equipo
  \item Un método para definir reconocimientos y remuneraciones
\end{enumerate}

\subsection*{Pregunta 397}
¿En qué documento del proyecto define Ud. los criterios de éxito y requisitos de aprobación?

\begin{enumerate}[label=\alph*)]
  \item El plan de calidad.
  \item \textbf{El acta de constitución del proyecto.}
  \item La documentación de los requisitos.
  \item El enunciado del alcance.
\end{enumerate}

\subsection*{Pregunta 398}
¿En qué grupo de procesos de dirección de proyectos se crea el acta de constitución del proyecto?

\begin{enumerate}[label=\alph*)]
  \item \textbf{Ejecución}
  \item Planificación
  \item Cierre
  \item Iniciación
\end{enumerate}

\subsection*{Pregunta 399}
¿En qué grupo de procesos de dirección de proyectos se puede identificar a los interesados?

\begin{enumerate}[label=\alph*)]
  \item Iniciación , planificación, ejecución y seguimiento y control
  \item Iniciación y planificación
  \item Planificación y seguimiento y control
  \item \textbf{Seguimiento y control y cierre}
\end{enumerate}

\subsection*{Pregunta 400}
¿En qué grupo de procesos de dirección de proyectos se realiza una estimación aproximada de orden de magnitud?

\begin{enumerate}[label=\alph*)]
  \item Planificación
  \item Cierre
  \item Ejecución
  \item \textbf{Iniciación 281}
\end{enumerate}

\subsection*{Pregunta 401}
¿En qué grupo de procesos de gestión del proyecto se elaboran las predicciones presupuestarias?

\begin{enumerate}[label=\alph*)]
  \item En la planificación
  \item En la fase de inicio
  \item En la supervisión y control
  \item En la de ejecución
\end{enumerate}

\subsection*{Pregunta 402}
En relación a las negociaciones en adquisiciones, ¿Cuál de las siguientes afirmaciones no es correcta?

\begin{enumerate}[label=\alph*)]
  \item Durante la fase de negociación y ejecución de contratos de compras, el Project manager estará implicado pero no necesariamente liderando las negociaciones.
  \item El responsable del departamento de compras designado las dirige y el Project manager se encargará de facilitarle datos técnicos y de gestión
  \item Una vez firmado el contrato no se admitirán más negociaciones
  \item Los aspectos habituales a negociar son el coste, el plazo o el alcance y los términos y condiciones de pago.
\end{enumerate}

\subsection*{Pregunta 403}
En relación al coste de la calidad cuál de los siguientes expertos concluyeron que el 85\% del coste de la calidad es un problema de gestión y no de mala ejecución?

\begin{enumerate}[label=\alph*)]
  \item Crosby
  \item Juran
  \item Deming
  \item Ishikawa
\end{enumerate}

\subsection*{Pregunta 404}
En relación con los cambios, la atención del director del proyecto está MEJOR enfocada en:

\begin{enumerate}[label=\alph*)]
  \item Realizar los cambios.
  \item Seguir y registrar los cambios.
  \item Informar al patrocinador de los cambios.
  \item Prevenir cambios innecesarios.
\end{enumerate}

\subsection*{Pregunta 405}
En una relación inicio-a-inicio entre dos tareas que no se encuentran en la ruta critica, el equipo del proyecto determina que el inicio de la tarea sucesora debe atrasarse 3 días después del inicio de su predecesora. A este atraso se le conoce como \_\_\_\_\_\_\_\_\_\_\_\_\_\_.

\begin{enumerate}[label=\alph*)]
  \item Tiempo de posposición.
  \item \textbf{Holgura libre.}
  \item Holgura total.
  \item Holgura.
\end{enumerate}

\subsection*{Pregunta 406}
En su programacién hay una actividad cuya fecha temprana de comienzo son 6 semanas, la fecha temprana de En su programación hay una actividad cuya fecha temprana de comienzo son 6 semanas, la fecha temprana de finalizacién 10 semanas, la fecha mas tardia de comienzo 14 semanas y la fecha mas tardia de finalizacién 18 finalización 10 semanas, la fecha más tardía de comienzo 14 semanas y la fecha más tardía de finalización 18 semanas. La holgura de esta actividad es: semanas. La holgura de esta actividad es: opción: a,4 semanas a.4 semanas cuestian b, 18 semanas b, 18 semanas Cc. 6 semanas c.6semanas

\begin{enumerate}[label=\alph*)]
  \item Que su Índice de desempeño del cronograma será positivo a. Que su indice de desempefio del cronograma sera positivo
  \item Está por encima del presupuesto b. Esta por encima del presupuesto cuestién
  \item \textbf{Está bajo presupuesto c. Esta bajo presupuesto}
  \item Que su indice de desempefi del cronograma sera negativo d. Que su Índice de desempeño del cronograma será negativo
\end{enumerate}

\subsection*{Pregunta 407}
En tu compañía el Project manager comparte responsabilidad con el jefe funcional o de departamento y asigna prioridades y dirige a los miembros del equipo de proyecto. en este caso, ¿cuál de las siguientes no es responsabilidad única del Project manager?:

\begin{enumerate}[label=\alph*)]
  \item Planificar los recursos
  \item Gestionar costes y presupuestos
  \item Llevar a cabo la valoración del desempeño del equipo
  \item Definir tareas
\end{enumerate}

\subsection*{Pregunta 408}
En un contrato de Precio Fijo ¿qué no será un resultado de la gestión del contrato?

\begin{enumerate}[label=\alph*)]
  \item La aprobación de los recursos destinados
  \item La aceptación de los Entregables y los correspondientes pagos
  \item Un riesgo de fuerza mayor ocurrido durante la ejecución del proyecto
  \item Cambios en el contrato
\end{enumerate}

\subsection*{Pregunta 409}
En un proyecto, un ‘entregable' se entiende como: aún

\begin{enumerate}[label=\alph*)]
  \item Un resultado para dividir el proyecto en elementos de trabajo manejables 1.00
  \item Una condición necesaria sobre el contenido, forma o funcionalidad de un producto o servicio
  \item Un hito o un momento importante durante la planificación y ejecución del proyecto
  \item Un resultado que puede ser medido, tangible o comprobable y que debe ser producido para completar una parte del proyecto
\end{enumerate}

\subsection*{Pregunta 410}
Entre los criterios de selección del tipo de contrato para proceder posteriormente a la selección de proveedores no está:

\begin{enumerate}[label=\alph*)]
  \item El bien que se va a adquirir y el grado de definición del alcance
  \item El nivel de seguimiento y control que el comprador va a llevar a cabo
  \item Al nivel de experiencia y capacidad técnica de los vendedores
  \item Si se quieren establecer o no incentivos
\end{enumerate}

\subsection*{Pregunta 411}
Eres designado como director del proyecto a la mitad del proyecto. El proyecto se encuentra dentro de las líneas base, pero el cliente no está contento con el desempeño del proyecto. ¿Qué debes hacer PRIMERO?

\begin{enumerate}[label=\alph*)]
  \item Analizarlo con el equipo del proyecto.
  \item Recalcular las líneas base.
  \item Renegociar el contrato.
  \item Reunirte con el cliente.
\end{enumerate}

\subsection*{Pregunta 412}
Eres el Project manager de una gran planta petroquímica de la cual has subcontratado buena parte de las instalaciones eléctrico-mecánicas y de control mediante un contrato de costes más honorarios fijos (Cost Plus Fixed Fee, CPFF). Dentro de los términos y condiciones del contrato está que el subcontratista debe enviar informes semanales de avance y estado real de sus trabajos. Sin embargo, el Project manager no ha recibido ninguno en las últimas dos semanas. ¿Qué crees que es lo mejor que puedes hacer?

\begin{enumerate}[label=\alph*)]
  \item Pedir a tu director que lo comunique al director de la otra compañía
  \item Enviar una orden de parada de los trabajos hasta recibir los informes
  \item Enviarles una carta de incumplimiento
  \item Certificar únicamente un 50\% de lo realizado durante ese período
\end{enumerate}

\subsection*{Pregunta 413}
Eres el Project manager de la construcción de un nuevo puente de 60 metros de luz. El objetivo es lograr tenerlo listo en 18 meses, fecha en la que, por haber elecciones autonómicas, debe de estar listo para que las autoridades lo inauguren. Pero los recursos que necesitas no estarán a tu disposición hasta dentro de 5 meses ya que están asignados de momento a otro proyecto. La máquina de pilotaje y su manipulador especialista deben estar en la obra en dos meses. ¿Cuál de las siguientes habilidades utilizarías para lograr que esté equipo esté en tu obra a tiempo?

\begin{enumerate}[label=\alph*)]
  \item Habilidades de negociación y capacidad de influencia
  \item Habilidades de comunicación y organización
  \item Habilidades de comunicación y liderazgo
  \item Habilidades de resolución de problemas
\end{enumerate}

\subsection*{Pregunta 414}
Eres el Project manager de la construcción de un parque eólico, la subestación y línea de alta conexión. Al valorar las compras y adquisiciones, se produce una bajada del precio del cobre que según tu cronograma tenías pensado comprar dentro de 8 meses cuando fueras a hacer el tendido. Ante esta coyuntura, reunidos el Project manager, el departamento de compras y el patrocinador (sponsor) se aprueba adelantar la compra del cable en este momento. ¿Qué estrategia estáis adoptando?

\begin{enumerate}[label=\alph*)]
  \item Aceptar la oportunidad
  \item Mejorar la oportunidad
  \item Explotar la oportunidad
  \item Compartir la oportunidad
\end{enumerate}

\subsection*{Pregunta 415}
Eres el Project manager de un proyecto de puesta en marcha de una nueva planta en China. Por vuestra experiencia, el elemento crítico en términos de plazo es la aprobación urbanística del terreno que, además forma parte del camino crítico. Preguntas a personas del sector y te responden que en el mejor de los casos llevará 6 meses, si se complican las cosas puede llevar hasta 12 meses pero que lo normal últimamente ha sido que en 8 meses estaba ya conseguido el permiso. ¿Cuál será el valor esperado o media ponderada?

\begin{enumerate}[label=\alph*)]
  \item Algo más de 8 meses
  \item 5 meses
  \item Algo más de 2 meses
  \item Algo menos de 10 meses
\end{enumerate}

\subsection*{Pregunta 416}
Eres el director de proyectos de un gran proyecto de instalación cuando te das cuenta de que hay más de 200 interesados potenciales en el proyecto. ¿Cuál de los siguientes sería el MEJOR curso de acción que podrías tomar?

\begin{enumerate}[label=\alph*)]
  \item Eliminar algunos interesados .
  \item Encontrar una manera efectiva de reunir las necesidades de todos los interesados.
  \item Contactarte con tu gerente y preguntarle qué interesados son los más importantes .
  \item \textbf{Reunir las necesidades de todos los interesados más influyentes. 542}
\end{enumerate}

\subsection*{Pregunta 417}
Eres un Project manager muy experimentado técnicamente en la estimación de instalaciones eléctricas en edificios comerciales por ello haces las estimaciones de tiempo y recursos por ti mismo. Cuál de las siguientes respuestas la consideras la más adecuada:

\begin{enumerate}[label=\alph*)]
  \item Tienes experiencia y tienes la autoridad por lo que es correcta tu forma de actuar.
  \item Deberías, si no realizas el trabajo de forma conjunta, al menos dejar que tu equipo revise y confirme tus estimaciones.
  \item Debes consensuar las estimaciones con el promotor del centro comercial
  \item No está considerando tiempos de tramitación que puede ser determinantes
\end{enumerate}

\subsection*{Pregunta 418}
Eres un Project manager que estás tratando de identificar la información que realmente necesita cada parte interesada del proyecto ¿Cuál de los siguientes elementos no son fuentes para recoger esa información según PMBOK?

\begin{enumerate}[label=\alph*)]
  \item Los organigramas de los departamentos y la compañía
  \item Los recursos asignados al proyecto y su ubicación física
  \item La repercusión social que pueda tener el proyecto
  \item Las tecnologías de comunicaciones disponibles del proyecto
\end{enumerate}

\subsection*{Pregunta 419}
Eres un Project manager y estás tratando de conocer bien las tareas a realizar desglosándolas en actividades, estimando el tiempo que puede llevarte ejecutar cada actividad y los recursos necesarios. ¿En qué grupo de procesos te encuentras?

\begin{enumerate}[label=\alph*)]
  \item Planificación
  \item Ejecución
  \item Inicio
  \item Supervisión y Control
\end{enumerate}

\subsection*{Pregunta 420}
Eres un director de proyectos en un proyecto de desarrollo de software de US\$5.000 .000. Mientras trabajas con tu equipo del proyecto para desarrollar un diagrama de red, te das cuenta de una serie de actividades que se pueden trabajar en para lelo pero deben terminar en una secuencia específica. ¿Qué tipo de método de secuencia de actividades se necesita para estas actividades?

\begin{enumerate}[label=\alph*)]
  \item Método de diagramación por precedencia
  \item Método de diagramación con flechas
  \item Método de la ruta crítica
  \item \textbf{Método de diagramación operacional}
\end{enumerate}

\subsection*{Pregunta 421}
Eres un director de proyectos en un proyecto de desarrollo de software de US\$5.000.000. Mientras trabajas con tu equipo del proyecto para desarrollar un diagrama de red, tu arquitecto de datos sugiere que se puede mejorar la calidad si la gerencia senior aprueba el modelo de datos antes de pasar a desarrollo de software. ¿Cuál de los siguientes describe MEJOR este tipo de entrada?

\begin{enumerate}[label=\alph*)]
  \item Dependencia obligatoria
  \item Dependencia discrecional
  \item Dependencia externa
  \item Heurística
\end{enumerate}

\subsection*{Pregunta 422}
Eres un director de proyectos que dirige un equipo de proyectos interfuncional en un ambiente matricial débil. Ninguno de los miembros del equipo de tu proyecto se reporta contigo funcionalmente y careces de la habilidad de recompensar de forma directa su desempeño. El proyecto es difícil pues incluye restricciones de tiempo muy estrictas y estándares de calidad que representan un desafío. ¿Cuál de los siguientes tipos de poderes de la dirección de proyectos será el MÁS efectivo si tomamos en cuenta estas circunstancias?

\begin{enumerate}[label=\alph*)]
  \item Referente
  \item Experto
  \item Penalidad
  \item \textbf{Formal 370}
\end{enumerate}

\subsection*{Pregunta 423}
Es muy importante ver si es la propia forma de trabajar la que está causando los problemas. El Project manager necesita dedicar un tiempo en pensar sobre cómo puede estar tranquilo de que el trabajo se está haciendo bien según unos procedimientos que lograrán un servicio o producto de calidad y que, además, el proceso es el más eficiente se realizan únicamente las actividades que añaden valor al producto o servicio. A qué proceso nos estamos refiriendo?

\begin{enumerate}[label=\alph*)]
  \item aseguramiento de la calidad
  \item control de la calidad
  \item gestionar la calidad
  \item planificar la calidad
\end{enumerate}

\subsection*{Pregunta 424}
Está determinando cómo medirá el rendimiento de su proyecto. Ha identificado varias áreas clave que incluyen el rendimiento de cumplir el tiempo, el control de costos y la frecuencia de defectos. Usted está estableciendo cuál de los siguientes:

\begin{enumerate}[label=\alph*)]
  \item Requisitos técnicos
  \item Métricas
  \item Especificaciones
  \item \textbf{Umbrales de control}
\end{enumerate}

\subsection*{Pregunta 425}
Está en el último mes de un proyecto de 14 meses. Uno de los miembros del equipo de su proyecto le informa que ha identificado un nuevo riesgo potencial que podría representar una amenaza importante para el éxito del proyecto. Usted debería:

\begin{enumerate}[label=\alph*)]
  \item Notificar rápidamente al patrocinador del proyecto.
  \item Implementar su plan de contingencia
  \item Concentrarse en completar el proyecto antes de que ocurra el evento
  \item \textbf{Analizar el riesgo}
\end{enumerate}

\subsection*{Pregunta 426}
Estamos hablando de desarrollar la línea base del desempeño de costos con la que se comparará el avance del presupuesto en la fase de ejecución. ¿De qué proceso estamos hablando?

\begin{enumerate}[label=\alph*)]
  \item Estimación de Costes
  \item Asignación de recursos a actividades
  \item Elaboración del Presupuesto
  \item Control de Costes
\end{enumerate}

\subsection*{Pregunta 427}
Están gestionando la ejecución de la construcción de unas naves industriales y en su secuencia de actividades se sabe que la colocación de los separadores entre estancias debe hacerse antes de colocar el parquet. Sin embargo como el proveedor es el mismo, le han suministrado todo a la vez. ¿Qué relación hay ahora entre la colocación de parquet y separadores?

\begin{enumerate}[label=\alph*)]
  \item Inicio a inicio
  \item Final a final (Finish-to-Finish, FF)
  \item Final a inicio (Finish-to-Start, FS)
  \item existe una dependencia externa
\end{enumerate}

\subsection*{Pregunta 428}
Estás a la mitad de un nuevo proyecto de construcción de unas nuevas instalaciones. El acero estructural está en su lugar y los conductos de calefacción están siendo instalados cuando un gerente senior te informa que está preocupado de que el proyecto no cumpla con los estándares de calidad. ¿Qué debes hacer en esta situación?

\begin{enumerate}[label=\alph*)]
  \item \textbf{Asegurarle a la gerencia senior que durante el proceso Planificar la Gestión de la Calidad se determinó que el proyecto cumpliría con los estándares de calidad.}
  \item Estimar análogamente los resultados futuros.
  \item Formar un equipo de aseguramiento de calidad.
  \item Revisar los resultados del último plan de gestión de la calidad.
\end{enumerate}

\subsection*{Pregunta 429}
Estás creando un proyecto para un diseño industrial. Has completado con la Lista de actividades los diagramas de red. ¿Cuál es la siguiente tarea que debes realizar?

\begin{enumerate}[label=\alph*)]
  \item Crear el cronograma
  \item Estimar las duraciones de cada actividad
  \item Asignar recursos a cada actividad
  \item En función del la línea base del alcance, realizar la secuencia de actividades
\end{enumerate}

\subsection*{Pregunta 430}
Estás dirigiendo un proyecto en un ambiente justo a tiempo. Esto requerirá más atención, porque la cantidad de inventario en tal ambiente generalmente es:

\begin{enumerate}[label=\alph*)]
  \item 45 por ciento.
  \item 10 por ciento.
  \item \textbf{12 por ciento.}
  \item O por ciento.
\end{enumerate}

\subsection*{Pregunta 431}
Estás en la mitad de la ejecución de una modificación grande a un producto existente cuando te enteras de que los recursos prometidos al principio del proyecto no están disponibles. Lo MEJOR que se puede hacer es:

\begin{enumerate}[label=\alph*)]
  \item Mostrar cómo los recursos fueron prometidos originalmente a tu proyecto.
  \item Volver a planificar el proyecto sin los recursos .
  \item Explicar el impacto si los recursos prometidos no están disponibles .
  \item Comprimir el proyecto.
\end{enumerate}

\subsection*{Pregunta 432}
Estás en una situación similar a la del caso anterior y decide solapar actividades. ¿Qué consecuencias puede conllevar su aplicación:

\begin{enumerate}[label=\alph*)]
  \item Aumentar los costes pero sin modificar el alcance y por tanto los riesgos
  \item Reducir costes al reducir el plazo final del proyecto
  \item Poder reducir el alcance del producto al dejar parte de los requerimientos no prioritarios sin hacer
  \item Mayor riesgo de tener que corregir actividades y tener que rehacerlas
\end{enumerate}

\subsection*{Pregunta 433}
Estás evaluando las probabilidades de los riesgos y el impacto en los costes de cada uno de los sucesos posibles para tomar finalmente en cuenta aquel que tenga un valor monetario del riesgo negativo menor. Con qué diagrama lo harás:

\begin{enumerate}[label=\alph*)]
  \item Diagrama de pez
  \item Diagrama de simulación
  \item Árbol de decisión
  \item Diagrama de influencia
\end{enumerate}

\subsection*{Pregunta 434}
Estás gestionando un proyecto de diseño y montaje de una nueva línea de producción con un presupuesto hasta la conclusión (budget at completion, BAC) de 1.000.000 euros, un Valor Ganado (EV) en el momento del análisis de 200.000 euros, un Valor Planificado (PV) de 250.000 y un Coste Real de 150.000 euros. ¿Cuál será el CPI?

\begin{enumerate}[label=\alph*)]
  \item 1,34
  \item 1,25
  \item 0,2
  \item 0,25
\end{enumerate}

\subsection*{Pregunta 435}
Estás gestionando un proyecto en el que ocurre un problema grave e inesperado que va a retrasar el cronograma mucho más allá del plazo de reserva previsto. ¿Qué harías?

\begin{enumerate}[label=\alph*)]
  \item Solicitar rápidamente un aumento del plazo al cliente para garantizar la calidad final
  \item Reducir los controles de calidad solapando fases y tareas
  \item Evaluar las posibles alternativas y su impacto sobre el plazo final
  \item Informar al patrocinador (sponsor) para que se reúna con el cliente y valoren la situación
\end{enumerate}

\subsection*{Pregunta 436}
Estás planificando con tu equipo las fases para la tramitación de un desarrollo urbanístico. Tienes definidas las actividades y las has secuenciado estableciendo precedencias y adelantos o retrasos. ¿qué es lo siguiente que debes hacer?

\begin{enumerate}[label=\alph*)]
  \item Estimar la duración de las actividades
  \item Estimar los recursos de las actividades
  \item Estimar el coste de las actividades
  \item Cerrar el ciclo desarrollando el cronograma
\end{enumerate}

\subsection*{Pregunta 437}
Estás planificando el cronograma y recursos para una subestación eléctrica. Para la parte de obra civil tienes claro el número de trabajadores que necesitarás –un encargado, un oficial de primera y un peón albañil así como unas cantidades precisas de hormigón, lacrillo, estucado,… pero la parte eléctrica, al no tener tanta experiencia, no tienes claro cómo estimar los recursos. ¿Qué harías para determinar los recursos de esta parte del proyecto?

\begin{enumerate}[label=\alph*)]
  \item Preguntar al departamento eléctrico
  \item Realizar una estimación preliminar con tu equipo de trabajo
  \item Cualquiera de las anteriores
\end{enumerate}

\subsection*{Pregunta 438}
Estas preocupado por determinar bien qué riesgos pueden afectar al proyecto. ¿Cuál es el proceso por el cual lograrlo?

\begin{enumerate}[label=\alph*)]
  \item Gestión de riesgos del proyecto
  \item Identificar riesgos
  \item Realizar el análisis cualitativo de riesgos
  \item Priorizar los riesgos relevantes
\end{enumerate}

\subsection*{Pregunta 439}
Estás preparando el presupuesto del proyecto para lo cual estás agregando los costes individuales de cada actividad. para hacerlo necesitas la siguientes información excepto:

\begin{enumerate}[label=\alph*)]
  \item El calendario de recursos
  \item La base de las estimaciones
  \item La línea de desempeño del coste
  \item La línea de base del alcance
\end{enumerate}

\subsection*{Pregunta 440}
Estás preparando la elaboración de un plato nuevo de cocina en el que los tres elementos del plato deberán estar listos, en el punto correcto y a la temperatura adecuada para emplatarlos y sacarlos al comedor. ¿Qué tipo de relación deberán tener las actividades de elaboración de cada elemento entre sí?

\begin{enumerate}[label=\alph*)]
  \item Final a final (Finish-to-Finish, FF)
  \item Inicio a Inicio (Start-to-Start, SS)
  \item Final a inicio (Finish-to-Start, FS)
  \item Inicio a Final (Start to finish, SF)
\end{enumerate}

\subsection*{Pregunta 441}
Estás preparando tu respuesta a los riesgos y de pronto identificas riesgos adiciona les. ¿Qué debes hacer?

\begin{enumerate}[label=\alph*)]
  \item Añadir reservas al proyecto con el fin de enfrentarte a nuevos riesgos, y notificar a la gerencia.
  \item Documentar las acciones de riesgo y calcular el valor monetario esperado en base a la probabilidad e impacto de los incidentes .
  \item Determinar los eventos de riesgo y costos asociados, y posteriormente sumarle este costo al presupuesto del proyecto a manera de reserva.
  \item \textbf{Sumar 10 por ciento de contingenc ia al presupuesto del proyecto y notificar al cliente.}
\end{enumerate}

\subsection*{Pregunta 442}
Estás recibiendo propuestas técnicas y económicas de potenciales proveedores para su evaluación. ¿en qué proceso estás?

\begin{enumerate}[label=\alph*)]
  \item Planificar las compras
  \item Administrar las compras
  \item Efectuar las adquisiciones
  \item Decidir si comprar o no Respuestas razonadas
\end{enumerate}

\subsection*{Pregunta 443}
Estas revisando ofertas y tratando de seleccionar la oferta más adecuada para la adquisición de un componente singular y de alto valor añadido que requiere una experiencia y conocimiento que no hay en la propia compañía. ¿Cuál de las siguientes técnicas o herramientas no deberías utilizar?

\begin{enumerate}[label=\alph*)]
  \item Realizar una subasta
  \item Juicios de expertos
  \item Estimaciones independientes
  \item Utilizar el método Delphi con personas experimentadas y de confianza
\end{enumerate}

\subsection*{Pregunta 444}
Estás teniendo problemas para estimar el costo de un proyecto . ¿Cuál de las siguientes opciones describe MEJOR la causa más probable de tu dificultad?

\begin{enumerate}[label=\alph*)]
  \item \textbf{Definición inadecuada del alcance}
  \item Falta de disponibilidad de los recursos deseados
  \item Falta de registros históricos de proyectos anteriores
  \item Falta de procesos de la compañía
\end{enumerate}

\subsection*{Pregunta 445}
Estas trabajando en una compañía y en un país en los que es práctica habitual recibir regalos cuando contratos importantes se firman. En este caso el valor del regalo supera el valor establecido por tu compañía para estos casos. Unos día antes tu habías entregado un regalo de menor valor en señal de buena voluntad y para como comienzo de una relación que durará los cuatro años que va a durar el contrato. ¿Qué debieras hacer?

\begin{enumerate}[label=\alph*)]
  \item Aceptar el regalo
  \item Contactar con la compañía y solicitar consejo
  \item Rechazar el regalo amablemente explicando la política de tu empresa
  \item Recomendarle al cliente que envíe el regalo la dirección de la compañía.
\end{enumerate}

\subsection*{Pregunta 446}
Estas trabajando en el departamento de marketing de tu compañía gestionando tres trabajos interrelacionados entre sí. Si utilizas la Gestión de Programas para su gestión una de las siguientes afirmaciones no será correcta:

\begin{enumerate}[label=\alph*)]
  \item Los programas son grupos de proyectos gestionados de forma conjunta para conseguir beneficios superiores a los obtenidos si se gestionaran de manera individual.
  \item Un proyecto puede ser o no parte de un programa
  \item Los gestores de programas monitorizan el desempeño de los proyectos para asegurarse de que cumplen plazos y presupuestos y satisfacen las expectativas de las partes interesadas
  \item La gestión de programas analiza las interdependencias de los proyectos y ayuda a determinar la forma de gestión más adecuada de los mismos.
\end{enumerate}

\subsection*{Pregunta 447}
Estás tratando de determinar si las actividades de los proyectos cumplen con las políticas de la organización. Esta actividad forma parte de:

\begin{enumerate}[label=\alph*)]
  \item La inspección
  \item Las auditorías de calidad
  \item La mejora de la calidad
  \item El análisis de tendencias
\end{enumerate}

\subsection*{Pregunta 448}
Este documento identifica el propósito del proyecto y la descripción del producto e incluye un resumen del presupuesto y un programa de hitos, ¿de qué documento estamos hablando?

\begin{enumerate}[label=\alph*)]
  \item El acta de constitución
  \item El plan de dirección del proyecto
  \item El presupuesto del proyecto
  \item El Enunciado del Trabajo
\end{enumerate}

\subsection*{Pregunta 449}
Federico gestiona la instalación e implementación de una aplicación en la empresa. Para gestionar el riesgo de sobrepasar el coste de este nuevo tipo de aplicación, necesita que el subcontratista ejecute aún un contrato de precio fijo cerrado. Utilizando este tipo de contrato, el contratista asume el riesgo en caso de sobrepasar los costes. Este tipo de estrategia de respuesta al riesgo es: 00

\begin{enumerate}[label=\alph*)]
  \item Evitar
  \item Transferir
  \item Mitigar
  \item Aceptar
\end{enumerate}

\subsection*{Pregunta 450}
Una forma de obtener la variación del cronograma (SV) es partir del valor ganado (EV) y...

\begin{enumerate}[label=\alph*)]
  \item \textbf{Restar el valor planificado (PV)}
  \item Restar el costo real (AC)
  \item Multiplicar por el coste real (AC)
  \item Multiplicarlo por el valor planificado (PV)
\end{enumerate}

\subsection*{Pregunta 451}
Generalmente las comunicaciones mejoran cuando el emisor\_\_\_\_\_\_\_\_\_\_\_ receptor

\begin{enumerate}[label=\alph*)]
  \item Utiliza gestos cuando le habla al
  \item \textbf{Muestra preocupación por la perspectiva del}
  \item Le habla lentamente al
  \item Habla más alto al
\end{enumerate}

\subsection*{Pregunta 452}
Gestión de la Calidad o e H o Calidad es:

\begin{enumerate}[label=\alph*)]
  \item Cumplir y superar las expectativas del cliente.
  \item Agregar extras para hacer feliz al cliente.
  \item \textbf{El grado en el que el proyecto cumple con los requisitos.}
  \item Cumplimiento de los objetivos de la dirección.
\end{enumerate}

\subsection*{Pregunta 453}
Gestión de la Calidad o e H o El equipo de proyectos ha creado un plan para la forma en que implementarán la política de calidad. Tiene en cuenta la estructura organizacional, responsabilidades, procedimientos y otra información con respecto a los planes de calidad. Si este plan cambia durante el proyecto, ¿CUÁL de los siguientes planes puede que cambie también?

\begin{enumerate}[label=\alph*)]
  \item Plan de aseguramiento de calidad
  \item Plan de gestión de la calidad
  \item Plan para la dirección del proyecto
  \item Plan de control de calidad
\end{enumerate}

\subsection*{Pregunta 454}
Gestión de la Integración e u A T R o Acaban de asignarte para que te hagas cargo de un proyecto de otro director de proyectos que deja la compañía. El director del proyecto anterior te dice que el proyecto está a tiempo, pero sólo porque él ha presionado constantemente al equipo para que rinda . ¿Qué es lo PRIMERO que debes hacer como el nuevo director del proyecto?

\begin{enumerate}[label=\alph*)]
  \item Revisar el estado del riesgo.
  \item \textbf{Revisar el desempeño de costos.}
  \item Determinar una estrategia de dirección.
  \item Informar tus objetivos al equipo.
\end{enumerate}

\subsection*{Pregunta 455}
Gestión de la Integración e u A T R o áctic..a La necesidad de \_\_ es una de las mayores razones para la comunicación en un proyecto .

\begin{enumerate}[label=\alph*)]
  \item Optimización
  \item Integridad
  \item Integración
  \item Diferenciación
\end{enumerate}

\subsection*{Pregunta 456}
Gestión de la Integración e u A T R o ¿Cuál de las siguientes opciones es VERDADERA acerca del desarrollo del acta de constitución del proyecto?

\begin{enumerate}[label=\alph*)]
  \item El patrocinador crea el acta de constitución del proyecto y el director del proyecto la aprueba.
  \item El equipo del proyecto crea el acta de constitución del proyecto y la PMO la aprueba.
  \item El gerente ejecutivo crea el acta de constitución del proyecto y el gerente funcional la aprueba .
  \item El director del proyecto crea el acta de constitución del proyecto y el patrocinador la aprueba.
\end{enumerate}

\subsection*{Pregunta 457}
Gestión de la Integración e u A T R o ¿Cuál de los siguientes refleja MEJOR la frase "influir en los factores que afectan al cambio"?

\begin{enumerate}[label=\alph*)]
  \item Decirles a las personas que los cambios no están permitidos después de finalizar la planificación
  \item Determinar las fuentes de los cambios y solucionar las causas raíz
  \item Agregar más actividades a la estructura de desglose del trabajo para adaptar los riesgos
  \item \textbf{Calcular el impacto de los cambios hasta la fecha sobre el proyecto 152}
\end{enumerate}

\subsection*{Pregunta 458}
Gestión de la Integración e u A T R o Este proyecto está concebido para determinar nuevas formas de extender la vida de uno de los productos producidos por la compañía . El director del proyecto viene del departamento de ingeniería y el equipo viene de los departamentos de gestión del producto y mercadeo . El enunciado del alcance del proyecto y la planificación del proyecto se han completado, pero un interesado notifica al equipo que hay una mejor manera de completar uno de los paquetes de trabajo. El interesado suministra una carta de revisión técnica de su departamento que prueba que la nueva manera de completar el paquete de trabajo será más rápida que la anterior. proyectos, y estaba esperando que esto sucediera en este proyecto. ¿Qué es lo PRIMERO que el director del proyecto debe hacer?

\begin{enumerate}[label=\alph*)]
  \item Contactar al departamento y reclamar de nuevo sobre la falta de cumplimiento de la fecha límite para la entrega del alcance.
  \item Revisar cómo este cambio impactará en el costo para completar el paquete de trabajo y la calidad del producto del paquete de trabajo.
  \item Ver si hay una manera de cambiar de una organización matricial a una funcional para eliminar la
  \item Preguntar al departamento si tiene algún otro cambio.
\end{enumerate}

\subsection*{Pregunta 459}
Gestión de los Costos s 1 e r e Acabas de completar los procesos de iniciación de un proyecto pequeño y te mueves hacia la planificación del proyecto cuando de pronto una interesada del proyecto te pide el presupuesto del proyecto y la línea base de costos. ¿Qué deberías decirle?

\begin{enumerate}[label=\alph*)]
  \item El presupuesto del proyecto se encuentra en el acta de constitución del proyecto, que acaba de completarse.
  \item El presupuesto del proyecto y su línea base no serán finalizados ni aceptados hasta que los procesos de planificación no sean completados .
  \item El plan para la dirección del proyecto no contendrá el presupuesto del proyecto ni su línea base; se trata de un proyecto pequeño.
  \item \textbf{Es imposible completar un estimado antes de crear el plan para la dirección del proyecto.}
\end{enumerate}

\subsection*{Pregunta 460}
Gestión de los Costos s 1 E r E ¿De qué proceso de gestión de costos es salida la línea base de costos?

\begin{enumerate}[label=\alph*)]
  \item Estimar los Recursos de las Actividades
  \item Estimar los Costos
  \item \textbf{Determinar el Presupuesto}
  \item Controlar los Costos
\end{enumerate}

\subsection*{Pregunta 461}
Gestión de los Costos s 1 e r e l. Una forma común de calcular la estimación a la conclusión (EAC) es tomar el presupuesto hasta la conclusión (BAC) y:

\begin{enumerate}[label=\alph*)]
  \item Dividir por el SPI.
  \item Multiplicar por el SPI.
  \item Multiplicar por el CPI.
  \item Dividir por el CPI.
\end{enumerate}

\subsection*{Pregunta 462}
Gestión de los Interesados T R E e E El punto hasta el cual un interesado en particular puede afectar negativa o positivamente un proyecto es su:

\begin{enumerate}[label=\alph*)]
  \item Nivel de participación .
  \item Nivel de interés.
  \item \textbf{Nivel de compromiso.}
  \item Nivel de influencia.
\end{enumerate}

\subsection*{Pregunta 463}
Gestión de los Recursos Humanos N u E v E Acabas de descubrir que un importante subcontratista de tu proyecto presenta los entregables con retraso en forma sistemática. El subcontratista se comunica contigo y te pide que continúes aceptando los entregables con retraso a cambio de una reducción en los costos del proyecto . Esta oferta es un ejemplo de:

\begin{enumerate}[label=\alph*)]
  \item Confrontar .
  \item Consentir.
  \item Suavizar.
  \item Forzar.
\end{enumerate}

\subsection*{Pregunta 464}
Gestión de los Recursos Humanos N u E v E ¿Qué muestra el histograma de recursos que no muestra una matriz de asignación de responsabilidades?

\begin{enumerate}[label=\alph*)]
  \item Tiempo
  \item Actividades
  \item Interrelaciones
  \item \textbf{La persona encargada de cada actividad}
\end{enumerate}

\subsection*{Pregunta 465}
Gestión de los Recursos Humanos N u e v e Un director de proyectos está tratando de darle fin a una disputa entre dos miembros del equipo. Uno dice que los sistemas deben ser integrados antes de ser probados y el otro mantiene que cada sistema debe ser probado antes de realizar la integración. El proyecto involucra a más de 30 personas y se necesita integrar 12 sistemas. El patrocinador exige que la integración se realice a tiempo. ¿Cuál es la MEJOR declaración que puede hacer el director de proyectos para resolver este conflicto?

\begin{enumerate}[label=\alph*)]
  \item Hazlo a mi manera.
  \item ¡Debemos calmarnos y hacer el trabajo!
  \item Hay que tratar este tema la próxima semana cuando ya estemos todos calmados.
  \item Hagamos una prueba limitada antes de la integración y terminemos la prueba después de la integración.
\end{enumerate}

\subsection*{Pregunta 466}
Gestión de los Riesgos o N e E l. Las siguientes opciones son factores para la evaluación del riesgo del proyecto, EXCEPTO:

\begin{enumerate}[label=\alph*)]
  \item Eventos de riesgo.
  \item Probabilidad del riesgo.
  \item Cantidad en juego.
  \item Primas de seguro.
\end{enumerate}

\subsection*{Pregunta 467}
Gestión de los Riesgos o N e E Las salidas del proceso Planificar la Respuesta a los Riesgos incluyen:

\begin{enumerate}[label=\alph*)]
  \item Riesgos residuales, planes de reserva y reservas para contingencias.
  \item Disparadores de riesgo, contratos y una lista de riesgos.
  \item Riesgos secundarios, actualizaciones al proceso y propietarios del riesgo.
  \item \textbf{Planes de contingencia, actualizaciones al plan para la dirección del proyecto y análisis de sensibilidad.}
\end{enumerate}

\subsection*{Pregunta 468}
Gestión de los Riesgos o N e E Te encontrabas a la mitad de un proyecto de dos años que buscaba implementar nuevas tecnologías en sucursales a lo largo del país. Un huracán ocasionó cortes en el suministro eléctrico justo cuando la actualización estaba a punto de completarse. Cuando se restableció la electricidad, todos los informes del proyecto y datos históricos se habían perdido sin que hubiera posibilidad de recuperarlos. ¿Qué debió hacerse para prevenir este problema?

\begin{enumerate}[label=\alph*)]
  \item Contratar un seguro.
  \item \textbf{Planificar un fondo de reserva.}
  \item Dar seguimiento al clima y tener un plan de contingencia.
  \item Programar la instalación para que no coincidiera con la temporada de huracanes .
\end{enumerate}

\subsection*{Pregunta 469}
Gestión de los Riesgos o N e E Un director de proyectos está diseñando un plan de respuesta a los riesgos. Sin embargo, cada vez que se sugiere una respuesta a los riesgos, se identifica otro riesgo causado por la respuesta. ¿Cuál de las siguientes opciones indica lo MEJOR que puede hacer el director del proyecto?

\begin{enumerate}[label=\alph*)]
  \item \textbf{Documentar los nuevos riesgos y continuar con el proceso Planificar la Respuesta a los Riesgos.}
  \item Asegurarse de que el trabajo del proyecto se entienda mejor.
  \item Dedicar más tiempo para asegurarse de que las respuestas al riesgo estén definidas claramente.
  \item Involucrar a más personas en el proceso Identificar los Riesgos dado que se han pasado de largo algunos riesgos.
\end{enumerate}

\subsection*{Pregunta 470}
Gestión del Alcance e 1 N e o ¿Cuál de las siguientes es una salida del proceso Recopilar los Requisitos?

\begin{enumerate}[label=\alph*)]
  \item La matriz de rastreabilidad de requisitos
  \item Enunciado del alcance del proyecto
  \item La estructura de desglose del trabajo
  \item \textbf{Solicitud de cambios}
\end{enumerate}

\subsection*{Pregunta 471}
Gestión del Alcance e 1 N e o Después de cinco años de trabajar en proyectos, te acabas de integrar a la oficina de la dirección de proyectos . Una de las cosas que quieres introducir en tu compañía es el valor de la creación y la utilización de estructuras de desglose del trabajo . Algunos de los directores de proyectos están enojados porque les estás pidiendo que hagan "trabajo de más' '. ¿Cuál de las siguientes sería la MEJOR cosa que les puedes decir a los directores de proyectos para convencerlos de usar las estructuras de desglose del trabajo?

\begin{enumerate}[label=\alph*)]
  \item Las estructuras de desglose del trabajo evitarán que el trabajo se omita .
  \item Las estructuras de desglose del trabajo sólo son necesarias en proyectos grandes .
  \item Las estructuras de desglose del trabajo sólo son necesarias si el proyecto implica contratos.
  \item \textbf{Las estructuras de desglose del trabajo son la única manera de identificar los riesgos.}
\end{enumerate}

\subsection*{Pregunta 472}
Gestión del Alcance e 1 N e o l. El sistema de numeración de una estructura de desglose del trabajo permite al equipo del proyecto:

\begin{enumerate}[label=\alph*)]
  \item Estimar en forma sistemática los costos de los elementos de la estructura de desglose del trabajo .
  \item Proporcionar la justificación del proyecto.
  \item Identificar el nivel en el cual se encuentran los elementos individuales.
  \item Usarlo en un software de dirección de proyectos.
\end{enumerate}

\subsection*{Pregunta 473}
Gestión del Tiempo s E 1 s ¿Cuándo se utiliza un diagrama de hitos en lugar de un diagrama de barras?

\begin{enumerate}[label=\alph*)]
  \item La planificación del proyecto
  \item Cuando se reporta a los miembros del equipo
  \item Cuando se reporta a la gerencia
  \item Análisis de riesgos
\end{enumerate}

\subsection*{Pregunta 474}
Gestión del Tiempo s e 1 s La gerencia senior se está quejando de que no son capaces de determinar fácilmente el estado de los proyectos en curso dentro de la organización. ¿Cuál de los siguientes tipos de informes ayudaría a proporcionar información resumida a la gerencia senior?

\begin{enumerate}[label=\alph*)]
  \item \textbf{Estimados de costos detallados}
  \item Planes para la dirección del proyecto
  \item Diagramas de barras
  \item Informes de hitos
\end{enumerate}

\subsection*{Pregunta 475}
Gestión del Tiempo s E 1 s Para controlar el cronograma, un director de proyectos está volviendo a analizar el proyecto para predecir su duración . Lo hace mediante el análisis de la secuencia de las actividades con la menor cantidad de flexibilidad de calendarización . ¿Qué técnica está utilizando?

\begin{enumerate}[label=\alph*)]
  \item Método de la ruta crítica
  \item Diagrama de flujo
  \item Diagramación por precedencia
  \item \textbf{La estructura de desglose del trabajo}
\end{enumerate}

\subsection*{Pregunta 476}
Gestión del Tiempo s e 1 s Tienes un proyecto con las siguientes actividades: La actividad A lleva 40 horas y puede comenzar después del inicio del proyecto . La actividad B lleva 25 horas y debe ocurrir después del inicio del proyecto. La actividad C debe ocurrir después de la actividad A y lleva 35 horas . La actividad D debe ocurrir después de las actividades B y C, y lleva 30 horas. La actividad E debe realizarse después de la actividad C y lleva 10 horas. La actividad F se realiza después de la Actividad E y lleva 22 horas. Las actividades F y D son las últimas actividades del proyecto. ¿Cuál de las siguientes es VERDADERA si la actividad B tomara 37 horas?

\begin{enumerate}[label=\alph*)]
  \item La ruta crítica sería de 67 horas.
  \item La ruta crítica cambiaría a Inicio, B, D, Final.
  \item La ruta crítica sería Inicio, A, C, E, F, Final.
  \item La ruta crítica se incrementaría por 12 horas.
\end{enumerate}

\subsection*{Pregunta 477}
Hacerle pruebas a la población entera podría:

\begin{enumerate}[label=\alph*)]
  \item Tomar demasiado tiempo.
  \item \textbf{Proveer más información de la requerida .}
  \item Ser mutuamente excluyente.
  \item Mostrar muchos defectos.
\end{enumerate}

\subsection*{Pregunta 478}
Has descompuesto los entregables del proyecto en un nivel que permite hacer estimaciones de horas y costes razonables. El departamento financiero te ha dado un único identificador para asignar a cada uno de los elementos desglosados.

\begin{enumerate}[label=\alph*)]
  \item Es un código de identificación que se asigna a cada elemento de la EDT
  \item Es un número de cuenta en al que se cargarán todos los costes imputados a cada nivel de la EDT
  \item Es un código cualitativo que permitirá imputar diferentes ratios de costes a las horas asignadas
  \item Es un código de identificación de la EDT que permita diferenciarla de proyectos similares
\end{enumerate}

\subsection*{Pregunta 479}
Has participado en la creación del acta de constitución del proyecto, pero no has logrado que sea aprobada. Tu director y su jefe pidieron que el proyecto iniciara inmediatamente. ¿Cuál de las siguientes es la MEJOR opción a realizar?

\begin{enumerate}[label=\alph*)]
  \item Crear un proceso de control integrado de cambios.
  \item Mostrarle a tu director el impacto de proceder sin una aprobación.
  \item Enfocarte en completar proyectos que tienen actas de constitución firmadas.
  \item Comenzar el trabajo sólo con las actividades de la ruta crítica.
\end{enumerate}

\subsection*{Pregunta 480}
Has sido asignado a un proyecto para la elaboración de un procedimiento y un sistema informático de apoyo para la elaboración y gestión de ofertas de contratos llave en mano a clientes. Tú reportas al Project manager responsable de este proyecto y al responsable de Desarrollo de Negocio quienes comparten la responsabilidad en este proyecto. ¿En qué tipo de estructura organizativa trabajas?

\begin{enumerate}[label=\alph*)]
  \item Organización funcional
  \item Organización matricial débil
  \item Organización enfocada a proyectos
  \item Organización matricial equilibrada
\end{enumerate}

\subsection*{Pregunta 481}
Has solicitado ya los cálculos de las cimentaciones de una turbina eólica, enviado el despiece de armaduras a la empresa suministradora que los está preparando. En ese momento, el cliente decide que, para mejorar la eficiencia energética de la turbina, va a modificar el diseño y éste alterará los avances realizados. ¿Qué es lo mejor que puedes, como Project manager, hacer?

\begin{enumerate}[label=\alph*)]
  \item Pedir al suministrador que deje de trabajar y espere instrucciones
  \item Evaluar e informar al cliente del impacto de esos cambios
  \item Continuar trabajando solamente en aquellas partes no afectadas.
  \item Preguntar al cliente cuando tendrá listos los nuevos datos definitivos para adecuar los trabajos.
\end{enumerate}

\subsection*{Pregunta 482}
Has terminado tu informe mensual pero te das cuenta de que el responsable de compras no ha reportado los avances en el diseño e instalación de un software que forma parte del proyecto. Ello conlleva entregar un informe incompleto e inexacto al Program manager. En este caso, ¿Cuál sería la mejor opción?

\begin{enumerate}[label=\alph*)]
  \item Evaluar el impacto con el responsable de compras sobre esas carencias en el proyecto
  \item Informar al Director del Departamento de Compras sobre la falta de colaboración del responsable de compras
  \item Cerrar el informe tal y como está y enviarlo a las personas que, según el plan de comunicación deben recibirlo
  \item Prepara un informe exacto, veraz y completo aunque se retrase
\end{enumerate}

\subsection*{Pregunta 483}
Hay veces que la comunicación no fluye adecuadamente entre los miembros de un equipo debido a diferentes barreras internas al individuo. ¿Cuál de las siguientes no es una barrera?

\begin{enumerate}[label=\alph*)]
  \item Los prejuicios culturales
  \item El estado de ánimo del que envía o recibe el mensaje
  \item La escucha asertiva
  \item Las experiencias pasadas con las personas
\end{enumerate}

\subsection*{Pregunta 484}
Una heurística se describe MEJOR como:

\begin{enumerate}[label=\alph*)]
  \item Herramienta de control.
  \item Método de calendarización .
  \item Herramienta de planificación.
  \item Regla generalmente aceptada.
\end{enumerate}

\subsection*{Pregunta 485}
Identifica cuál de las siguientes no es una herramienta o técnica del proceso de la realización del control de calidad

\begin{enumerate}[label=\alph*)]
  \item Diagramas de afinidad
  \item Muestreo
  \item Evaluación del desempeño
  \item Diagrama de Pareto
\end{enumerate}

\subsection*{Pregunta 486}
Identifica cuál de los siguientes aspectos no será producto del proceso de desarrollar el equipo de proyecto

\begin{enumerate}[label=\alph*)]
  \item Informe de la eficiencia del equipo de proyecto y de su trabajos
  \item Informe del desarrollo de las habilidades y competencias del equipo
  \item lecciones aprendidas para mejorar en futuros proyectos
  \item La capacitación necesaria para ejecutar sus actividades
\end{enumerate}

\subsection*{Pregunta 487}
Incrementar los recursos en las actividades del camino crítico no siempre acorta la duración del proyecto por: aún 00

\begin{enumerate}[label=\alph*)]
  \item El Director de Proyecto siempre selecciona los recursos mejor preparados
  \item Las actividades no dependen del tiempo ni de los recursos
  \item La duración de las actividades siempre se calcula con el número óptimo de recursos
  \item Incrementar el número de recursos puede generar más trabajo y producir ineficiencias Luis ha identificado un riesgo en donde se utiliza el antiguo hardware para ejecutar un software más actual. Luis ha recibido la aprobación para proseguir actualizando el hardware para asegurar que el aún
\end{enumerate}

\subsection*{Pregunta 488}
Indique cuál de las siguientes afirmaciones es cierta, acerca de una actividad que es predecesora de otra que está en el camino crítico:

\begin{enumerate}[label=\alph*)]
  \item Ninguna de las anteriores
  \item Su margen total es mayor o igual que su margen libre
  \item Su margen total es igual a su margen libre
  \item Su margen libre es igual a cero
\end{enumerate}

\subsection*{Pregunta 489}
Indique qué afirmación sobre la planificación de un proyecto largo y complejo en el que hay dos caminos críticos es verdad

\begin{enumerate}[label=\alph*)]
  \item Los dos caminos críticos podrían no tener ningún hito ni actividad en común
  \item Para gestionar mejor el proyecto se suele aumentar o disminuir la duración de alguna actividad para que sólo haya un camino crítico
  \item En algún caso podría haber un camino crítico de duración menor que el otro
  \item Todas las afirmaciones anteriores son falsas
\end{enumerate}

\subsection*{Pregunta 490}
Informar sobre el desempeño en la fase de cierre administrativo es importante ya que…

\begin{enumerate}[label=\alph*)]
  \item Muestra los problemas que se han producido en cada actividad
  \item Se comunica a los miembros del equipo el resultado del proyecto
  \item Se obtienen la aprobación de la dirección para identificar las lecciones aprendidas
  \item Se puede certificar la aprobación formal por parte del cliente
\end{enumerate}

\subsection*{Pregunta 491}
La acción de dar seguimiento a los costos incurridos hasta la fecha con el fin de detectar variaciones respecto del plan tiene lugar durante:

\begin{enumerate}[label=\alph*)]
  \item La creación del plan de gestión de los cambios en el costo.
  \item Recomendación de acciones correctivas.
  \item \textbf{Actualización de la línea base de costos.}
  \item Revisiones del desempeño del proyecto.
\end{enumerate}

\subsection*{Pregunta 492}
La acción de reorganizar los recursos de modo que se utilice cada mes un número constante de recursos se le llama:

\begin{enumerate}[label=\alph*)]
  \item \textbf{Compresión .}
  \item Holgura.
  \item Nivelación.
  \item Ejecución rápida.
\end{enumerate}

\subsection*{Pregunta 493}
La cantidad presupuestada para el trabajo realmente completado de la actividad del cronograma o el componente de la EDT durante un periodo de tiempo determinado es...

\begin{enumerate}[label=\alph*)]
  \item El valor ganado
  \item El costo real
  \item El valor planificado
  \item La variación del costo
\end{enumerate}

\subsection*{Pregunta 494}
La categoría que se asigna a productos y servicios que tienen el mismo uso funcional pero características diferentes es el

\begin{enumerate}[label=\alph*)]
  \item Las características funcionales
  \item la calidad de un producto
  \item Los requisitos técnicos de un producto
  \item El Grado de un producto
\end{enumerate}

\subsection*{Pregunta 495}
La comunicación durante la fase de proyecto debe, muy especialmente, satisfacer las necesidades de

\begin{enumerate}[label=\alph*)]
  \item El Project manager
  \item El cliente
  \item La alta dirección
  \item Las partes interesadas
\end{enumerate}

\subsection*{Pregunta 496}
La comunicación formal escrita es menos importante para cuál de las siguientes comunicaciones:

\begin{enumerate}[label=\alph*)]
  \item Cambios en el alcance inicial del proyecto
  \item Hacer una excepción sobre un requerimiento técnico del proyecto
  \item Un contrato de adquisición de equipos
  \item Un cambio en la convocatoria de la reunión semanal
\end{enumerate}

\subsection*{Pregunta 497}
La comunicación relativa a los contratos de adquisiciones o compras conviene que utilicen:

\begin{enumerate}[label=\alph*)]
  \item Comunicación formal escrita
  \item Comunicación formal verbal
  \item Comunicación informal escrita
  \item Comunicación informal verbal
\end{enumerate}

\subsection*{Pregunta 498}
La creencia o la imagen mental de un interesado acerca del futuro es un/a:

\begin{enumerate}[label=\alph*)]
  \item Requisito.
  \item Heurística.
  \item Expectativa.
  \item \textbf{Restricción. 543}
\end{enumerate}

\subsection*{Pregunta 499}
La desviación estándar es una medida de qué tan/tanto:

\begin{enumerate}[label=\alph*)]
  \item Lejos está el estimado del estimado más alto.
  \item Lejos está la medición de la media.
  \item Correcta es la muestra .
  \item \textbf{Tiempo queda del proyecto. 324}
\end{enumerate}

\subsection*{Pregunta 500}
La diferencia entre la línea base de costos y el presupuesto de costos puede describirse MEJOR como:

\begin{enumerate}[label=\alph*)]
  \item Las reservas de gestión .
  \item \textbf{Las reservas para contingencias.}
  \item El estimado del costo del proyecto.
  \item La cuenta de costo.
\end{enumerate}

\subsection*{Pregunta 501}
La diferencia entre un proyecto, un programa y un portafolio es que:

\begin{enumerate}[label=\alph*)]
  \item \textbf{Un proyecto es una tarea temporal con un principio y un final, un programa puede incluir otro trabajo no relacionado con el proyecto, y un portafolio incluye todos los proyectos de un determinado departamento o división.}
  \item Un proyecto es una tarea prolongada con un principio y un final, un programa combina dos o más proyectos no relacionados, y un portafolio combina dos o más programas.
  \item Un proyecto es una tarea temporal con un principio y un final, un programa es un conjunto de proyectos relacionados, y un portafolio es un grupo de proyectos y programas relacionados con un objetivo estratégico específico.
  \item Un proyecto es una tarea contratada con un principio y un final, un portafolio es un grupo de proyectos con fechas de conclusión de final más abierto, y un programa combina dos o más portafolios .
\end{enumerate}

\subsection*{Pregunta 502}
La EDT, los estimados para cada paquete de trabajo y el diagrama de red están concluidos. Lo SIGUIENTE que debe hacer el director del proyecto es:

\begin{enumerate}[label=\alph*)]
  \item Secuenciar las actividades.
  \item Validar que tengan el alcance correcto .
  \item Crear un cronograma preliminar y obtener la aprobación del equipo.
  \item Completar la gestión de los riesgos.
\end{enumerate}

\subsection*{Pregunta 503}
La EDT y el diccionario de la EDT han sido completados. El equipo del proyecto ha comenzado a trabajar para identificar riesgos. El patrocinador contacta al director del proyecto y le solicita que la matriz de asignación de responsabilidades sea presentada. El proyecto tiene un presupuesto de US\$100.000 y se llevará a cabo en tres países con la ayuda de 14 recursos humanos . Se pronostica poco riesgo para el proyecto y el director del proyecto ha dirigido muchos proyectos parecidos a este. ¿Qué es lo SIGUIENTE que debemos hacer?

\begin{enumerate}[label=\alph*)]
  \item Entender la experiencia del patrocinador en proyectos similares .
  \item Crear una lista de actividades .
  \item Asegurarse . de tener un alcance del proyecto definido .
  \item Completar la gestión de riesgos y emitir la matriz de asignación de responsabilidades.
\end{enumerate}

\subsection*{Pregunta 504}
La estimación a la conclusión (EAC) es una evaluación periódica de:

\begin{enumerate}[label=\alph*)]
  \item El costo del trabajo completado .
  \item \textbf{El valor del trabajo llevado a cabo.}
  \item El costo total anticipado a la conclusión del proyecto.
  \item Lo que costará terminar el proyecto.
\end{enumerate}

\subsection*{Pregunta 505}
La Estimación ascendente utiliza:

\begin{enumerate}[label=\alph*)]
  \item Datos de proyectos similares
  \item La suma de los costes individuales de cada actividad y paquete de trabajo
  \item Utiliza parámetros estadísticos y simulaciones de Montecarlo
  \item Se utiliza para imputar las horas realmente dedicadas en la fase de proyecto
\end{enumerate}

\subsection*{Pregunta 506}
La estructura de desglose del trabajo puede considerarse MEJOR como una ayuda efectiva para las comunicaciones del \_\_

\begin{enumerate}[label=\alph*)]
  \item Equipo
  \item Director del proyecto
  \item \textbf{Cliente}
  \item Interesado
\end{enumerate}

\subsection*{Pregunta 507}
La evaluación de la probabilidad e impacto de los riesgos, representado o no en una matriz probabilidad-impacto, la evaluación de la calidad de los datos sobre los riesgos, la categorización de los riesgos y la evaluación de su urgencia son Técnicas y Herramientas para el proceso de:

\begin{enumerate}[label=\alph*)]
  \item análisis cuantitativo de riesgos
  \item análisis cualitativo de riesgos
  \item priorización de riesgos
  \item seguimiento y control de riesgos
\end{enumerate}

\subsection*{Pregunta 508}
La fase de construcción de un nuevo producto de software está por finalizar. Las siguientes fases son prueba e implementación. El proyecto lleva dos semanas de adelanto con respecto al cronograma . ¿Por cuál de los siguientes procesos el director de proyectos debe preocuparse MÁS antes de pasar a la fase final?

\begin{enumerate}[label=\alph*)]
  \item Validar el alcance
  \item Realizar el Control de Calidad
  \item Gestionar las comunicaciones
  \item \textbf{Controlar los Costos 188}
\end{enumerate}

\subsection*{Pregunta 509}
La frase que MEJOR describe el rol del patrocinador del proyecto es:

\begin{enumerate}[label=\alph*)]
  \item Ayudar a planificar actividades.
  \item Ayudar a prevenir cambios innecesarios a los objetivos del proyecto.
  \item \textbf{Identificar restricciones innecesarias al proyecto.}
  \item Ayudar a desarrollar el plan para la dirección del proyecto.
\end{enumerate}

\subsection*{Pregunta 510}
La gestión de los recursos humanos del proyecto del proyecto consta de los siguientes procesos:

\begin{enumerate}[label=\alph*)]
  \item Asignación, Formación, Dirección y Evaluación de los recursos humanos
  \item Planificar los recursos humanos, adquirir y Desarrollar el equipo de proyecto y gestionar el equipo de proyecto.
  \item Planificar, adquirir, desarrollar y evaluar el equipo de proyecto
  \item Identificar a los interesados, preparar el plan de comunicación, asignar las actividades y dirigir y gestionar al equipo
\end{enumerate}

\subsection*{Pregunta 511}
La gestión de los recursos humanos del proyecto en la Gestión de proyectos incluye: A las partes interesadas, patrocinadores, clientes, equipo de trabajo, proveedores,… a.

\begin{enumerate}[label=\alph*)]
  \item A los miembros del equipo de proyecto y Project managers
  \item Los miembros del equipo de proyecto
  \item Identificar, documentar y asignar funciones y responsabilidades
\end{enumerate}

\subsection*{Pregunta 512}
La holgura de una actividad se determina a través de:

\begin{enumerate}[label=\alph*)]
  \item \textbf{Realizar el análisis Monte Carlo.}
  \item Determinar el tiempo de espera entre actividades.
  \item Determinar el retraso .
  \item Determinar la cantidad de tiempo que la actividad se puede atrasar sin que se retrase la ruta crítica.
\end{enumerate}

\subsection*{Pregunta 513}
La integración es realizada por:

\begin{enumerate}[label=\alph*)]
  \item \textbf{El director del proyecto.}
  \item El equipo.
  \item El patrocinador.
  \item Los interesados.
\end{enumerate}

\subsection*{Pregunta 514}
La longitud media establecida para una pieza es de 4 mm. con una tolerancia de 0,2 mm. Qué afirmación no es correcta en relación a la aceptación de las piezas?

\begin{enumerate}[label=\alph*)]
  \item Se rechazarán las que tengan una longitud de 4,3 mm.
  \item Se aceptarán las piezas que tenga una longitud menor de 4,2 mm.
  \item Se rechazarán las piezas que tengan 3,6 mm.
  \item Se aceptarán las piezas que tengan 4,1 mm.
\end{enumerate}

\subsection*{Pregunta 515}
¿La medición del desempeño de costos se lleva a cabo MEJOR si se realiza de cuál de las siguientes maneras?

\begin{enumerate}[label=\alph*)]
  \item \textbf{Pedir a cada miembro del equipo un porcentaje completado e informándolo en un informe mensual de progreso realizar proyecciones del desempeño futuro}
  \item Usar la regla de y garantizar que el costeo del ciclo de vida sea menor que el costo del proyecto
  \item Enfocarse en la cantidad que fue gastada el mes pasado y en la que será gastada el mes siguiente
\end{enumerate}

\subsection*{Pregunta 516}
La metodología de gestión moderna de calidad creada y elaborada por el grupo japonés Toyota para lograr un proceso de mejora continua en la empresa se denomina

\begin{enumerate}[label=\alph*)]
  \item kanban
  \item kaizen
  \item PDCA
  \item Seis Sigma
\end{enumerate}

\subsection*{Pregunta 517}
La numeración de la Estructura de desglose de tareas permite

\begin{enumerate}[label=\alph*)]
  \item Imputar las horas de personal dedicadas por actividad y realizar contabilidad analítica
  \item Facilitar la programación de software de seguimiento del avance de tareas
  \item Identificarlas con más facilidad en el Diccionario de la EDT
  \item Conocer el orden en el que van a ejecutar las actividades
\end{enumerate}

\subsection*{Pregunta 518}
¿La programación lineal es un ejemplo de qué tipo de criterios de selección de proyectos?

\begin{enumerate}[label=\alph*)]
  \item Optimización por restricciones
  \item Enfoque comparativo
  \item Medición de beneficios
  \item \textbf{Análisis de impacto}
\end{enumerate}

\subsection*{Pregunta 519}
La reunión de inicio es todo menos:

\begin{enumerate}[label=\alph*)]
  \item una reunión antes de cerrar el proceso de desarrollo del plan para la dirección del proyecto y comenzar con la ejecución del proyecto. En que participan todas las partes implicadas en el proyecto –clientes, Project manager b. y equipo de proyecto,…-
  \item una reunión que se mantiene en la fase de inicio y se definen los objetivos del proyecto para que estén alineados con los de la compañía y las reglas básicas de comportamiento del equipo, entre otras cosas.
  \item Se tratan aspectos especialmente relevantes del proyecto: los riesgos, la comunicación, distribución de información, programas de reuniones,…
\end{enumerate}

\subsection*{Pregunta 520}
La siguiente definición a qué concepto se refiere: ``nivel en el que un conjunto de características inherentes satisface los requisitos?

\begin{enumerate}[label=\alph*)]
  \item Grado
  \item Calidad
  \item Aseguramiento de la calidad
  \item Satisfacción del cliente
\end{enumerate}

\subsection*{Pregunta 521}
La teoría de dirección que establece que las personas pueden dirigir sus propios esfuerzos es:

\begin{enumerate}[label=\alph*)]
  \item Teoría Y.
  \item Teoría de Herzberg.
  \item \textbf{Jerarquía de Maslow.}
  \item Teoría X. equipo expresaron su falta de certeza respecto del trabajo que estaban por completar . En una situación así,
\end{enumerate}

\subsection*{Pregunta 522}
¿Las acciones correctivas aprobadas son una entrada a cuáles de los siguientes procesos?

\begin{enumerate}[label=\alph*)]
  \item Validar el Alcance
  \item Dirigir y Gestionar el Trabajo del Proyecto
  \item Desarrollar el Acta de Constitución del Proyecto
  \item Desarrollar el Cronograma
\end{enumerate}

\subsection*{Pregunta 523}
Las causas MÁS frecuentes de conflicto en un proyecto son los cronogramas, prioridades del proyecto y:

\begin{enumerate}[label=\alph*)]
  \item \textbf{Personalidad.}
  \item Recursos .
  \item Costo.
  \item Gerencia.
\end{enumerate}

\subsection*{Pregunta 524}
Las evaluaciones del desempeño del proyecto difieren de las evaluaciones del desempeño del equipo en que las evaluaciones del desempeño del proyecto se enfocan en:

\begin{enumerate}[label=\alph*)]
  \item El desempeño de un miembro individual del equipo en un proyecto.
  \item Una evaluación de la efectividad del equipo del proyecto.
  \item Un esfuerzo de formación de equipo.
  \item \textbf{Reducir la tasa de rotación de personal.}
\end{enumerate}

\subsection*{Pregunta 525}
Las materias primas son un ejemplo de...

\begin{enumerate}[label=\alph*)]
  \item \textbf{Costo variable}
  \item Costo fijo
  \item Costo directo
  \item Costo indirecto
\end{enumerate}

\subsection*{Pregunta 526}
Las reservas de contingencias: Las reservas de contingencias: opción:

\begin{enumerate}[label=\alph*)]
  \item Se calcula como un porcentaje del coste estimado a a. Se calcula como un porcentaje del coste estimado
  \item Puede utilizarse a medida que avanza el proyecto para cubrir imprevistos no previstos b. Puede utilizarse a medida que avanza el proyecto para cubrir imprevistos no previstos
  \item \textbf{Puede utilizarse a medida que avanza el proyecto para cubrir imprevistos previstos (5) c. Puede utilizarse a medida que avanza el proyecto para cubrir imprevistos previstos}
  \item Se utiliza cuando hay un cambio aprobado en la linea base de costes del proyecto d. Se utiliza cuando hay un cambio aprobado en la línea base de costes del proyecto
\end{enumerate}

\subsection*{Pregunta 527}
Las siguientes son características de un proyecto, EXCEPTO:

\begin{enumerate}[label=\alph*)]
  \item Es temporal.
  \item Tiene un comienzo y final definitivos .
  \item Tiene actividades interrelacionadas.
  \item Se repite cada mes.
\end{enumerate}

\subsection*{Pregunta 528}
Las siguientes son todas partes del esfuerzo de gestión de los interesados del equipo, EXCEPTO:

\begin{enumerate}[label=\alph*)]
  \item Determinar las necesidades de los interesados.
  \item Identificar a los interesados.
  \item Dar extras a los interesados.
  \item \textbf{Gestionar las expectativas de los interesados . 541}
\end{enumerate}

\subsection*{Pregunta 529}
¿Una lista de supervisión para riesgos de baja prioridad es una salida de qué proceso de gestión de riesgos?

\begin{enumerate}[label=\alph*)]
  \item Planificar la Respuesta a los Riesgos
  \item \textbf{Realizar el Análisis Cuantitativo de Riesgos}
  \item Realizar el Análisis Cualitativo de Riesgos
  \item Planificar la Gestión de los Riesgos
\end{enumerate}

\subsection*{Pregunta 530}
Lo MEJOR sería que las lecciones aprendidas fueran completadas por:

\begin{enumerate}[label=\alph*)]
  \item El director del proyecto.
  \item El equipo .
  \item El patrocinador.
  \item Los interesados.
\end{enumerate}

\subsection*{Pregunta 531}
Los costes de corrección, ajuste, parada,… debidos a la detección de errores durante el tiempo en el que el producto o servicio se está elaborando son:

\begin{enumerate}[label=\alph*)]
  \item Costes de prevención
  \item Costes de evaluación
  \item Costes de fallos internos
  \item Costes de fallos externos
\end{enumerate}

\subsection*{Pregunta 532}
Los costos de instalación del proyecto son un ejemplo de:

\begin{enumerate}[label=\alph*)]
  \item Costos variables.
  \item Costos fijos.
  \item \textbf{Costos generales.}
  \item Costos de oportunidad.
\end{enumerate}

\subsection*{Pregunta 533}
Los costes de no conformidad de un proyecto los componen:

\begin{enumerate}[label=\alph*)]
  \item Los costes de fallos internos y externos
  \item Solamente los costes de prevención
  \item Los costes de evaluación y los fallos internos
  \item Solamente los fallos internos
\end{enumerate}

\subsection*{Pregunta 534}
Los diagramas de Pareto ayudan al director de proyectos a:

\begin{enumerate}[label=\alph*)]
  \item Centrarse en las polémicas más críticas para mejorar la calidad.
  \item Centrarse en estimular el pensamiento .
  \item Explorar una futura salida deseada.
  \item \textbf{Determinar si un proceso está fuera de control.}
\end{enumerate}

\subsection*{Pregunta 535}
Los Project manager no siempre pueden crear equipos de trabajo perfectamente adecuados a las características y complejidad del proyecto. Para evitar conflictos a la hora de asignar las funciones y responsabilidades, el Project manager debería:

\begin{enumerate}[label=\alph*)]
  \item Elaborar una matriz de asignación de responsabilidades con el equipo del proyecto
  \item Evitar compartir actividades entre miembros de diferentes departamentos funcionales
  \item Estar muy pendiente del desarrollo de la ejecución de las actividades más conflictivas
  \item Dirigir personalmente las actividades equipo de trabajo
\end{enumerate}

\subsection*{Pregunta 536}
Los estándares –ISO, PMBOK®, estándares internos de la organización,… que sin ser legalmente obligatorios recogen buenas prácticas en el desarrollo de proyecto-, la legislación vigente– en diferentes áreas como seguridad laboral, medio ambiente, protección de datos, legislación laboral,…- y otras reglas o guías metodológicas que puedan ser de interés son:

\begin{enumerate}[label=\alph*)]
  \item Factores ambientales de la empresa a considerar al planificar la calidad
  \item activos de los procesos de la compañía con los que planificar la calidad
  \item son pautas que la compañía debe obligatoriamente cumplir
  \item son documentos que la compañía puede optar por aplicar no
\end{enumerate}

\subsection*{Pregunta 537}
Los incentivos al proveedor en el contrato tienen como objetivo:

\begin{enumerate}[label=\alph*)]
  \item Alinear las objetivos de vendedor y comprador
  \item Asegurar que no se realizan trabajos no incluidos
  \item Reducir costes para el comprador
  \item Mejorar los beneficios del vendedor
\end{enumerate}

\subsection*{Pregunta 538}
Los informes de desempeño son una entrada al proceso de Control de Costes. Incluye todo excepto:

\begin{enumerate}[label=\alph*)]
  \item El cálculo de CV, SV, CPI, SPI para los componentes de la EDT
  \item Estimaciones de las actividades programadas finalizadas
  \item Porcentaje de actividades sin comenzar
  \item Costes autorizados e incurridos
\end{enumerate}

\subsection*{Pregunta 539}
Los patrocinadores del proyecto deben tener la capacidad de tomar decisiones de manera efectiva. ¿Cuál de los siguientes aspectos es importante para facilitarles este trabajo?

\begin{enumerate}[label=\alph*)]
  \item Los patrocinadores deben poseer toda la información técnica detallada a considerar en la toma de una decisión
  \item Los miembros del equipo deben ayudarle al patrocinador a tomar decisiones en todo momento
  \item Los patrocinadores deben consultar a las partes interesadas antes de tomar cualquier decisión
  \item Los patrocinadores y los Project managers deben establecer una relación y unas expectativas claras desde el inicio del proceso de planificación.
\end{enumerate}

\subsection*{Pregunta 540}
Los procedimientos de la compañía requieren de la creación de un documento de lecciones aprendidas . ¿Cuál de los siguientes es el MEJOR uso que se le puede dar a las lecciones aprendidas?

\begin{enumerate}[label=\alph*)]
  \item Registros históricos para proyectos futuros
  \item Registro de planificación para el proyecto actual
  \item Informar al equipo acerca de lo que ha hecho el director del proyecto
  \item \textbf{Informar al equipo sobre el plan para la dirección del proyecto}
\end{enumerate}

\subsection*{Pregunta 541}
Los procesos de gestión de las comunicaciones son:

\begin{enumerate}[label=\alph*)]
  \item Planificar, divulgar y hacer un seguimiento de las comunicaciones
  \item Identificar grupos de interés, planificar estrategia de comunicación y ejecutarla
  \item Planificar, gestionar y monitorizar las comunicaciones
  \item Planificar, asignar recursos y monitorizar las comunicaciones
\end{enumerate}

\subsection*{Pregunta 542}
Los procesos de gestión de los costos y sus herramientas se seleccionan generalmente...

\begin{enumerate}[label=\alph*)]
  \item A medida que se ejecuta el proyecto
  \item \textbf{En el plan de gestión de costos}
  \item Durante la definición del ciclo de vida del proyecto
  \item Depende de cada proyecto
\end{enumerate}

\subsection*{Pregunta 543}
Los requisitos de comunicación con las partes interesadas, las razones para la distribución de la información, frecuencia, métodos para preparar la información y diagramas de flujos de información son elementos de (elige la respuesta más adecuada):

\begin{enumerate}[label=\alph*)]
  \item Plan de gestión de la comunicación
  \item Plan de gestión del proyecto
  \item Análisis de los requerimientos de la comunicación
  \item Plan de distribución de la información
\end{enumerate}

\subsection*{Pregunta 544}
Los requisitos de recursos de la actividad y el calendario de recursos son:

\begin{enumerate}[label=\alph*)]
  \item Entradas para la estimación de la duración de actividades
  \item Salidas de la estimación de los recursos de las actividades
  \item Entradas de la estimación de los recursos de las actividades
  \item Solamente una de ellas es entrada para el desarrollo del cronograma
\end{enumerate}

\subsection*{Pregunta 545}
Los riesgos identificados son:

\begin{enumerate}[label=\alph*)]
  \item Una entrada para el proceso Estimar los Costos.
  \item Una salida del proceso Estimar los Costos.
  \item \textbf{No están relacionados con el proceso Estimar los Costos.}
  \item Tanto una entrada como una salida del proceso Estimar los Costos.
\end{enumerate}

\subsection*{Pregunta 546}
Una manera de estimar los costes a la conclusión (EAC) es:

\begin{enumerate}[label=\alph*)]
  \item \textbf{Dividiendo el Presupuesto a la conclusión (BAC) entre el índice de desempeño del coste}
  \item Dividiendo el Presupuesto a la conclusión (BAC) entre el índice de desempeño del cronograma
  \item Sumando al Presupuesto a la conclusión (BAC) la variación del cronograma
  \item Sumando al Presupuesto a la conclusión (BAC) la variación del coste
\end{enumerate}

\subsection*{Pregunta 547}
Marco de Referencia para la Dirección de Proyectos o o s Dos directores de proyectos se acaban de dar cuenta de que están en una organización matricial débil y que su poder como directores de proyectos está bastante limitado. Uno de ellos cae en cuenta de que realmente trabaja de expedidor del proyecto mientras que el otro se da cuenta de que en realidad es un coordinador del proyecto. ¿En qué se diferencia un expedidor de proyectos de un coordinador de proyectos?

\begin{enumerate}[label=\alph*)]
  \item El expedidor del proyecto no puede tomar decisiones.
  \item El expedidor del proyecto puede tomar más decisiones.
  \item El expedidor del proyecto reporta a un director de mayor nivel.
  \item \textbf{El expedidor del proyecto tiene cierta autoridad .}
\end{enumerate}

\subsection*{Pregunta 548}
Marco de Referencia para la Dirección de Proyectos o o s Un equipo de proyectos está trabajando en la manufactura de un nuevo producto, pero está teniendo problemas para crear el acta de constitución del proyecto. ¿Cuál es la MEJOR descripción del verdadero problema?

\begin{enumerate}[label=\alph*)]
  \item No han identificado los objetivos del proyecto.
  \item Están trabajando en un proceso y no un proyecto.
  \item No se ha establecido una fecha de finalización.
  \item No han identificado el producto del proyecto.
\end{enumerate}

\subsection*{Pregunta 549}
Una matriz de evaluación del compromiso de los interesados puede ser utilizada para identificar:

\begin{enumerate}[label=\alph*)]
  \item Brechas de comunicación.
  \item Más interesados.
  \item Relaciones clave entre los interesados.
  \item \textbf{Niveles de habilidades de los interesados.}
\end{enumerate}

\subsection*{Pregunta 550}
Mónica está dirigiendo un proyecto orientado a la elaboración de un diagnóstico sobre inmigración. Hasta el momento se han elaborado cuatro productos entregables que se van a entregar al cliente: tres informes preliminares y uno definitivo. Mónica ha cliente son preliminares y por tanto no tienen tanta importancia. ¿Está de acuerdo con su modo de actuar?

\begin{enumerate}[label=\alph*)]
  \item Depende del tiempo que tenga Mónica para hacer esa tarea
  \item Sí, porque Mónica se imagina que el cliente no va a mirar los informes preliminares, sino solamente el definitivo
  \item Sí, porque el informe definitivo es más importante que los preliminares
  \item \textbf{No, es necesario revisar todos los entregables para asegurarse de cada uno se complete satisfactoriamente. Esta es una de las actividades básicas del proceso verificar el alcance}
\end{enumerate}

\subsection*{Pregunta 551}
Mostrar un comportamiento profesional, informar de comportamientos no éticos, mostrar consideración por otras realidades socioculturales y tratar de tener una comunicación cuidada para evitar malas interpretaciones forma parte de qué característica del Project manager:

\begin{enumerate}[label=\alph*)]
  \item Responsabilidad
  \item Respeto
  \item Imparcialidad
  \item Honestidad
\end{enumerate}

\subsection*{Pregunta 552}
N u E v E Gestión de los Recursos Humanos l. Todas las siguientes son formas de poder que se derivan del puesto del directo r de proyectos, EXCEPTO:

\begin{enumerate}[label=\alph*)]
  \item Formal.
  \item Recompensa.
  \item Penalidad.
  \item Experto.
\end{enumerate}

\subsection*{Pregunta 553}
N u E v E Gestión de los Recursos Humanos Un miembro del equipo no está rindiendo bien en el proyecto porque carece de experiencia en el trabajo realizar el trabajo. ¿Cuál es la MEJOR solución para el director de proyectos?

\begin{enumerate}[label=\alph*)]
  \item Consultar al gerente funcional para determinar incentivos de finalización del proyecto para el miembro del equipo.
  \item Brindar capacitación al miembro del equipo.
  \item Asignar parte de la reserva del cronograma del proyecto.
\end{enumerate}

\subsection*{Pregunta 554}
N u E v E Gestión de los Recursos Humanos Varios paquetes de trabajo han sido completados con éxito en el proyecto y el patrocinador ha hecho recomendaciones para mejoras. El proyecto avanza a tiempo para cumplir con una fecha límite muy ajustada cuando de pronto la actividad sucesora a una de las actividades de la ruta crítica sufre un retraso miembros del equipo que cuentan con las habilidades para ayudar a la actividad en problemas, si así se requiere. están pidiendo que se les retire del proyecto pues no creen que éste pueda ser exitoso. Cuando el director del proyecto indaga más al respecto, descubre que esos miembros del equipo tienen polémicas que no han sido atendidas. ¿Cuál de las siguientes es la MEJOR acción que se puede tomar para mejorar el proyecto?

\begin{enumerate}[label=\alph*)]
  \item Hacer que los miembros del equipo ayuden de inmediato a la actividad en problemas.
  \item Investigar por qué el cronograma del proyecto es agresivo.
  \item Ver quién puede reemplazar a los tres miembros del equipo.
  \item \textbf{Crear un registro de polémicas . 375}
\end{enumerate}

\subsection*{Pregunta 555}
Normalmente, como lo habitual es trabajar en una organización jerarquizada funcional, los Project manager no tienen autoridad sobre su equipo ya que reportan a sus respectivos jefes funcionales. Por lo tanto necesitan ejercer un poder:

\begin{enumerate}[label=\alph*)]
  \item coercitivo
  \item emanado de sus conocimientos técnicos
  \item que le da que determinar la remuneración del equipo
\end{enumerate}

\subsection*{Pregunta 556}
Oyes un comentario machista de una persona clave e imprescindible de tu equipo hacia una compañera recién llegada. ¿Qué harías en esta situación?

\begin{enumerate}[label=\alph*)]
  \item Hacer como que no lo has oído y esperar a ver si se vuelve a repetir
  \item Informar al jefe jerárquico de esta persona
  \item Documentarlo en el informe de desempeño del próximo mes
  \item Tener una reunión privada con esa persona y recordarle que los comentarios sexistas son inaceptables y que se vuelve a hacerlo tendrá consecuencias negativas sobre él.
\end{enumerate}

\subsection*{Pregunta 557}
PAN Engineering Prot Manufacture Manufacture Si un director de proyecto esta utilizando la estructura de desglose del trabajo Si un director de proyecto está utilizando la estructura de desglose del trabajo como la mostrada en el grafico, entonces la forma de la EDT es mas probable como la mostrada en el gráfico, entonces la forma de la EDT es más probable que sea: que sea: :

\begin{enumerate}[label=\alph*)]
  \item El enfoque de ensamblaje ascendente. a. El enfoque de ensamblaje ascendente.
  \item El enfoque de subproyectos. b. El enfoque de subproyectos.
  \item El enfoque de fases de ciclo de vida. c. El enfoque de fases de ciclo de vida.
  \item \textbf{El enfoque de entregables principales. —d. El enfoque de entregables principales.}
\end{enumerate}

\subsection*{Pregunta 558}
Para conocer la situación real de un proyecto en términos económicos ¿qué valor debe calcular?

\begin{enumerate}[label=\alph*)]
  \item El índice del costo real del trabajo realizado
  \item El índice del desempeño del cronograma (SPI).
  \item El índice de desempeño del coste (CPI)
  \item La estimación a la conclusión (VAC)
\end{enumerate}

\subsection*{Pregunta 559}
Para desarrollar el equipo se definen protocolos, formas de actuar del equipo, situaciones que no se deben dar y que el equipo acepte como válidas.¿a qué técnica nos referimos?

\begin{enumerate}[label=\alph*)]
  \item Procedimientos de trabajo en equipo
  \item Reglas básicas
  \item Normalización del trabajo
  \item Alineamiento con los objetivos del proyecto
\end{enumerate}

\subsection*{Pregunta 560}
Para obtener apoyo para el proyecto en toda la organización ejecutante, es MEJOR si el director del proyecto:

\begin{enumerate}[label=\alph*)]
  \item Se asegura de que haya un plan de gestión de las comunicaciones.
  \item Correlaciona la necesidad del proyecto con el plan estratégico de la organización.
  \item Conecta el proyecto con los objetivos personales del patrocinador.
  \item Se asegura de que el plan de gestión incluye la gestión de los miembros del equipo.
\end{enumerate}

\subsection*{Pregunta 561}
Para poder estimar el número de meses que va a llevar la tramitación de un expediente habitual y con ellas calculas tu mejor estimación. ¿Qué método has utilizado?

\begin{enumerate}[label=\alph*)]
  \item La estimación paramétrica
  \item Simulación por Montecarlo
  \item Juicio de expertos
  \item El método PERT
\end{enumerate}

\subsection*{Pregunta 562}
¿Para qué sirve una matriz de asignación de responsabilidades?

\begin{enumerate}[label=\alph*)]
  \item Para saber el orden de intervención de cada miembro del equipo
  \item Para saber quien es responsable de cada actividad
  \item Para conocer cual debe ser el flujo de información
  \item Tener claro quien pertenece y quien no pertenece al equipo de proyecto
\end{enumerate}

\subsection*{Pregunta 563}
¿Para qué utilizarías en tu proyecto un código de cuenta?

\begin{enumerate}[label=\alph*)]
  \item Para identificar riesgos potenciales
  \item Para identificar cada elemento de la EDT
  \item Para facilitar la estimación de duraciones
  \item Para actualizar ficheros del proyecto
\end{enumerate}

\subsection*{Pregunta 564}
Para supervisar y revisar el avance de tu proyecto, tienes en cuenta todos las siguientes elementos menos uno ¿cuál?

\begin{enumerate}[label=\alph*)]
  \item plan para la dirección del proyecto
  \item informes de desempeño
  \item activos de procesos de la organización
  \item solicitudes de cambios aprobadas
\end{enumerate}

\subsection*{Pregunta 565}
Como parte de dirigir su proyecto Angel está implementando respuestas al riesgo, reevaluando el entorno de riesgo del proyecto y evaluando el proceso de riesgo ¿Qué proceso de gestión de proyectos está realizando?

\begin{enumerate}[label=\alph*)]
  \item Planificar la gestión de riesgos
  \item Controlar riesgos
  \item Realizar análisis cuantitativo del riesgo
  \item Realizar análisis cualitativo del riesgo Su proyecto de dos años tiene un avance de ocho meses. Usted está preparando el informe de estado para la gerencia hasta el término del 8avo. mes. A lo largo del 7o mes, el CPI fue de 1.05 y SPI de 1.07. En el 8o mes, SV(\$) fue -\$5,000 y
\end{enumerate}

\subsection*{Pregunta 566}
Como parte de su plan de proyecto debe desarrollar un método efectivo de comunicación para su equipo de interesados multinacional. Dispone de varias opciones de medios. La acción adecuada a tomar en el desarrollo del Plan de comunicación sería:

\begin{enumerate}[label=\alph*)]
  \item \textbf{Realizar un análisis de requisitos de comunicación con la contribución de los interesados del proyecto}
  \item Utilizar múltiples opciones de medios para asegurarse de que todos reciben la información
  \item Utilizar los medios estándar efectivos en sus proyectos previos
  \item Obtener fondos adicionales del patrocinador del proyecto y desarrollar una infraestructura de comunicaciones especificas para el proyecto
\end{enumerate}

\subsection*{Pregunta 567}
Como parte del cierre del proyecto están teniendo contactos con tu equipo de proyecto para identificar y documentar las lecciones aprendidas. ¿Cuál de estos elementos no incluirás en las lecciones aprendidas?

\begin{enumerate}[label=\alph*)]
  \item Causas de variaciones
  \item Justificación de la acción correctiva elegida
  \item Solamente aspectos positivos que sean de utilidad en futuros proyectos
  \item Razones de las variaciones de coste y del sobrecoste final del proyecto
\end{enumerate}

\subsection*{Pregunta 568}
Como PMP® recibes una notificación de que hay una queja por comportamiento no ético probada contra ti. ¿Como PMP® que no puedes esperar que no ocurra?

\begin{enumerate}[label=\alph*)]
  \item Una reprimenda pública por parte de la organización
  \item Que se suspenda tu certificación PMP®
  \item La pérdida del carnet de miembro del PMI®
  \item Una sanción económica de acuerdo a la gravedad del comportamiento
\end{enumerate}

\subsection*{Pregunta 569}
¿Por qué crees que es mejor planificar el control de calidad que directamente realizar inspecciones conforme se van ejecutando los trabajos?

\begin{enumerate}[label=\alph*)]
  \item Mejora la calidad final y es más barato
  \item Reduce la calidad pero es menos caro
  \item Mejora la calidad aunque es más caro
  \item Mejora la calidad aunque los plazos se alarguen
\end{enumerate}

\subsection*{Pregunta 570}
Procesos de Dirección de Proyectos 1 n e s Un miembro del equipo notifica al director del proyecto que las actividades que comprenden un paquete de trabajo ya no resultan apropiadas . Lo MEJOR para el director del proyecto sería estar en qué parte del proceso de la dirección de proyectos:

\begin{enumerate}[label=\alph*)]
  \item Acción correctiva
  \item \textbf{Control integrado de cambios}
  \item Seguimiento y control
  \item Cierre del proyecto
\end{enumerate}

\subsection*{Pregunta 571}
Procesos de Dirección de Proyectos r n e s l. ¿En qué grupo de procesos de dirección de proyectos se crea el presupuesto detallado del proyecto?

\begin{enumerate}[label=\alph*)]
  \item Iniciación
  \item Antes del proceso de dirección de proyectos
  \item \textbf{Planificación}
  \item Ejecución
\end{enumerate}

\subsection*{Pregunta 572}
Qué aspecto es mas importante \%n el cierre del proyecto: Qué aspecto es más importante Bn el cierre del proyecto: :

\begin{enumerate}[label=\alph*)]
  \item Trabajar con el cliente para determinar los criterios de aceptacion a. Trabajar con el cliente para determinar los criterios de aceptación
  \item \textbf{Confirmar que todos los requisitos del proyecto se hayan cumplido b. Confirmar que todos los requisitos del proyecto se hayan cumplido cuestién}
  \item Recolectar informacion histérica de proyectos anteriores c. Recolectar información histórica de proyectos anteriores
  \item Obtener la aprobacién formal de los planes de gestión d. Obtener la aprobación formal de los planes de gestión
\end{enumerate}

\subsection*{Pregunta 573}
¿Qué datos de entrada de la planificación de los recursos hay que tener en cuenta, más en unos

\begin{enumerate}[label=\alph*)]
  \item Los factores políticos
  \item Las condiciones económicas
  \item Los convenios laborales
  \item Las estructuras organizativas
\end{enumerate}

\subsection*{Pregunta 574}
¿Qué describe mejor el objetivo de una matriz de trazabilidad?

\begin{enumerate}[label=\alph*)]
  \item Que no se pierda algún requerimiento o no se tenga la documentación descriptiva que aclare sus razones y alcance y no se lleve a cabo.
  \item relaciona cada actividad con su paquete de trabajo y éste con su Entregable
  \item es un sistema que permite detectar el entregable que contiene errores para que no quede sin corrección
  \item permite comprender el origen de cada cambio y la modificación que implicó en el plan de dirección del proyecto.
\end{enumerate}

\subsection*{Pregunta 575}
¿Que es el análisis del flujo de valor añadido de un producto?

\begin{enumerate}[label=\alph*)]
  \item El análisis de los flujos de materiales e información que se requieren para poner a disposición del cliente un producto o servicio de calidad
  \item lo que la empresa y sus proveedores hace para fabricar un producto, entregarlo al cliente y cobrarlo.
  \item el análisis de las actividades que no añaden valor al producto, no las va apreciar el cliente y no va a pagar por ellas.
  \item es el análisis del flujo de caja del margen de beneficio obtenido con la venta del producto.
\end{enumerate}

\subsection*{Pregunta 576}
Qué es el coste del ciclo de vida de un proyecto?

\begin{enumerate}[label=\alph*)]
  \item Es el valor a fecha de hoy de todos los costes directos e indirectos del proyecto y puesta en marcha del producto.
  \item Es el coste que nos da una tasa interna de rentabilidad del proyecto igual a cero.
  \item Es una estimación de costes que incluye el mantenimiento, instalación, asistencia
  \item Es el valor residual del producto tras las amortizaciones por el uso que se hace durante toda su vida útil.
\end{enumerate}

\subsection*{Pregunta 577}
¿Qué estrategia se debe seguir con el vecindario de una población importante que tiene muchas preocupaciones sobre los impactos ambientales de un Proyecto?

\begin{enumerate}[label=\alph*)]
  \item \textbf{Tratar de mantenerles en lo posible más bien satisfechos}
  \item Solo se puede gestionarlo de forma muy activa
  \item Tenerles muy controlados y evaluados, en todo momento
  \item Mantenerles muy bien informados, de forma muy frecuente
\end{enumerate}

\subsection*{Pregunta 578}
¿Qué facilita más la transmisión de la comunicación entre el patrocinador y el Project manager?

\begin{enumerate}[label=\alph*)]
  \item Cuando se habla de forma clara y vocalizando bien
  \item Cuando se utiliza un nivel de lenguaje adecuado a ambos
  \item Cuando ambas personas empatizan
  \item Cuando tienen formación multicultural
\end{enumerate}

\subsection*{Pregunta 579}
¿Qué grupo de procesos contiene los procesos de Controlar el alcance y Verificar el Alcance?

\begin{enumerate}[label=\alph*)]
  \item Planificación
  \item Cierre
  \item Supervisión y Control
  \item Ejecución
\end{enumerate}

\subsection*{Pregunta 580}
¿Qué grupo de procesos de dirección de proyectos es el que generalmente consume MÁS tiempo y recursos del proyecto?

\begin{enumerate}[label=\alph*)]
  \item Planificación
  \item Diseño
  \item \textbf{Integración}
  \item Ejecución
\end{enumerate}

\subsection*{Pregunta 581}
¿Qué grupo de procesos se enfoca en completar los requisitos del proyecto?

\begin{enumerate}[label=\alph*)]
  \item \textbf{Iniciación}
  \item Planificación
  \item Ejecución
  \item Cierre
\end{enumerate}

\subsection*{Pregunta 582}
¿Qué grupos de procesos deben estar incluidos en todos los proyectos?

\begin{enumerate}[label=\alph*)]
  \item \textbf{Planificación, ejecución y cierre}
  \item Iniciación, planificación y ejecución
  \item Iniciación, planificación , ejecución, seguimiento y control, y cierre
  \item Planificación, ejecución, y seguimiento y control
\end{enumerate}

\subsection*{Pregunta 583}
¿Qué harías si descubrieras que ninguno de los miembros del equipo de proyecto está revisando las actualizaciones semanales del avance del proyecto?

\begin{enumerate}[label=\alph*)]
  \item Establecer un registro de revisión que debe ser firmado tras revisarlo
  \item Simplificar el plan de gestión de la comunicación manteniendo las necesidades de comunicación de los interesados
  \item Crear plantillas y formatos más sencillos que faciliten el trabajo a los miembros del equipo
  \item Establecer reconocimientos a aquellos que realicen correctamente la revisiones
\end{enumerate}

\subsection*{Pregunta 584}
¿qué herramienta utilizarías para reunir a un grupo de partes interesadas y expertos técnicos para conocer sus expectativas y actitud respecto a un producto propuesto?

\begin{enumerate}[label=\alph*)]
  \item Talleres facilitadores
  \item Grupos focalizados
  \item Técnicas de creatividad en grupo
  \item Técnicas de toma de decisión
\end{enumerate}

\subsection*{Pregunta 585}
Qué proyecto elegirías para su desarrollo:

\begin{enumerate}[label=\alph*)]
  \item el proyecto A con una TIR del 15\%,
  \item el proyecto B con un ratio coste beneficio de 1,4,
  \item el proyecto C con un coste de oportunidad de 800.000 euros
  \item o el Proyecto D que tiene un periodo de payback de 6 meses.
\end{enumerate}

\subsection*{Pregunta 586}
¿Qué técnica de solución de conflictos utiliza un director de proyectos cuando dice "No puedo ocuparme de esto ahora''?

\begin{enumerate}[label=\alph*)]
  \item Solución de problemas
  \item Forzar
  \item Eludir
  \item \textbf{Consentir 369}
\end{enumerate}

\subsection*{Pregunta 587}
¿Qué técnica o herramienta se utiliza en el proceso de Control del Alcance?

\begin{enumerate}[label=\alph*)]
  \item Análisis de variación
  \item Análisis de producto
  \item Inspección
  \item Listas de control
\end{enumerate}

\subsection*{Pregunta 588}
¿Qué teoría propone que los esfuerzos de los empleados tienen como consecuencia un desempeño efectivo y que los empleados serán recompensados por sus logros?

\begin{enumerate}[label=\alph*)]
  \item \textbf{Refuerzo condicional}
  \item Jerarquía de Maslow
  \item La teoría de McGregor
  \item Expectativas
\end{enumerate}

\subsection*{Pregunta 589}
¿Qué tipo de costo constituye la capacitación del equipo?

\begin{enumerate}[label=\alph*)]
  \item Directo
  \item Valor actua l neto
  \item Indirecto
  \item \textbf{Fijo}
\end{enumerate}

\subsection*{Pregunta 590}
¿Qué tres atributos utiliza el Modelo de prominencia de los Interesados del proyecto para ayudar al equipo de dirección del proyecto a clasificar a los interesados?

\begin{enumerate}[label=\alph*)]
  \item Poder, influencia, urgencia.
  \item \textbf{Poder, interés, inquietud.}
  \item Poder, urgencia, legitimidad.
  \item Poder, influencia, legitimidad.
\end{enumerate}

\subsection*{Pregunta 591}
¿Quién de los siguientes no es una parte interesada en el proyecto?:

\begin{enumerate}[label=\alph*)]
  \item El Project manager, como responsable de ejecutar el proyecto
  \item El equipo de proyecto que realizará las actividades del proyecto
  \item El cliente que define sus necesidades y requisitos del proyecto
  \item La competencia que aspira al mismo proyecto
\end{enumerate}

\subsection*{Pregunta 592}
¿Quién es el principal responsable de la gestión de la calidad del proyecto?

\begin{enumerate}[label=\alph*)]
  \item El ingeniero del proyecto
  \item El director del proyecto
  \item El gerente de calidad
  \item El miembro del equipo
\end{enumerate}

\subsection*{Pregunta 593}
¿Quién es responsable de la correcta elaboración de la Estructura de Desglose de Tareas?

\begin{enumerate}[label=\alph*)]
  \item El Project manager
  \item La alta dirección del área de proyectos
  \item El cliente
  \item El Project manager y su equipo de proyecto.
\end{enumerate}

\subsection*{Pregunta 594}
¿Quién tiene el MAYOR poder en una organización orientada a proyectos?

\begin{enumerate}[label=\alph*)]
  \item El director del proyecto
  \item El gerente funcional
  \item El equipo
  \item \textbf{Todos comparten el poder}
\end{enumerate}

\subsection*{Pregunta 595}
¿quién tiene el riesgo en un contrato de Tiempo y Materiales?

\begin{enumerate}[label=\alph*)]
  \item El comprador
  \item El vendedor
  \item El equipo de proyecto
  \item Tanto el comprador como el vendedor a partes iguales
\end{enumerate}

\subsection*{Pregunta 596}
Retraso (lag) significa:

\begin{enumerate}[label=\alph*)]
  \item \textbf{La cantidad de tiempo que una actividad se puede atrasar sin retrasar la fecha de finalización del proyecto.}
  \item La cantidad de tiempo que una actividad se puede atrasar sin retrasar la fecha de inicio temprano de su sucesora .
  \item Tiempo de espera.
  \item El producto de un recorrido hacia adelante y hacia atrás. 240
\end{enumerate}

\subsection*{Pregunta 597}
Una salida del proceso Cerrar el Proyecto o Fase es la creación de:

\begin{enumerate}[label=\alph*)]
  \item Los archivos del proyecto .
  \item Un acta de constitución del proyecto .
  \item Un plan para la dirección del proyecto.
  \item \textbf{Un plan de gestión de los riesgos.}
\end{enumerate}

\subsection*{Pregunta 598}
¿Cómo se denomina el documento que como Project manager has preparado para que los solicitantes tengan claro el alcance de los trabajos a realizar dentro del contrato?:

\begin{enumerate}[label=\alph*)]
  \item Enunciado del alcance del trabajo
  \item Contrato
  \item Documentos de compra
  \item Enunciado del alcance del contrato o adquisición
\end{enumerate}

\subsection*{Pregunta 599}
Se está llevando a cabo un proyecto importante cuando de pronto uno de los miembros del equipo revisa el informe del estado del proyecto. Se encuentra con que el proyecto va con retraso. Conforme lee más del informe, descubre que el retraso causará que una de sus actividades sea programada en un momento en el que estará fuera del país y no podrá trabajar en la actividad. Es de gran importancia para el miembro del equipo porque está muy comprometido con el éxito del proyecto y no quiere ser el culpable de que el proyecto se retrase más. ¿Qué es lo MEJOR que podría hacer?

\begin{enumerate}[label=\alph*)]
  \item Contactar inmediatamente al director de proyectos para entregarle su cronograma .
  \item Incluir la información en su próximo informe.
  \item Solicitar que la polémica sea incorporada al registro de polémicas del proyecto .
  \item \textbf{Recomendar acciones preventivas. 374}
\end{enumerate}

\subsection*{Pregunta 600}
Se estimó que un proyecto tomaría seis meses. Después de tres meses, el análisis del valor ganado muestra la siguiente información. ¿Cuáles son las variaciones del cronograma y del costo?

\begin{enumerate}[label=\alph*)]
  \item \textbf{SV (DP) = +\$100,000 and CV (DC) = -\$50,000}
  \item SV (DP) = +\$150,000 and CV (DC) = \$100,000
  \item SV (DP) = -\$50,000 and CV (DC) = +\$150,000
  \item SV (DP) = - \$100,000 y CV (DC) = + \$150,000
\end{enumerate}

\subsection*{Pregunta 601}
Se le ha alertado sobre una oportunidad en su proyecto. Una respuesta al riesgo para esta oportunidad es:

\begin{enumerate}[label=\alph*)]
  \item Considerar las oportunidades cuando lo permita el tiempo.
  \item Enfocarse únicamente en el lado negativo de los riesgos.
  \item \textbf{Explotar el riesgo.}
  \item Transferir el riesgo.
\end{enumerate}

\subsection*{Pregunta 602}
¿Cómo se llama el nivel más bajo de la Estructura de Desglose de Tareas?

\begin{enumerate}[label=\alph*)]
  \item Código de identificación
  \item Paquetes de trabajo
  \item Entregables
  \item Listas de actividad
\end{enumerate}

\subsection*{Pregunta 603}
¿cómo se llama la técnica mediante la cual recomiendas una acción específica para identificar el problema, descubres las causas que hay detrás de él y desarrollas acciones preventivas?

\begin{enumerate}[label=\alph*)]
  \item Análisis de tendencias
  \item Análisis de procesos
  \item Análisis de riesgos
  \item Análisis de causas raíz
\end{enumerate}

\subsection*{Pregunta 604}
Se va a lanzar una nueva línea cosmética que comprende 4 diferentes productos, cada uno de los cuales saldrá cada seis meses para mantener la marca con imagen de innovación continua. El primer y segundo producto se elaborarán con mayor detalle ya mientras que los dos últimos, se dejarán para más adelante. Como Project manager tienes que documentar la matriz de trazabilidad. De las siguientes cuatro afirmaciones solamente una es falsa:

\begin{enumerate}[label=\alph*)]
  \item origen y la razón del mismo y el responsable
  \item Vincula los requisitos con su origen y los sigue a lo largo de la vida del proyecto tratando de que los requisitos añadan valor al proyecto de la empresa y, si no es así, eliminándolos
  \item Evita que en sucesivos cambios que se produzcan se pierda algún requerimiento o no se tenga la documentación descriptiva que aclare sus razones y alcance
  \item Es una técnica utilizada en el proceso de recopilación de requisitos del producto
\end{enumerate}

\subsection*{Pregunta 605}
Según Deming, la mayoría de los problemas de calidad y los errores en un proyecto se podrían resolver por

\begin{enumerate}[label=\alph*)]
  \item El equipo de trabajo
  \item El Project manager y su plan de gestión de calidad
  \item La Dirección de la empresa
  \item Los jefes funcionales técnicos y de compras bien coordinados
\end{enumerate}

\subsection*{Pregunta 606}
¿sería lo MEJOR crear para la segunda mitad del proyecto?

\begin{enumerate}[label=\alph*)]
  \item Una estructura de desglose del trabajo
  \item Un histograma de recursos
  \item Un plan de gestión del personal
  \item Una matriz de asignación de responsabilidades
\end{enumerate}

\subsection*{Pregunta 607}
Si analizas el impacto de un riesgo sobre un objetivo de tu proyecto manteniendo fijas el resto de las variables te ayuda a:

\begin{enumerate}[label=\alph*)]
  \item Tener un análisis cuantitativo del riesgo
  \item Identificar qué riesgos tendrán un mayor impacto sobre el proyecto
  \item Evaluar la probabilidad de que el evento ocurra
  \item Determinar qué implica elegir una alternativa u otra
\end{enumerate}

\subsection*{Pregunta 608}
Si asumimos que los extremos de un rango de estimados son+/- 3 sigmas respecto de la media, ¿cuál de los siguientes estimados de rango tiene el MENOR riesgo?

\begin{enumerate}[label=\alph*)]
  \item 30 días, más o menos 5 días
  \item 22 a 30 días
  \item Optimista = 26 días, más probable= 30 días, pesimista = 33 días
  \item Media de 28 días
\end{enumerate}

\subsection*{Pregunta 609}
Si dado un proyecto que tiene un camino crítico con una duración determinada nos pidieran reducir su duración, eso implicaría necesariamente:

\begin{enumerate}[label=\alph*)]
  \item Incrementar el coste del proyecto
  \item Re-planificar el proyecto de nuevo
  \item Incrementar el personal asignado al proyecto
  \item Ninguna de las anteriores
\end{enumerate}

\subsection*{Pregunta 610}
Si debido al mal desempeño de un proveedor decides cambiarlo y el nuevo proveedor elegido llega mañana al lugar de trabajo, ¿qué sería lo menos efectivo que podrías hacer?

\begin{enumerate}[label=\alph*)]
  \item Tomar las riendas de la nueva adquisición
  \item Enviarle al director el plan de dirección del proyecto
  \item Tener previsto un tercer proveedor por si este segundo no cumple
  \item Llevarles a que conozcan el lugar de trabajo en el que se van a desenvolver.
\end{enumerate}

\subsection*{Pregunta 611}
Si como Project manager tienes que subcontratar un servicio de suministro de material electrónico cuyas especificaciones y alcance están bien definidos y el presupuesto se puede estimar con suficiente precisión. ¿Qué tipo de contrato sería más adecuado?

\begin{enumerate}[label=\alph*)]
  \item Orden de compra
  \item contrato de Tiempo y materiales
  \item contrato de Precio fijo
  \item contrato de costes más honorarios fijos
\end{enumerate}

\subsection*{Pregunta 612}
Si como Project manager ves que un miembro del equipo no está realizando correctamente su trabajo, ¿qué es lo mejor que puedes hacer inicialmente?

\begin{enumerate}[label=\alph*)]
  \item Comunicación formal escrita
  \item Comunicación formal verbal
  \item Comunicación informal escrita
  \item Comunicación informal verbal
\end{enumerate}

\subsection*{Pregunta 613}
Si el cliente decide proponer un cambio que puede afectar, entre otras cosas, al presupuesto, como Project manager deberás seguir el procedimiento integrado de gestión de cambios para lo cual, en caso de que el cambio se apruebe deberás actualizar:

\begin{enumerate}[label=\alph*)]
  \item El plan de dirección del proyecto
  \item La línea de base del coste del proyecto
  \item Las estimaciones de costes de las actividades
  \item Los informes de desempeño
\end{enumerate}

\subsection*{Pregunta 614}
Si el Project manager recibe un largo listado de errores en el proyecto, con el número de veces que ocurre y el coste de reparación de cada error y quiere priorizar las actuaciones de mejora del procedimiento de calidad ¿qué herramienta o técnica utilizará?

\begin{enumerate}[label=\alph*)]
  \item Diagramas de causa efecto
  \item Diagrama de Pareto
  \item Gráfica de dispersión
  \item Correlaciones lineales entre frecuencia y coste.
\end{enumerate}

\subsection*{Pregunta 615}
Si el PV es mayor que el EV significa que el proyecto está

\begin{enumerate}[label=\alph*)]
  \item retrasado
  \item adelantado
  \item con ahorros
  \item con sobrecostes
\end{enumerate}

\subsection*{Pregunta 616}
Si el tiempo y costos del proyecto no son tan importantes como el número de recursos utilizados cada mes, ¿cuál de las siguientes es la MEJOR cosa por hacer?

\begin{enumerate}[label=\alph*)]
  \item Realizar el análisis Monte Carlo.
  \item Implementar la técnica de ejecución rápida.
  \item Realizar la optimización de los recursos.
  \item \textbf{Analizar los costos del ciclo de vida. 241}
\end{enumerate}

\subsection*{Pregunta 617}
Si el valor ganado (EV) = 350, el costo real (AC) = 400 y el valor planificado (PV) = 325, ¿cuál es la variación del costo (CV)?

\begin{enumerate}[label=\alph*)]
  \item 350
  \item -75
  \item 400
  \item -50
\end{enumerate}

\subsection*{Pregunta 618}
Si en el gráfico anterior duplicamos los recursos para hacer la actividad E detrayéndolos de la actividad H, nos encontramos con que la Duración de la actividad E pasará a ser de 5 y la de H de 12. ¿Cuál sería entonces el camino crítico?

\begin{enumerate}[label=\alph*)]
  \item I-B-C-D-L
  \item I-E-C-D-L
  \item I-E-F-L
  \item I-G-H-L
\end{enumerate}

\subsection*{Pregunta 619}
Si en tu proyecto, tienes una actividad de la cual no estás seguro que puedas terminarla en plazo y correctamente, te supone un riesgo a la hora de cumplir con los objetivos del proyecto, y decides subcontratarla a una empresa que tenga el conocimiento, ¿qué estrategia estás adoptando?

\begin{enumerate}[label=\alph*)]
  \item Aceptar el riesgo
  \item Mitigar el riesgo
  \item Transferir el riesgo
  \item Evitar el riesgo
\end{enumerate}

\subsection*{Pregunta 620}
Si eres el Project manager de una compañía que va a suministrar servicios a otra compañía. ¿Cuál de los siguientes tipos de contrato tendrá más riesgo económico para tu compañía?

\begin{enumerate}[label=\alph*)]
  \item El contrato de precio fijo, CPF (Fixed Price, FP)
  \item contratos de Tiempo y Materiales, CTM (Time \& Materials)
  \item contrato de Costes Reembolsables, CCR, (Cost reimbursable, CR)
  \item Coste reembolsable más un variable por desempeño (Costs plus award fee, CPAF)
\end{enumerate}

\subsection*{Pregunta 621}
Si estás como Project manager de un proyecto internacional, ¿qué es lo mejor que podrías hacer para prevenir problemas culturales entre miembros del equipo?

\begin{enumerate}[label=\alph*)]
  \item Formar al equipo en los procesos de gestión del proyecto a utilizar
  \item Exigir que todos los miembros del equipo hablen un mismo idioma
  \item Formar al equipo en las diferencias culturales de los diferentes países implicados
  \item Formar al equipo en la historia y costumbres de los países implicados
\end{enumerate}

\subsection*{Pregunta 622}
Si estás teniendo problemas con tu equipo en la fase de inicio, habiéndole enviado toda la información de lo que deben hacer pero sin éxito, cuál será la mejor opción:

\begin{enumerate}[label=\alph*)]
  \item porqué las cosas no están funcionando correctamente.
  \item Reportar al patrocinador para que tome las medidas disciplinarias pertinentes
  \item Al estar en la fase de inicio, conviene optar por una actitud de colaboración y resolver en equipo los problemas de comunicación o falta de claridad en los objetivos y actividades.
  \item Solicitar un cambio de equipo más acorde con las necesidades del proyecto y el ritmo necesario para poder llegar a tiempo.
\end{enumerate}

\subsection*{Pregunta 623}
Si no estás contento con el desempeño de un miembro de tu equipo, a) ¿qué tipo de comunicación utilizarías para comentar tu opinión con el empleado? B) Si a pesar del primer contacto, el empleado sigue actuando de una manera que consideras incorrecta, qué forma de comunicación utilizarías para informar a su superior.

\begin{enumerate}[label=\alph*)]
  \item A) informal verbal, B) informal escrita
  \item A) informal verbal, B) formal escrita
  \item A) formal escrita, B) formal escrita
  \item A) informal escrita, B) formal escrita
\end{enumerate}

\subsection*{Pregunta 624}
Si nos dan las siguientes estimaciones del presupuesto de una actividad del proyecto que estás ejecutando, ¿cuál de ellas tendrá menos riesgo de equivocación?:

\begin{enumerate}[label=\alph*)]
  \item 500.000 euros más/menos 100.000 euros
  \item Entre 450.000 y 550.000 euros
  \item Valor optimista: 350.000, valor pesimista: 650, valor probable: 500.000
  \item Valor medio 520.000
\end{enumerate}

\subsection*{Pregunta 625}
Si quieres analizar los datos de un proceso para ver si el proceso está bajo control, qué herramienta utilizarías:

\begin{enumerate}[label=\alph*)]
  \item diagramas de comportamiento,
  \item diagramas de dispersión,
  \item diagramas de flujo,
  \item diagramas de control
\end{enumerate}

\subsection*{Pregunta 626}
Si seguimos con el caso anterior, en el que sabemos que el valor esperado es 8 meses y la desviación estándar 1 mes. Puedes decir con una probabilidad del 95\% que la duración estará en el intervalo

\begin{enumerate}[label=\alph*)]
  \item Entre 5 y 11 días
  \item Entre 7 y 9 días
  \item Entre 6 y 10 días
  \item Entre 4 y 12 días
\end{enumerate}

\subsection*{Pregunta 627}
Si te das un tiempo extra de un dos días más para la realización de la actividad de extensión de aglomerado para tener en cuenta que puede haber problemas meteorológicos extraordinarios por tratarse de las fechas que se tratan, estás

\begin{enumerate}[label=\alph*)]
  \item Aceptando el riesgo
  \item Haciendo una reserva de tiempo o contingencia
  \item Mitigando el riesgo
  \item Siendo muy conservador y quitándole validez al cronograma.
\end{enumerate}

\subsection*{Pregunta 628}
Si trabajas en una organización jerarquizada y necesitas crear un equipo de trabajo en una época en la que la contratación externa está parada por problemas económicos, qué es lo mejor que puedes hacer?

\begin{enumerate}[label=\alph*)]
  \item Identificar las personas adecuadas y solicitar al patrocinador que se incluyan en el acta de constitución
  \item Negociar con los jefes de departamentos funcionales para ver cómo dotar al equipo del personal más cualificado
  \item Consultar los activos de los procesos de la compañía para ver cómo proceder.
  \item Negarte a hacer el proyecto ya que no tienes los medios y no puedes garantizar su éxito.
\end{enumerate}

\subsection*{Pregunta 629}
Si un director de proyecto está utilizando datos históricos sobre 3 proyectos similares para realizar la estimación de duración de las actividades qué herramienta está utilizando:

\begin{enumerate}[label=\alph*)]
  \item Estimación ascendente
  \item Estimación análoga
  \item Estimación basada en tres valores
  \item Estimación paramétrica
\end{enumerate}

\subsection*{Pregunta 630}
Si un Project manager necesita estimar los resultados a final de proyecto, necesitará hacer:

\begin{enumerate}[label=\alph*)]
  \item Un análisis de tendencias
  \item Una predicción
  \item Un informe de variaciones
  \item Un cálculo por analogía
\end{enumerate}

\subsection*{Pregunta 631}
Si un director de proyectos está preocupado por recolectar, integrar y difundir las salidas de todos los procesos de dirección de proyectos, debe concentrarse en mejorar:

\begin{enumerate}[label=\alph*)]
  \item Estructura de desglose del trabajo (EDT).
  \item Plan de gestión de las comunicaciones.
  \item Sistema de información para la dirección de proyectos (PMIS).
  \item El plan de gestión del alcance.
\end{enumerate}

\subsection*{Pregunta 632}
Si un equipo externo al equipo de proyecto, ya interno ya externo a la compañía, proceden a comprobar que el equipo de proyecto cumple con la política de calidad de la compañía, los estándares y procedimientos y que estos son efectivos y permiten asegurar el cumplimiento de los requisitos del cliente y dar sugerencias para su mejora, estamos ante una:

\begin{enumerate}[label=\alph*)]
  \item Inspecciones de trabajo
  \item Control de procesos
  \item Auditorias de calidad
  \item Procedimientos de mejora continua
\end{enumerate}

\subsection*{Pregunta 633}
Si un evento de riesgo tiene 90 por ciento de posibilidades de ocurrir y su consecuencia seria de US\$10.000, Si un evento de riesgo tiene 90 por ciento de posibilidades de ocurrir y su consecuencia sería de US\$10.000, Adn sin équé representa la cifra de US\$9,000? ¿qué representa la cifra de US\$9.000? opción:

\begin{enumerate}[label=\alph*)]
  \item Valor actual a. Valor actual
  \item Presupuesto de contingencias Ob. Presupuesto de contingencias
  \item Valor del riesgo Cc. Valor del riesgo
  \item \textbf{Valor monetario esperado 3d. Valor monetario esperado}
\end{enumerate}

\subsection*{Pregunta 634}
Si un miembro de tu equipo que trabaja en tu oficina está siendo poco eficiente, se está desentendiendo de su trabajo y llega tarde a las reuniones, ¿qué es lo mejor que puede hacer el Project manager para tratar de solucionarlo?

\begin{enumerate}[label=\alph*)]
  \item Llamarle por teléfono y decirle que esto no puede continuar así
  \item Convocarle a una reunión y abordar las causas y buscar soluciones
  \item Tener una reunión a 3 bandas: él, el Gerente funcional y el Project manager.
  \item Comentarlo en la próxima reunión de seguimiento del trabajo con el equipo de proyecto
\end{enumerate}

\subsection*{Pregunta 635}
Si un proyecto tiene una probabilidad del 60 por ciento de obtener ganancias de US\$100.000, y una probabilidad del 40 por ciento de que sus pérdidas sean de US\$100.000, el valor monetario esperado (EMV) del proyecto es:

\begin{enumerate}[label=\alph*)]
  \item Ganancia de \$100.000
  \item Pérdida de \$60.000
  \item Ganancia de \$20.000
  \item Pérdida de \$40.000
\end{enumerate}

\subsection*{Pregunta 636}
Si un riesgo tiene una probabilidad de un 20 por ciento de ocurrir en cierto mes, y se espera que el proyecto dure cinco meses, ¿cuál es la probabilidad de que el evento de riesgo ocurra en el cuarto mes del proyecto?

\begin{enumerate}[label=\alph*)]
  \item Menos de 1 por ciento
  \item \textbf{20 por ciento}
  \item 60 por ciento
  \item 80 por ciento
\end{enumerate}

\subsection*{Pregunta 637}
Si utilizas un diagrama de Ishikawa para identificar de forma estructurada las causas de un problema encontrado tratando de identificar los posibles orígenes de los problemas y sus causas raíz, estás

\begin{enumerate}[label=\alph*)]
  \item Asegurando la calidad
  \item Planificando la calidad
  \item Inspeccionando la calidad
  \item Realizando el control y el aseguramiento de la calidad
\end{enumerate}

\subsection*{Pregunta 638}
SV(\%) fue -2\%. CV(\$) fue -\$12,000 y CV(\%) fue -7\%. ¿Qué deberá hacer a continuación del director del proyecto para preparar el informe ante los interesados?

\begin{enumerate}[label=\alph*)]
  \item Recalcular el CPI y el SPI utilizando todos los datos hasta el 8o mes.
  \item Calcular el CPI y SPI utilizando sólo los datos del 8o mes.
  \item Reportar que el CPI es ahora 0,98 y el SPI de 1.05.
  \item Reportar que las cosas iban bien hasta el 8o mes y que los datos del 8o mes pueden representar una anomalía. Dejar el CPI y el SPI como estaban hasta el 7o mes y esperar a tener los datos correspondientes al 9o mes para determinar si algo malo está sucediendo.
\end{enumerate}

\subsection*{Pregunta 639}
T R E e E Gestión de los Interesados l. ¿Cuándo tiene el interesado una MAYOR influencia en un proyecto?

\begin{enumerate}[label=\alph*)]
  \item Al inicio del proyecto
  \item En la mitad del proyecto
  \item Al final del proyecto
  \item \textbf{A lo largo de todo el proyecto}
\end{enumerate}

\subsection*{Pregunta 640}
T R E e E Gestión de los Interesados La información del registro de los interesados debe ser:

\begin{enumerate}[label=\alph*)]
  \item Accesible solamente para el director del proyecto.
  \item Disponible para el director del proyecto y el personal de la PMO.
  \item \textbf{Disponible para todos los interesados y miembros del equipo.}
  \item Compartida con terceros a discreción del director del proyecto.
\end{enumerate}

\subsection*{Pregunta 641}
Te está costando trabajo evaluar el costo exacto del impacto de los riesgos. Es necesario que evalúes utilizando una:

\begin{enumerate}[label=\alph*)]
  \item Base cuantitativa.
  \item \textbf{Base numérica.}
  \item Base cualitativa.
  \item Base econométrica. 445
\end{enumerate}

\subsection*{Pregunta 642}
Te estás haciendo cargo de un proyecto durante la planificación del proyecto y descubres que seis personas han firmado el acta de constitución del proyecto. ¿Cuál de las siguientes opciones te preocupa MÁS?

\begin{enumerate}[label=\alph*)]
  \item Quién será miembro del comité de control de cambios
  \item Gastar más tiempo en la gestión de la configuración
  \item Obtener a un solo patrocinador del proyecto
  \item Determinar la estructura de informe
\end{enumerate}

\subsection*{Pregunta 643}
Te han hecho director de un proyecto nuevo, grande y complejo. Dado que el proyecto es de crítica importancia para el negocio y muy visible, la gerencia senior te ha pedido que analices los riesgos del proyecto y que prepares estrategias de respuesta para ellos lo antes posible. La organización tiene ciertos procedimientos de gestión de los riesgos, pero rara vez lo usan o siguen, además de que históricamente tiene malas referencias en cuanto al manejo de los riesgos. El primer hito del proyecto es dentro de dos semanas . Al momento de preparar el plan de respuesta a los riesgos, ¿de cuál de las siguientes opciones resulta MENOS importante obtener entradas?

\begin{enumerate}[label=\alph*)]
  \item Miembros del equipo del proyecto
  \item Patrocinador del proyecto
  \item Individuos responsables de las políticas de gestión de los riesgos y las plantillas de riesgos
  \item \textbf{Interesados clave 447}
\end{enumerate}

\subsection*{Pregunta 644}
Te han nombrado por primera vez Project manager dentro de una compañía con una estructura organizativa matricial. ¿Cómo esperas que sean las comunicaciones dentro del equipo, partes interesadas, organización,…?

\begin{enumerate}[label=\alph*)]
  \item Sencillas
  \item Transparentes y precisas
  \item Complejas
  \item Difíciles de procedimentar en un plan de gestión de la comunicación
\end{enumerate}

\subsection*{Pregunta 645}
Te piden que prepares un presupuesto para completar un proyecto que había comenzado el año pasado y después fue suspendido por seis meses. Todo lo siguiente sería incluido en el presupuesto del proyecto, EXCEPTO:

\begin{enumerate}[label=\alph*)]
  \item Costos fijos.
  \item \textbf{Costos irrecuperables .}
  \item Costos directos.
  \item Costos variables.
\end{enumerate}

\subsection*{Pregunta 646}
Te piden que selecciones las herramientas y técnicas para implementar un programa de aseguramiento de calidad para complementar las actividades de control de calidad existentes. ¿Cuál de los siguientes NO sería apropiado para este propósito?

\begin{enumerate}[label=\alph*)]
  \item Auditorías de calidad
  \item Muestreo estadístico
  \item Diagramas de Pareto
  \item Grupos focales
\end{enumerate}

\subsection*{Pregunta 647}
Todas las opciones siguientes forman parte del proceso Dirigir y Gestionar el Trabajo del Proyecto, EXCEPTO:

\begin{enumerate}[label=\alph*)]
  \item Identificar los cambios.
  \item Utilizar una estructura de desglose de trabajo.
  \item Implementar acciones correctivas.
  \item \textbf{Implementar un sistema de control del proyecto. 146}
\end{enumerate}

\subsection*{Pregunta 648}
Todas las siguientes acciones deben realizarse durante la iniciación del proyecto, EXCEPTO :

\begin{enumerate}[label=\alph*)]
  \item \textbf{Identificar y documentar las necesidades del negocio.}
  \item Crear un enunciado del alcance del proyecto .
  \item Dividir un proyecto grande en fases.
  \item Acumular y evaluar información histórica.
\end{enumerate}

\subsection*{Pregunta 649}
Todas las siguientes afirmaciones están relacionadas con la Responsabilidad excepto:

\begin{enumerate}[label=\alph*)]
  \item Implica asumir la autoría y las consecuencias de sus acciones,
  \item Tratar de ser proactivo para hacer las cosas de forma correcta
  \item Hacer todo lo necesario para resolver una situación que haya surgido por una falta inicial de responsabilidad.
  \item Informar de comportamientos no éticos de alguna persona del equipo o partes interesadas
\end{enumerate}

\subsection*{Pregunta 650}
Todas las siguientes afirmaciones muestran que el Project manager es una persona respetuosa menos una

\begin{enumerate}[label=\alph*)]
  \item Le cuesta controlar sus emociones y reacciones cuando pueden ser ofensivas para otra persona,
  \item Da por supuesto que el equipo conoce las políticas de empresa relativas a conflictos de intereses, diversidad cultural, igualdad de género o atención al cliente
  \item Ayuda y forma al equipo para que muestre también estas actitudes de respeto en su comportamiento profesional
  \item Muestra siempre obediencia a las órdenes de sus superiores aunque el no esté de acuerdo
\end{enumerate}

\subsection*{Pregunta 651}
Todas las siguientes opciones ocurren durante el grupo de proceso de planificación , EXCEPTO :

\begin{enumerate}[label=\alph*)]
  \item Desarrollar el acta de constitución del proyecto.
  \item Crear la EDT.
  \item Estimar los Costos.
  \item \textbf{Secuenciar las Actividades .}
\end{enumerate}

\subsection*{Pregunta 652}
Todas las siguientes SIEMPRE son entradas del proceso de gestión de los riesgos, EXCEPTO :

\begin{enumerate}[label=\alph*)]
  \item Información histórica .
  \item Lecciones aprendidas.
  \item \textbf{Estructura de desglose del trabajo .}
  \item Informes del estado del proyecto .
\end{enumerate}

\subsection*{Pregunta 653}
Todas las siguientes son formas de poder que derivan del puesto del director del proyecto EXCEPTO:

\begin{enumerate}[label=\alph*)]
  \item Penalidad
  \item \textbf{Recompensa}
  \item Formal
  \item Experto
\end{enumerate}

\subsection*{Pregunta 654}
Todas las siguientes son herramientas de Control de Calidad excepto:

\begin{enumerate}[label=\alph*)]
  \item Inspección
  \item \textbf{Costes de la calidad}
  \item Diagrama de Pareto
  \item Diagrama de espina de pescado
\end{enumerate}

\subsection*{Pregunta 655}
Todas las siguientes son responsabilidades del director de proyectos, EXCEPTO:

\begin{enumerate}[label=\alph*)]
  \item Seleccionar cuidadosamente a los interesados adecuados para el proyecto.
  \item Hacer que los interesados aprueben que se han finalizado todos los requisitos .
  \item Determinar cuándo y en qué medida van a participar los interesados en el proyecto.
  \item \textbf{Pedir a los interesados que te informen acerca de los problemas que surjan en las comunicaciones y las relaciones del proyecto.}
\end{enumerate}

\subsection*{Pregunta 656}
Todas las siguientes son salidas del proceso Estimar los Costos EXCEPTO :

\begin{enumerate}[label=\alph*)]
  \item \textbf{Un entendimiento del riesgo de costo en el trabajo que se ha estimado .}
  \item La prevención de que se incluyan cambios inadecuados en la línea base de costos .
  \item Una indicación del rango de costos posibles para el proyecto .
  \item Documentación de cualquier supuesto que se hizo durante el proceso Estimar los Costos . 280
\end{enumerate}

\subsection*{Pregunta 657}
Todas son partes del grupo de procesos de inicio del proyecto excepto

\begin{enumerate}[label=\alph*)]
  \item Definir los Criterios de aceptación del producto
  \item Definir los límites del proyecto
  \item Definir los objetivos del proyecto
  \item Definir líneas de base para la medición del desempeño técnico
\end{enumerate}

\subsection*{Pregunta 658}
Todo lo que sigue son partes de la línea base del alcance EXCEPTO:

\begin{enumerate}[label=\alph*)]
  \item Plan de gestión del alcance .
  \item Enunciado del alcance del proyecto .
  \item \textbf{Estructura de desglose del trabajo.}
  \item Diccionario de la EDT.
\end{enumerate}

\subsection*{Pregunta 659}
Todo lo siguiente sucede durante el proceso Cerrar el Proyecto o Fase, EXCEPTO:

\begin{enumerate}[label=\alph*)]
  \item Crear lecciones aprendidas.
  \item Aceptación formal.
  \item \textbf{Informes de desempeño.}
  \item Realizar el análisis costo-beneficio.
\end{enumerate}

\subsection*{Pregunta 660}
Todos los elementos que se enumeran a continuación forman parte de la gestión de las expectativas de los interesados menos una:

\begin{enumerate}[label=\alph*)]
  \item Remunerarles para que participen
  \item Identificar las partes interesadas
  \item Determinar las necesidades de las partes interesadas
  \item Gestionar sus expectativas
\end{enumerate}

\subsection*{Pregunta 661}
Todos los siguientes son ejemplo de costes de incumplimiento de calidad EXCEPTO:

\begin{enumerate}[label=\alph*)]
  \item Capacitación de calidad
  \item Costes de garantía
  \item Reproceso
  \item Desperdicio
\end{enumerate}

\subsection*{Pregunta 662}
Todos los siguientes son ejemplos de los costos de incumplimiento EXCEPTO:

\begin{enumerate}[label=\alph*)]
  \item Reproceso.
  \item Capacitación de calidad.
  \item \textbf{Desperdicio .}
  \item Costos de garantía.
\end{enumerate}

\subsection*{Pregunta 663}
Todos los siguientes son herramientas de Realizar el Control de Calidad EXCEPTO:

\begin{enumerate}[label=\alph*)]
  \item \textbf{Inspección.}
  \item Costo de la calidad.
  \item Diagrama de Pareto.
  \item Diagrama de espina de pescado.
\end{enumerate}

\subsection*{Pregunta 664}
Todos los siguientes son partes de un plan de gestión de cambios efectivo, EXCEPTO:

\begin{enumerate}[label=\alph*)]
  \item Procedimientos.
  \item Estándares de los informes.
  \item Reuniones.
  \item Lecciones aprendidas .
\end{enumerate}

\subsection*{Pregunta 665}
Todos los siguientes son pilares en los que se basa la integridad personal del Project manager excepto:

\begin{enumerate}[label=\alph*)]
  \item Responsabilidad
  \item Respeto
  \item Imparcialidad
  \item Eficiencia
\end{enumerate}

\subsection*{Pregunta 666}
Todos los siguientes son resultados normales de la gestión de los riesgos, EXCEPTO:

\begin{enumerate}[label=\alph*)]
  \item Establecimiento de los términos y condiciones del contrato.
  \item Se cambia el plan para la dirección del proyecto.
  \item \textbf{Se cambia el plan de gestión de las comunicaciones.}
  \item Se cambia el acta de constitución del proyecto.
\end{enumerate}

\subsection*{Pregunta 667}
Trabajas con un equipo de proyecto descentralizado. Unos están en las oficinas de Chicago, notas explicativas para comunicarse entre los miembros del equipo. Estos son ejemplos de comunicación:

\begin{enumerate}[label=\alph*)]
  \item Formal escrita
  \item Informal escrita
  \item Formal verbal
  \item Informal verbal
\end{enumerate}

\subsection*{Pregunta 668}
Trabajas en el desarrollo de un parque eólico en otro país por lo que las partes interesadas están dispersas en diferentes países y oficinas. Para evitar problemas, en la reunión de inicio, el Project manager ha definido y consensuado unas reglas básicas de funcionamiento. ¿Cuál de las siguientes medidas crees tú que tendrá mayor impacto en el proyecto?

\begin{enumerate}[label=\alph*)]
  \item Elegir la tecnología y modo de comunicación más adecuado
  \item Establecer un orden del día, de preguntas y respuestas que garantice unas reuniones y discusiones ordenadas
  \item Un buen plan de comunicación
  \item Una buena identificación de requisitos de comunicación
\end{enumerate}

\subsection*{Pregunta 669}
Trabajas como PMP® autónomo y tienes costumbre de tomarte un vaso de vino durante las comidas. Actualmente estás en un proyecto en una compañía en que está mal visto beber y máxime durante horas de trabajo. Quedas con el cliente para comer y pides tu vaso de vino. ¿Cómo definirías el problema?

\begin{enumerate}[label=\alph*)]
  \item Eres un trabajador autónomo por lo que no te afecta esa cultura
  \item Estás violando el código ético relativo a respetar la realidad sociocultural en la que trabajas
  \item Estás poniendo en peligro el proyecto por lo que no debieras beber
  \item No está demostrando un comportamiento profesional
\end{enumerate}

\subsection*{Pregunta 670}
Tras reunirte con las partes interesadas, tienes ya preparada la documentación con los requerimientos del proyecto que podrás utilizar para

\begin{enumerate}[label=\alph*)]
  \item El lista de interesados y el Enunciado del Trabajo
  \item El acta de constitución y la EDT
  \item El mapa de interesados y el acta de constitución
  \item El Enunciado del Alcance del Trabajo y la EDT
\end{enumerate}

\subsection*{Pregunta 671}
Tu compañía acaba de ganar un proyecto nuevo importante . Comenzará en tres meses y está valuado en US\$2.000.000. Tú eres el director de proyectos de un proyecto existente. ¿Qué es lo PRIMERO que debes hacer una vez que te enteras del proyecto nuevo?

\begin{enumerate}[label=\alph*)]
  \item Preguntar a la gerencia cómo utilizará los recursos el proyecto nuevo.
  \item Nivelar los recursos de tu proyecto .
  \item Comprimir tu proyecto .
  \item \textbf{Preguntar a la gerencia cómo afectará a tu proyecto el nuevo proyecto. 147}
\end{enumerate}

\subsection*{Pregunta 672}
Tu compañía quiere contratar un grupo de delineantes para un trabajo cuya duración no está clara. La compañía pagará 50 euros por hora trabajada. El coste total será calculado a partir del precio unitario de la hora multiplicado por el número de horas imputadas. ¿De qué tipo de contrato estamos hablando?

\begin{enumerate}[label=\alph*)]
  \item contrato de Tiempo y Materiales (T \& M)
  \item Coste Reembolsable (RC)
  \item Precio Fijo (FP)
  \item Costes reembolsables con incentivo por desempeño
\end{enumerate}

\subsection*{Pregunta 673}
Tu compañía tiene como política que cada proyecto que se elabora, entre otras cosas, pase por un comité técnico de revisión previo al paso a la fase de construcción del proyecto. En este momento estas preparando la secuencia de actividades. ¿Cómo denominarías esta revisión?

\begin{enumerate}[label=\alph*)]
  \item Punto de Decisión de aprobación formal
  \item Hito
  \item Nuevo alcance del producto
  \item Entregable
\end{enumerate}

\subsection*{Pregunta 674}
Tu equipo ha presentado 434 riesgos e identificado 16 causas importantes de esos riesgos. El proyecto es el último en una serie de proyectos en los que el equipo ha trabajado en conjunto. El patrocinador brinda mucho apoyo, y se invirtió mucho tiempo en asegurar que el trabajo del proyecto estuviera terminado y que los interesados clave firmaran su autorización . Durante la planificación del proyecto, el equipo no puede encontrar una manera efectiva de mitigar o asegurar un riesgo. No es trabajo que pueda ser subcontratado, y tampoco puede eliminarse. ¿Cuál sería la MEJOR solución?

\begin{enumerate}[label=\alph*)]
  \item \textbf{Aceptar el riesgo.}
  \item Continuar investigando maneras de mitigar el riesgo.
  \item Buscar maneras de evitar el riesgo.
  \item Buscar maneras de transferir el riesgo.
\end{enumerate}

\subsection*{Pregunta 675}
Tú eres un director de proyectos para un proyecto de sistemas de información importante. Alguien del departamento de calidad viene a verte para hablar acerca de comenzar una auditoría de calidad de tu proyecto. El equipo, ya bajo presión para completar el proyecto lo más pronto posible, se opone a la auditoría. Debes explicarle al equipo cuál es el propósito de la auditoría de calidad:

\begin{enumerate}[label=\alph*)]
  \item \textbf{Satisfacer parte de una investigación ISO 9000.}
  \item Revisar si el cliente está siguiendo el proceso de calidad.
  \item Identificar políticas ineficientes e ineficaces.
  \item Para controlar la precisión de los costos entregados por el equipo.
\end{enumerate}

\subsection*{Pregunta 676}
Tu gerencia ha decidido que todos los pedidos sean tratados como "proyectos" y que los directores de proyectos actualicen los pedidos todos los días, para resolver polémicas, y para asegurar que el cliente acepte formalmente el producto dentro de los 30 días posteriores a la fecha en que fue completado . Las ganancias provenientes de los pedidos individuales pu.:dt::n variar de los US \$100 hasta los US \$150.000. No se le solicitará al director del proyecto que lleve a cabo planificación o que haga entrega de documentación que vaya más allá del estado diario. ¿Cómo definirías esta situación?

\begin{enumerate}[label=\alph*)]
  \item Dado que cada orden individual es una "tarea temporal': cada orden es un proyecto.
  \item Se trata de dirección de programas dado que hay muchos proyectos involucrados.
  \item Se trata de un proceso recurrente.
  \item Los pedidos que superen \$100.000 en ganancias serían considerados proyectos e involucrarían dirección de proyectos.
\end{enumerate}

\subsection*{Pregunta 677}
Tu plan para la dirección del proyecto resulta en un cronograma del proyecto que es demasiado largo. Si el diagrama de red del proyecto no puede cambiar pero tienes recursos de personal adicionales, ¿qué es lo MEJOR que puedes hacer?

\begin{enumerate}[label=\alph*)]
  \item Implementar la técnica de ejecución rápida.
  \item \textbf{Nivelar los recursos .}
  \item Implementar la técnica de compresión .
  \item Realizar el análisis Monte Cario.
\end{enumerate}

\subsection*{Pregunta 678}
Tú provees un estimado de costos del proyecto al patrocinador del proyecto . Él no está contento con el estimado porque piensa que el precio debería ser menor. Te pide que recortes en 15 por ciento el estimado del proyecto. ¿Qué deberías hacer?

\begin{enumerate}[label=\alph*)]
  \item Comenzar el proyecto y buscar ahorros en el costo de forma constante.
  \item Decirles a todos los miembros del equipo que reduzcan sus estimados en un 15 por ciento.
  \item Informar al patrocinador acerca de las actividades que se recortarán.
  \item Agregar recursos adicionales con tarifas por hora bajas.
\end{enumerate}

\subsection*{Pregunta 679}
Tu proyecto de aglomerado para el tramo de carretera que tienes contratado tiene un presupuesto total de 1.000 euros/metro lineal de calzada de 7 metros de ancho y 10 kilómetros de longitud. El plazo previsto es de 10 semanas, suponiendo un rendimiento constante todas las semanas. En este momento estás en la quinta semana y tu equipo ha aglomerado 4 kilómetros. Cuál es el SPI del proyecto?

\begin{enumerate}[label=\alph*)]
  \item 0,6
  \item 0,5
  \item 0,4
  \item 0,8
\end{enumerate}

\subsection*{Pregunta 680}
Tu proyecto de puesta en marcha de una nueva campaña de genéricos está en plena fase de ejecución. Como buen Project manager están monitorizando y controlando los riesgos para minimizar el impacto negativo y maximizar el impacto positivo de las oportunidades. ¿Cuál es el principal objetivo de las auditorias de riesgos?

\begin{enumerate}[label=\alph*)]
  \item Examinar y documentar la efectividad de las respuestas a los riesgos
  \item Ayudar en la reevaluación de los riesgos ante un cambio significativo
  \item Dar información válida para elaborar el plan de gestión de riesgos en la fase de planificación
  \item Informar de las tendencias en la ejecución del proyecto y el cumplimiento de sus objetivos Respuestas razonadas
\end{enumerate}

\subsection*{Pregunta 681}
Tu proyecto está en la fase inicial, has creado ya el acta de constitución y estás trabajando en la identificación de las partes interesadas. ¿Cuál de las siguientes actividades no tienes que abordar en esta fase?

\begin{enumerate}[label=\alph*)]
  \item Revisar las lecciones aprendidas en proyectos anteriores
  \item Realizar un análisis de partes interesadas
  \item Revisar los documentos de adquisiciones
  \item Crear el sistema de control de cambios
\end{enumerate}

\subsection*{Pregunta 682}
Tu proyecto tiene una cantidad media de riesgo y no está muy bien definido . El patrocinador te entrega un acta de constitución y te pide que confirmes que el proyecto puede completarse dentro del presupuesto del costo del proyecto. ¿Cuál es el MEJOR método para lidiar con esto?

\begin{enumerate}[label=\alph*)]
  \item Constru ir un estimado, dándole forma de un rango de posibles resultados.
  \item Pedirles a los miembros del equipo que te ayuden a estimar los costos en base al acta de constit ución del proyecto.
  \item En base a la información que tienes, calcular una estimac ión paramétrica.
  \item \textbf{Entregar una estimación análoga basada en la historia pasada.}
\end{enumerate}

\subsection*{Pregunta 683}
Tus proyecciones de costo muestran que tendrás sobrecostos al final del proyecto . ¿Cuál de las siguientes acciones debes llevar a cabo?

\begin{enumerate}[label=\alph*)]
  \item Eliminar los riesgos en los estimados y volver a estimar.
  \item Reunirte con el patrocinador para ver qué trabajo se puede hacer antes.
  \item Recortar la calidad .
  \item \textbf{Disminuir el alcance .}
\end{enumerate}

\subsection*{Pregunta 684}
Un acta de constitución del proyecto es un documento que debe ser preparado y autorizado por \_\_\_\_\_\_\_\_\_\_\_\_\_\_ quién le otorga al director de proyecto la aún autoridad para aplicar recursos de la organización. 00

\begin{enumerate}[label=\alph*)]
  \item La alta gerencia / patrocinador del proyecto.
  \item El director del proyecto.
  \item El cliente.
  \item Los gerentes funcionales.
\end{enumerate}

\subsection*{Pregunta 685}
Un diagrama de control ayuda al director de proyectos a:

\begin{enumerate}[label=\alph*)]
  \item Centrarse en las polémicas más críticas para mejorar la calidad.
  \item Centrarse en estimular el pensamiento.
  \item Explorar una futura salida deseada.
  \item Determinar si un proceso está funcionando dentro los límites.
\end{enumerate}

\subsection*{Pregunta 686}
Un diagrama de control muestra siete puntos de datos alineados de un lado de la media. ¿Qué debe hacerse?

\begin{enumerate}[label=\alph*)]
  \item Realizar un diseño de experimentos .
  \item Ajustar el diagrama para reflejar la nueva media.
  \item Encontrar una causa asignable .
  \item Nada. Esta es la regla de siete y puede ser ignorada .
\end{enumerate}

\subsection*{Pregunta 687}
Un directo r de proyectos recibe una llamada de parte de un miembro del equipo notificándole que existe una variación entre la velocidad de un sistema en el proyecto y la velocidad deseada o planificada. El director del proyecto se sorprende ya que esa medición del desempeño tendría que haberse identificado durante la planificación. Si el director del proyecto evalúa si la variación amerita una respuesta, ¿en qué parte de los grupos de procesos de la dirección de proyectos se encuentr a?

\begin{enumerate}[label=\alph*)]
  \item Iniciac ión
  \item Ejecución
  \item Seguimiento y control
  \item \textbf{Cierre 101}
\end{enumerate}

\subsection*{Pregunta 688}
Un Project manager acaba de elaborar el Enunciado del Trabajo. ¿Qué es lo siguiente que debe hacer?

\begin{enumerate}[label=\alph*)]
  \item Crear el Documento de Requisitos del producto
  \item Definir el alcance del producto.
  \item Identificar las actividades concretas del proyecto
  \item Crear la Estructura de Desglose de Tareas
\end{enumerate}

\subsection*{Pregunta 689}
Un Project manager entra a trabajar en una nueva compañía. Al poco tiempo uno de sus miembros del equipo se queja y se niega a hacer el trabajo que le ha sido asignado. Para resolver el conflicto, convoca una reunión con el gerente funcional de esta persona que es de tu mismo nivel. Finalmente esa persona acaba realizando el trabajo que su gerente funcional y el Project manager han acordado. Qué tipo de estructura organizativa tiene la compañía:

\begin{enumerate}[label=\alph*)]
  \item Matriz débil
  \item Matriz fuerte
  \item Funcional
  \item Matriz equilibrada
\end{enumerate}

\subsection*{Pregunta 690}
Un Project manager que ha trabajado durante año y medio en un complejo proyecto de ejecución de un puente atirantado de 500 metros de luz, actualiza su curriculum para presentarlo ante una nueva compañía indicando que es un experto en proyectos de puentes atirantados. ¿Esta siendo honesto?

\begin{enumerate}[label=\alph*)]
  \item No ya que año y medio no es tiempo suficiente para considerarse un experto en este tipo de obras.
  \item No tenemos información suficiente sobre su experiencia.
  \item Sí, trabajar año y medio como Project manager en una obra de esas características se puede considerar suficiente
  \item Si él considera que es capaz de abordar proyectos de este tipo, no está siendo deshonesto.
\end{enumerate}

\subsection*{Pregunta 691}
Un director de proyecto quiere analizar las varias causas y sub-causas que están creando un problema. ¿Qué tipo de diagrama le será de mayor beneficio?

\begin{enumerate}[label=\alph*)]
  \item Diagrama de sistema.
  \item Diagrama costo beneficio.
  \item \textbf{Diagrama de Pareto.}
  \item Diagrama de espina de pescado. Hay varias actividades de ejecución que están en marcha en tu proyecto. Comienzas a preocuparte sobre la previsión de los reportes
\end{enumerate}

\subsection*{Pregunta 692}
Un director de proyecto s ha organizado al equipo del proyecto. Ha identificado 56 riesgos en el proyecto y ha determinado los disparadores de esos riesgos. También los calificó usando una matriz de calificación de riesgos, hizo una prueba de los supuestos y evaluó la calidad de los datos utilizados . El equipo sigue avanzando a través del proceso de la gestión de los riesgos. ¿Qué se le ha olvidado hacer al director del proyecto?

\begin{enumerate}[label=\alph*)]
  \item Simulación
  \item Mitigar el riesgo
  \item \textbf{Calificación general de riesgos para el proyecto 446}
\end{enumerate}

\subsection*{Pregunta 693}
Un Project manager sabe que los cambios en un proyecto son inevitables pero debe transmitir a todas las partes interesadas que:

\begin{enumerate}[label=\alph*)]
  \item El coste de los cambios es mayor en la fase inicial
  \item El coste de los cambios es menor en la fase inicial
  \item El valor añadido de un cambio en la fase inicial es menor
  \item El coste del cambio es menor en la fase de ejecución
\end{enumerate}

\subsection*{Pregunta 694}
Un Project manager tiene que gestionar las expectativas de varias partes interesadas. el patrocinador del proyecto y el cliente son dos partes interesadas muy importantes en cualquier proyecto. en este contexto, patrocinador y cliente hacen todo menos:

\begin{enumerate}[label=\alph*)]
  \item Tener en cuenta la tolerancia al riesgo
  \item Definir reuniones, hitos y entregables clave
  \item Aceptar formalmente el proyecto
  \item Financiar el proyecto
\end{enumerate}

\subsection*{Pregunta 695}
Un Project manager va a comenzar la ejecución de un nuevo edificio de oficinas singular con técnicas utilizadas anteriormente por la empresa en otro proyecto. Al comentarlo con el Project manager del anterior proyecto le informa de que, efectivamente, el coste y el plazo fue muy superior porque tuvieron problemas con la estabilización del terreno previa a la ejecución de las cimentaciones. ¿Dónde debe buscar la información relativa a ese hecho ocurrido y las soluciones adoptadas en aquel proyecto?

\begin{enumerate}[label=\alph*)]
  \item En el sistema de gestión de registros del proyecto
  \item En los activos de procesos de organización de la compañía
  \item En la información de desempeño del trabajo realizada en el proyecto
  \item En los informes de desempeño del proyecto
\end{enumerate}

\subsection*{Pregunta 696}
Un director de proyectos analizó la calidad de los datos de riesgos y le pidió a varios interesados que determinen la probabilidad y el impacto de ciertos riesgos. Está a punto de pasar al siguiente proceso de la gestión de los riesgos. En base a esta información , ¿qué se le ha olvidado hacer al director del proyecto?

\begin{enumerate}[label=\alph*)]
  \item \textbf{Evaluar tendencias en el análisis de riesgos.}
  \item Identificar disparadores.
  \item Proporcionar una matriz estandarizada para calificación del riesgo.
  \item Crear un plan de reserva.
\end{enumerate}

\subsection*{Pregunta 697}
Un director de proyectos debe publicar el cronograma del proyecto. Ya se identificaron las actividades, los tiempos de comienzo/finalización y los recurso s. ¿Qué debe hacer A CONTINUACIÓN el director del proyecto'

\begin{enumerate}[label=\alph*)]
  \item Distribuir el cronograma del proyecto de acuerdo con el plan de gestión de las comunicaciones.
  \item \textbf{Confirmar la disponibilidad de los recursos.}
  \item Refinar el plan para la dirección del proyecto con el fin de que refleje información más precisa sobre los costos.
  \item Publicar un diagrama de barras que ilustre la línea del tiempo .
\end{enumerate}

\subsection*{Pregunta 698}
Un director de proyectos en un proyecto de implementación de un sitio Web multinacional está en una fiesta y habla con unos amigos que utilizarán este nuevo sitio Web con regularidad cuando se termine el proyecto y se implemente el sitio. Ellos describen algunos aspectos molestos del sitio Web actual. El director del proyecto le lleva esta retroalimentación al patrocinador, y propone cambios en el diseño y el alcance. ¿Cuál de los siguientes describe MEJOR lo que ha hecho el director del proyecto?

\begin{enumerate}[label=\alph*)]
  \item Validación del alcance
  \item Control integrado de cambios
  \item Análisis de los interesados
  \item \textbf{Planificación del alcance 544}
\end{enumerate}

\subsection*{Pregunta 699}
Un director de proyectos está cuantificando el riesgo para su proyecto. Varios de sus expertos se encuentran fuera del lugar pero quieren ser incluidos. ¿De qué manera se puede hacer esto?

\begin{enumerate}[label=\alph*)]
  \item Utilizar el análisis Monte Cario e Internet como herramienta.
  \item Aplicar el método de la ruta crítica.
  \item Determinar las opciones para una acción correctiva recomendada.
  \item Usar la técnica Delphi .
\end{enumerate}

\subsection*{Pregunta 700}
Un director de proyectos está en medio de la ejecución de un proyecto de construcción muy grande cuando descubre que el tiempo que se necesita para concluir el proyecto es mayor al tiempo disponible. ¿Qué es lo MEJOR que puedes hacer?

\begin{enumerate}[label=\alph*)]
  \item Reducir el alcance del producto .
  \item Reunirse con la gerencia y decirles que no se puede cumplir con la fecha requerida.
  \item Trabajar horas extras.
  \item \textbf{Determinar las opciones para la compresión del cronograma y presentar a la gerencia su opción recomendada. 242}
\end{enumerate}

\subsection*{Pregunta 701}
Un director de proyectos está intentando completar un proyecto de desarrollo de software, pero no puede captar atención suficiente para el proyecto. Los recursos están enfocados en la finalización del trabajo relativo al proceso y el director del proyecto tiene poca autoridad para hacer una asignación de recursos. ¿En qué tipo de organización está trabajando el director del proyecto?

\begin{enumerate}[label=\alph*)]
  \item Funcional
  \item Matricial
  \item Expedidor
  \item \textbf{Coordinador}
\end{enumerate}

\subsection*{Pregunta 702}
Un director de proyectos está usando los estimados de duración del promedio ponderado para realizar el análisis de la red del cronograma. ¿Qué tipo de análisis matemático se está utilizando?

\begin{enumerate}[label=\alph*)]
  \item \textbf{Método de la ruta crítica}
  \item Distribución beta
  \item Monte Cario
  \item Nivelación de recursos
\end{enumerate}

\subsection*{Pregunta 703}
Un director de proyectos experimentado acaba de comenzar a trabajar para una empresa grande dedicada a la integración de tecnología de la información. Su gerente le proporciona al director de proyectos una versión preliminar del acta de constitución del proyecto y de inmediato le pide que le entregue un análisis de los riesgos del proyecto. ¿Cuál de las siguientes opciones es la que MEJOR le ayudará a lograrlo?

\begin{enumerate}[label=\alph*)]
  \item Un artículo de la revista PM Network
  \item Su enunciado del alcance del proyecto obtenido del proceso de planificación del proyecto
  \item Su plan de recursos obtenido del proceso de planificación del proyecto
  \item Una conversación con un miembro del equipo de un proyecto similar que fracasó en el pasado
\end{enumerate}

\subsection*{Pregunta 704}
Un director de proyectos ha recibido los estimados de duración de las actividades de su equipo. ¿Cuál de las siguientes necesita con el fin de concluir el proceso Desarrollar el Cronograma?

\begin{enumerate}[label=\alph*)]
  \item Solicitud de cambios
  \item Sistema de control de cambios al cronograma
  \item \textbf{Acciones correctivas recomendadas}
  \item Reservas
\end{enumerate}

\subsection*{Pregunta 705}
Un director de proyectos necesita analizar los costos del proyecto para encontrar maneras de disminuir los costos. Lo MEJOR sería que el director pusiera atención en:

\begin{enumerate}[label=\alph*)]
  \item Los costos variables y fijos.
  \item Los costos fijos e indirectos .
  \item Los costos directos y variables.
  \item \textbf{Los costos indirectos y directos. 285}
\end{enumerate}

\subsection*{Pregunta 706}
Un director de proyectos no tiene mucho tiempo para hacer la planificación, pues ya se aproxima la fecha límite para comenzar . Por lo tanto, busca planificar de la manera más efectiva posible. ¿Qué consejo ofrecerías?

\begin{enumerate}[label=\alph*)]
  \item Asegurarse de que tenga un acta de constitución del proyecto firmada y comenzar con la EDT.
  \item Crear una lista de actividades antes de crear un diagrama de red.
  \item Documentar todos los riesgos conocidos antes de documentar los supuestos de alto nivel.
  \item Finalizar el plan de gestión de calidad antes de determinar las métricas de calidad .
\end{enumerate}

\subsection*{Pregunta 707}
Un director de proyectos nuevo te ha pedido un consejo sobre la creación de una estructura de desglose del trabajo . Después de explicarle el proceso, te pregunta qué software debe usar para crear la EDT, y qué debe hacer con ella cuando esté completa. Quizás respondas que la imagen no es el resultado más valioso de crear una EDT. El resultado más valioso de una EDT es:

\begin{enumerate}[label=\alph*)]
  \item Un diagrama de barras .
  \item Aceptación del equipo.
  \item Actividades.
  \item Una lista de riesgos.
\end{enumerate}

\subsection*{Pregunta 708}
Un director de proyectos nuevo tiene como mentor a un director de proyectos certificado como PMP con más experiencia . Al nuevo director del proyecto se le está dificultando encontrar tiempo suficiente para dirigir el proyecto porque el alcance del proyecto se está elaborando progresivamente. El director de proyectos certificado como PMP aconseja que las herramientas básicas para la dirección de proyectos, como la estructura de desglose del trabajo, puedan usarse durante la ejecución del proyecto para ayudar al director del proyecto. ¿Para cuál de las siguientes puede usarse una estructura de desglose del trabajo?

\begin{enumerate}[label=\alph*)]
  \item Comunicación con el cliente
  \item Mostrar las fechas del calendario para cada paquete de trabajo
  \item Mostrar los gerentes funcionales para cada miembro del equipo
  \item \textbf{Mostrar la necesidad de negocio para el proyecto}
\end{enumerate}

\subsection*{Pregunta 709}
Un director de proyectos para una pequeña compañía de construcción tiene un proyecto que se presupuestó en US\$130.000 durante un período de seis semanas. Según su cronograma, el proyecto debería haber costado US\$60.000 a la fecha. Sin embargo, ha costado US\$90.000. El proyecto también está retrasado con respecto al cronograma, debido a que los estimados originales no fueron exactos. ¿Quién tiene la PRINCIPAL responsabilidad para resolver este problema?

\begin{enumerate}[label=\alph*)]
  \item El director del proyecto
  \item La gerencia senior
  \item El patrocinador del proyecto
  \item \textbf{El director de la oficina de dirección de proyectos 245}
\end{enumerate}

\subsection*{Pregunta 710}
Un director de proyectos quiere involucrar más a los interesados en el proyecto. ¿Cuál de las siguientes opciones sería la MEJOR maner a de lograrlo?

\begin{enumerate}[label=\alph*)]
  \item Hacer que los interesados revisen periódicamente la lista de los requisitos del proyecto.
  \item Invitar a los interesados a asistir a las reuniones de estado del proyecto.
  \item Enviar informes de estado a los interesados.
  \item \textbf{Actualizar a los interesados con respecto al estado de todos los cambios del proyecto.}
\end{enumerate}

\subsection*{Pregunta 711}
Un director de proyectos quizás use \_\_\_ para asegurarse de que los miembros del equipo conocen claramente qué trabajo se incluye en cada uno de sus paquetes de trabajo.

\begin{enumerate}[label=\alph*)]
  \item El enunciado del alcance del proyecto
  \item El alcance del producto
  \item El diccionario de la EDT
  \item El cronograma
\end{enumerate}

\subsection*{Pregunta 712}
Un director de proyectos se está haciendo cargo de un proyecto de otro director de proyectos durante la planificación del proyecto. Si el nuevo director de proyectos quiere ver lo que planificó el director de proyectos anterior para gestionar los cambios al cronograma, sería MEJOR buscar en:

\begin{enumerate}[label=\alph*)]
  \item Plan de gestión de las comunicaciones.
  \item Plan de gestión de las actualizaciones.
  \item Plan de asignación de personal.
  \item Plan de gestión del cronograma.
\end{enumerate}

\subsection*{Pregunta 713}
Un director de proyectos tenía que resolver un problema complejo y facilitó una decisión del equipo acerca de lo que debía hacerse. Unos meses después, el problema reapareció. ¿Cuál de las siguientes acciones es MÁS probable que el director del proyecto NO haya llevado a cabo?

\begin{enumerate}[label=\alph*)]
  \item Realizar el análisis de riesgo adecuado
  \item \textbf{Confirmar que la decisión haya resuelto el problema}
  \item Hacer que el patrocinador del proyecto validara la decisión
  \item Utilizar un diagrama de Ishikawa
\end{enumerate}

\subsection*{Pregunta 714}
Un director de proyectos tiene muy poca experiencia en proyectos, pero ha sido asignado como director de proyectos de un nuevo proyecto. Dado que estará trabajando en una organización matricial para completar su proyecto , puede esperar que las comunicaciones resulten :

\begin{enumerate}[label=\alph*)]
  \item Simples.
  \item Abiertas y precisas.
  \item Complejas.
  \item Difíciles de automatizar.
\end{enumerate}

\subsection*{Pregunta 715}
Un director de proyectos y su equipo de una firma que diseña equipamiento para ferrocarriles tienen la tarea de diseñar una máquina para cargar piedra a los vagones de ferrocarril . El diseño permite el 2 por ciento de derrame, sumando hasta dos toneladas de piedras derramadas por día. ¿En cuál de los siguientes el director del proyecto documenta los procesos de control de calidad, aseguramiento de calidad y mejora de calidad para este proyecto?

\begin{enumerate}[label=\alph*)]
  \item Plan de gestión de la calidad
  \item Política de calidad
  \item Diagramas de control
  \item \textbf{Plan para la dirección del proyecto}
\end{enumerate}

\subsection*{Pregunta 716}
Un gerente se da cuenta de que un director de proyectos está llevando a cabo una reunión con algunos miembros del equipo y algunos interesados para discutir la calidad del proyecto. El cronograma del proyecto ha sido comprimido y el índice del desempeño del costo es 1,1. Han trabajado mucho en el proyecto, el equipo ha sido recompensado de acuerdo con el sistema de recompensas que el director de proyectos estableció, y hay un gran sentimiento de equipo. El gerente sugiere que el director de proyectos no tiene suficiente tiempo para llevar a cabo reuniones sobre la calidad cuando el cronograma está tan comprimido. ¿Cuál de las siguiente s describe MEJOR por qué el gerente está equivocado?

\begin{enumerate}[label=\alph*)]
  \item La calidad mejorada lleva a un aumento de productividad, a un aumento de efectivid ad de costos y a un menor riesgo de costos .
  \item La calidad mejorada lleva a un aumento de productividad, a una disminución de efectividad de costos y a un mayor riesgo de costos .
  \item \textbf{La calidad mejorada lleva a un aume nto de productividad, a un aumento en la efectividad de costos y a un mayor riesgo de costos.}
  \item La calidad mejorada lleva a un aumento de productividad, a una disminución en la efectividad de costos y a un menor riesgo de costos . 323
\end{enumerate}

\subsection*{Pregunta 717}
Un índice de desempeño del costo (CPI) de 0,89 significa:

\begin{enumerate}[label=\alph*)]
  \item En este momento esperamos que el proyecto total cueste 89 por ciento más de lo planificado .
  \item Una vez completado el proyecto, habremos gastado 89 por ciento más de lo planificado .
  \item \textbf{El proyecto progresa a 89 por ciento del ritmo previsto.}
  \item El proyecto está obteniendo 89 centavos de cada dólar invertido .
\end{enumerate}

\subsection*{Pregunta 718}
Un índice de desempeño del cronograma (SPI) de 0,76 significa:

\begin{enumerate}[label=\alph*)]
  \item Te has excedido del presupuesto .
  \item \textbf{Te has adelantado al cronograma .}
  \item Estás progresando a 76 por ciento del ritmo que se planificó originalmente .
  \item Estás progresando a 24 por ciento del ritmo que se planificó originalmente.
\end{enumerate}

\subsection*{Pregunta 719}
Un interesado del proyecto ha sugerido un cambio en el alcance. Después de un examen cuidadoso y mucha argumentación, el comité de control de cambios ha decidido rechazar el cambio. ¿Qué debe hacer el director del proyecto?

\begin{enumerate}[label=\alph*)]
  \item Apoyar al interesado preguntando al comité la razón del rechazo.
  \item Insinuar al interesado que el siguiente cambio que solicite será aprobado.
  \item Documentar el resultado de la solicitud de cambio.
  \item Aconsejar al comité de control de cambios para asegurar que elaboren procesos de aprobación antes de que el siguiente cambio sea propuesto.
\end{enumerate}

\subsection*{Pregunta 720}
Un interesado en particular tiene la reputación de solicitar muchos cambios en los proyectos . ¿Cuál es el MEJOR enfoque que un director de proyectos puede adoptar en el inicio del proyecto para dirigir esta situación?

\begin{enumerate}[label=\alph*)]
  \item \textbf{Decir "no" al interesado unas cuantas veces para disuadirlo de presentar más cambios.}
  \item Hacer que el interesado se involucre en el proyecto lo antes posible .
  \item Hablar con el jefe del interesado para encontrar formas de dirigir las actividades del interesado hacia otro proyecto.
  \item Pedir que no se incluya al interesado en la lista de los interesados.
\end{enumerate}

\subsection*{Pregunta 721}
Un marco de referencia para mantener la organización enfocada en su estrategia general es:

\begin{enumerate}[label=\alph*)]
  \item \textbf{Dirección de proyectos organizacional.}
  \item La Guía del PMBOK º.
  \item Gobierno del proyecto .
  \item Gestión del portafolio .
\end{enumerate}

\subsection*{Pregunta 722}
Un método de estimación de costes utilizando la "regla general" o comparando actividades similares se conoce como estimación\_\_\_\_\_\_\_\_ y se trata de un enfoque \_\_\_\_\_\_\_\_\_. :

\begin{enumerate}[label=\alph*)]
  \item Análoga, descendente
  \item Definitiva, ascendente
\end{enumerate}

\subsection*{Pregunta 723}
Un miembro del equipo constata y así notifica al Project manager que algunos entregables han sido realizados sin haber acabado todos los paquetes de trabajo que se habían definido en la EDT. ¿Qué es lo primero que debe hacer el Project manager?

\begin{enumerate}[label=\alph*)]
  \item Comunicar inmediatamente los cambios a las partes interesadas y miembros del equipo y revisar el desempeño de los miembros del equipo
  \item Revisar los paquetes de trabajo y entregables para evaluar, si lo hay, el impacto en el proyecto y recomendar cambios el plan de dirección del proyecto.
  \item Cambiar las líneas de base de desempeño del proyecto y la EDT
  \item Cambiar el procedimiento de reporte para tratar de identificar estos cambios con más antelación.
\end{enumerate}

\subsection*{Pregunta 724}
Un miembro del equipo de investigación y desarrollo te dice que su trabajo es muy creativo como para proporcionarte un sólo estimado fijo para la actividad. Entre los dos deciden usar las horas de trabajo promedio para desarrollar un prototipo (de proyectos anteriores). ¿Éste es un ejemplo de cuál de las siguientes?

\begin{enumerate}[label=\alph*)]
  \item Estimación paramétrica
  \item Estimación por tres valores
  \item Estimación análoga
  \item \textbf{Análisis Monte Carlo}
\end{enumerate}

\subsection*{Pregunta 725}
Un miembro del equipo del proyecto está hablando con otro miembro del equipo y quejándose de que muchas personas le piden que haga cosas. Si trabaja en una organización funcional, ¿quién tiene el poder de dirigir al miembro del equipo?

\begin{enumerate}[label=\alph*)]
  \item El director del proyecto
  \item El gerente funcional
  \item El equipo
  \item La oficina de dirección de proyectos
\end{enumerate}

\subsection*{Pregunta 726}
Un nuevo proyecto de desarrollo de producto tiene cuatro niveles en la estructura de desglose del trabajo y ha sido secuenciado utilizando el método de diagramación por precedencia. Se han recibido los estimados de duración de las actividades. ¿Qué se debe hacer A CONTINUACIÓN?

\begin{enumerate}[label=\alph*)]
  \item Crear una lista de actividades .
  \item Comenzar la estructura de desglose del trabajo.
  \item Finalizar el cronograma.
  \item \textbf{Comprimir el cronograma . 244}
\end{enumerate}

\subsection*{Pregunta 727}
Un plan de gestión de costos incluye una descripción de:

\begin{enumerate}[label=\alph*)]
  \item Los costos del proyecto .
  \item \textbf{De qué manera se asignan los recursos .}
  \item Los presupuestos y de qué manera se calculan .
  \item El nivel de la EDT en el que se calculará el valor ganado.
\end{enumerate}

\subsection*{Pregunta 728}
Un plan para la dirección del proyecto debe ser realista para poder utilizarse en la dirección de un proyecto. ¿Cuál de los siguientes es el MEJOR método para obtener un plan para la dirección del proyecto realista?

\begin{enumerate}[label=\alph*)]
  \item El patrocinador crea el plan para la dirección del proyecto basado en las entradas del director del proyecto.
  \item El gerente funcional crea el plan para la dirección del proyecto basado en las entradas del director del proyecto.
  \item El director del proyecto crea el plan para la dirección del proyecto basado en las entradas de la gerencia senior .
  \item El director del proyecto crea el plan para la dirección del proyecto basado en las entradas del equipo.
\end{enumerate}

\subsection*{Pregunta 729}
Un presupuesto por fases de tiempo que se utiliza como la base contra la cual se mide el desempeño general del costo se conoce como:

\begin{enumerate}[label=\alph*)]
  \item Costo presupuestado del trabajo realizado.
  \item \textbf{Línea base del costo.}
  \item Estimación hasta la conclusión.
  \item Estimación a la conclusión.
\end{enumerate}

\subsection*{Pregunta 730}
Un producto único o verificable que debe ser producido para considerar el proyecto o la fase del proyecto completa está mejor descrito como:

\begin{enumerate}[label=\alph*)]
  \item Paquetes de trabajo
  \item Requerimientos
  \item Entregables
  \item Actividades
\end{enumerate}

\subsection*{Pregunta 731}
Un proyecto cuenta con varios equipos. En el pasado, el equipo Cha incumplido varias veces con las fechas límites. Esto ha ocasionado que el equipo D tenga que comprimir la ruta crítica en varias ocasiones. Como líder del equipo D, es necesario que te reúnas con:

\begin{enumerate}[label=\alph*)]
  \item El líder del equipo C.
  \item El director del proyecto.
  \item El director del proyecto y la gerencia.
  \item El director del proyecto y el líder del equipo C.
\end{enumerate}

\subsection*{Pregunta 732}
Un proyecto de desarrollo de sistema está cerca del cierre del proyecto cuando de pronto se descubre un riesgo que no había sido identificado previamente. Esto podría afectar potencialmente la posibilidad de cumplir del proyecto. ¿Qué se debe hacer A CONTINUACIÓN?

\begin{enumerate}[label=\alph*)]
  \item Alertar al patrocinador del proyecto acerca de potenciales impactos al costo, alcance o cronograma.
  \item \textbf{Calificar el riesgo.}
  \item Mitigar el riesgo a partir de la creación de un plan de respuesta a los riesgos.
  \item Establecer una solución temporal.
\end{enumerate}

\subsection*{Pregunta 733}
Un proyecto de manufactura tiene un índice de desempeño del cronograma (SPI) de 0,89 y un índice de desempeño del costo (CPI) de 0,91. Por lo general, ¿cuál es la MEJOR explicación de por qué ocurrió esto?

\begin{enumerate}[label=\alph*)]
  \item \textbf{Se cambió el alcance.}
  \item Uno de los proveedores cerró sus puertas y se tuvo que encontrar a uno nuevo.
  \item Se tuvo que comprar equipo adicional.
  \item Una actividad de la ruta crítica tardó más tiempo y se necesitaron más horas de trabajo para completarla.
\end{enumerate}

\subsection*{Pregunta 734}
Un proyecto está a la mitad del esfuerzo de ejecución cuando un interesado sugiere un nuevo cambio importante . Este cambio causará la tercera revisión importante del proyecto. Al mismo tiempo, el director del proyecto descubre que un paquete de trabajo importante no se completó debido a que el jefe de uno de los miembros del equipo lo trasladó a un proyecto de mayor prioridad. ¿Cuál de las siguientes personas es la MÁS adecuada para discutir este asunto con el director del proyecto?

\begin{enumerate}[label=\alph*)]
  \item El equipo
  \item \textbf{La gerencia senior}
  \item El cliente
  \item El patrocinador
\end{enumerate}

\subsection*{Pregunta 735}
Un proyecto está plagado de cambios al acta de constitución del proyecto. ¿Quién tiene la responsabilidad principal de decidir si estos cambios son necesarios?

\begin{enumerate}[label=\alph*)]
  \item El director del proyecto
  \item El equipo del proyecto
  \item El patrocinador
  \item \textbf{Los interesados 144}
\end{enumerate}

\subsection*{Pregunta 736}
Un proyecto tiene tres caminos críticos. ¿Cuál de las siguientes afirmaciones describe mejor cómo afecta esto al proyecto?

\begin{enumerate}[label=\alph*)]
  \item Lo hace más caro
  \item Incrementa los riesgos del proyecto
  \item Lo hace más fácil de dirigir
  \item Requiere de más personas
\end{enumerate}

\subsection*{Pregunta 737}
Un proyecto tiene tres rutas críticas. ¿Cuál de las siguientes describe MEJOR cómo afecta esto el proyecto?

\begin{enumerate}[label=\alph*)]
  \item \textbf{Lo hace más fácil de dirigir .}
  \item Incrementa los riesgos del proyecto.
  \item Requiere de más personas .
  \item Lo hace más caro.
\end{enumerate}

\subsection*{Pregunta 738}
Un sistema de autorización de trabajo puede utilizarse para:

\begin{enumerate}[label=\alph*)]
  \item \textbf{Gestionar quién hace cada actividad.}
  \item Gestionar cuándo y en qué orden se realiza el trabajo.
  \item Gestionar cuándo se realiza cada actividad.
  \item Gestionar quién hace cada actividad y cuándo se realiza.
\end{enumerate}

\subsection*{Pregunta 739}
Uno de los entregables ya está listo y revisado por la compañía. ¿Cuál será el siguiente paso a dar?

\begin{enumerate}[label=\alph*)]
  \item Enviarlo al cliente para que controle el alcance
  \item Enviarlo al cliente para que verifique el alcance
  \item Solicitar al cliente la aceptación formal del proyecto
  \item Comprobar que no es necesario gestionar algún Cambio en el alcance
\end{enumerate}

\subsection*{Pregunta 740}
Uno de los grandes problemas de los proyectos es poder cuadrar las agendas y compromisos de todas las partes interesadas y es una de las tareas que más tiempo puede llevar a un Project manager. ¿Cuál de los siguientes resultados no forma parte de las salidas del proceso de planificación de recursos humanos?

\begin{enumerate}[label=\alph*)]
  \item RACI
  \item RAM
  \item Evaluación de la eficiencia y disponibilidad del personal
  \item Plan de recursos humanos
\end{enumerate}

\subsection*{Pregunta 741}
Uno de los interesados del proyecto se pone en contacto con el director de proyectos para hablar sobre cierto alcance adicional que quieren añadir al proyecto. El director de proyectos pide los detalles por escrito y luego trabaja en el proceso Controlar el Alcance. ¿Qué es lo que debe hacer el director del proyecto DESPUÉS, cuando la evaluación del alcance solicitado está finalizada?

\begin{enumerate}[label=\alph*)]
  \item Preguntar al interesado si hay más cambios esperados.
  \item Completar el control integrado de cambios.
  \item \textbf{Asegurarse de que el interesado comprenda el impacto del cambio .}
  \item Hallar la causa raíz por la cual no se descubrió el alcance durante la planificación del proyecto.
\end{enumerate}

\subsection*{Pregunta 742}
Uno de los miembros de su equipo ha presentado una solicitud de cambio al alcance. El cliente se muestra interesado y le ha pedido que reúna al Comité de ana Control de Cambios. Como parte del protocolo, usted debe acompañar al miembro del equipo del proyecto a la reunión con el Comité. ¿Cuál de los bre 1.00

\begin{enumerate}[label=\alph*)]
  \item documentos mencionados es el más importante y que el director de proyecto debe llevar consigo a dicha reunión?
  \item Estimados de los paquetes de trabajo.
  \item \textbf{Línea de base de alcance. Blan del proyecto.}
  \item EDT.
\end{enumerate}

\subsection*{Pregunta 743}
Uno de los miembros de tu equipo te informa que no sabe cuál de los muchos proyectos en los que está trabajando es el más importante. ¿Quién debe determinar prioridades entre los proyectos de una compañía?

\begin{enumerate}[label=\alph*)]
  \item El director del proyecto
  \item El equipo de dirección del proyecto
  \item La oficina de dirección de proyectos (PMO)
  \item El equipo
\end{enumerate}

\subsection*{Pregunta 744}
Uno de los proveedores que ha presentado una oferta para un contrato de adquisición desea quedar contigo como Project manager del proyecto para revisar algunos aspectos de la oferta y del contrato. ¿Qué es lo mejor que puedes hacer?

\begin{enumerate}[label=\alph*)]
  \item Indicarle que los cauces para las aclaraciones al concurso están establecidas y que debe seguir esos procedimientos y, al acabar la conversación, registrar su solicitud.
  \item Resolverle las dudas que pueda tener para posteriormente enviar una nota al resto de ofertantes con la respuesta a esas dudas.
  \item Tratar de eludir el tema y desviar la conversación a aspectos no relacionados con el contrato en sí.
  \item Solicitarle que esta conversación no se haga en horario de trabajo y que la tratéis en un almuerzo informal.
\end{enumerate}

\subsection*{Pregunta 745}
Usted acaba de pasar por una difícil revisión del proyecto con su patrocinador. Hubo más problemas de lo que esperaban y parecía no haber un orden lógico en la revisión de las cuestiones que se examinaron. Durante la próxima revisión, ¿qué herramientas debe utilizar para determinar la prioridad? :

\begin{enumerate}[label=\alph*)]
  \item Plan de gestión de riesgos para ayudar a determinar cómo se puede evaluar el riesgo.
  \item Diagramas de espina de pescado.
  \item Diagrama de Pareto para ayudarle a concentrarse en los asuntos más críticos.
  \item Proceso de inspección eficiente.
\end{enumerate}

\subsection*{Pregunta 746}
Usted debe ejecutar medidas de control de calidad durante la ejecución. El resultado de estos procesos son:

\begin{enumerate}[label=\alph*)]
  \item Acciones correctivas recomendadas.
  \item \textbf{Mediciones de control de calidad.}
  \item Entregables verificados.
  \item Solicitudes de cambio.
\end{enumerate}

\subsection*{Pregunta 747}
Usted e explica que una variación en el cronograma positiva acompañada de una variación de costo aún negativa puede ser indicativa de todo lo siguiente EXCEPTO: 00

\begin{enumerate}[label=\alph*)]
  \item Tener miembros del equipo con pagos por hora y con trabajo en horas extras
  \item Tasas de productividad menores a las esperadas
  \item Uso de gente con salarios más altos que trabajan en forma más eficiente
  \item Agregar recursos adicionales al proyecto
\end{enumerate}

\subsection*{Pregunta 748}
Usted es el director de proyecto para la construcción de un importante centro de datos cerca de San Andrés en el Área de la Bahía de San Francisco, donde los terremotos son altamente probables. ¿Cuál es una estrategia razonable de respuesta al riesgo para lidiar con el impacto de los terremotos?

\begin{enumerate}[label=\alph*)]
  \item Mitigar
  \item Mejorar
  \item Compartir
  \item \textbf{Aceptar}
\end{enumerate}

\subsection*{Pregunta 749}
Usted esta a cargo de un equipo global de proyecto. Cada vez que se presenta Usted está a cargo de un equipo global de proyecto. Cada vez que se presenta un problema, su equipo decide enviar correos electrénicos a todos los que se un problema, su equipo decide enviar correos electrónicos a todos los que se incluyen en la lista de distribucidn del proyecto. Esta situacidn esta fuera de incluyen en la lista de distribución del proyecto. Esta situación está fuera de control y hay mucha informacién en manos de personas que no necesitan control y hay mucha información en manos de personas que no necesitan conocerla o que no han mostrado interés en el proyecto. Qué debera hacer conocerla o que no han mostrado interés en el proyecto. ¿Qué deberá hacer ¥ Marear para corregir esta situacién? para corregir esta situación? cuestién :

\begin{enumerate}[label=\alph*)]
  \item Pedirle a su equipo que investigue este problema y que le presente una recomendacion. a. Pedirle a su equipo que investigue este problema y que le presente una recomendación.
  \item \textbf{Si ya se establecieron las reglas, hacer que el equipo les preste atencidn. De lo b. Si ya se establecieron las reglas, hacer que el equipo les preste atención. De lo contrario, actualizar el plan de comunicacién con los lineamientos para comunicacidn a contrario, actualizar el plan de comunicación con los lineamientos para comunicación a través de correo electrénico. través de correo electrónico.}
  \item Publicar y distribuir los lineamientos y pedir a los miembros del equipo que lean la c. Publicar y distribuir los lineamientos y pedir a los miembros del equipo que lean la seccioén correspondiente a los correos electrénicos. sección correspondiente a los correos electrónicos.
  \item Enviar un correo electronico a los miembros del equipo y hacer que detengan esta d. Enviar un correo electrónico a los miembros del equipo y hacer que detengan esta practica relacionada con la distribucién de informacidn. práctica relacionada con la distribución de información.
\end{enumerate}

\subsection*{Pregunta 750}
Usted ha estimado y realizado un cronograma de su proyecto en función de los miembros del equipo que se le asignaron, lo que le ha permitido determinar la línea base de costo. Sin embargo, ahora tiene nuevos miembros del equipo que tienen más experiencia y salarios más altos que los que fueron incluidos en su estimado inicial. El impacto más probable en el desempeño del proyecto sería: (SV= Shedule variance o desviación en tiempos CV=Cost variance o desviación en costes)

\begin{enumerate}[label=\alph*)]
  \item CV Negativa, SV negativa
  \item CV positiva, SV negativa
  \item \textbf{SV positiva, CV negativa}
  \item CV positiva
\end{enumerate}

\subsection*{Pregunta 751}
Usted ha sido contratado como experto para realizar el plan de gestión de las comunicaciones de un proyecto industrial, ¿sobre cuál de los siguientes elementos NO centrará su trabajo? o más de una:

\begin{enumerate}[label=\alph*)]
  \item \textbf{Definir cómo se comunicará la información}
  \item Identificar los requisitos de información de los interesados
  \item Monitorear la efectividad de las comunicaciones
  \item Establecer el método y la tecnología de comunicación
\end{enumerate}

\subsection*{Pregunta 752}
Usted se reúne con el equipo del proyecto para revisar el Análisis Costo- Beneficio, Gráficos de Control, Diseño de Experimentos, y otras metodologías de gestión de calidad corporativas. Lo más probable es que usted y su equipo estén enfocados en:

\begin{enumerate}[label=\alph*)]
  \item Realizar el Aseguramiento de Calidad.
  \item \textbf{Control la Calidad.}
  \item Gestionar la Calidad.
  \item Planificar la Gestión de la Calidad.
\end{enumerate}

\subsection*{Pregunta 753}
Utilizar escalas numéricas para evaluar la probabilidad e impacto se considera como parte del proceso de \_\_\_\_\_\_\_\_\_\_\_\_\_.

\begin{enumerate}[label=\alph*)]
  \item Análisis cuantitativo de riesgo
  \item Planificación de respuesta al riesgo
  \item Identificación de riesgo
  \item Análisis cualitativo de riesgo
\end{enumerate}

\subsection*{Pregunta 754}
Validar el Alcance está estrechamente relacionado con:

\begin{enumerate}[label=\alph*)]
  \item Realizar el Control de Calidad.
  \item Secuenciar las Actividades.
  \item \textbf{Realizar el Aseguramiento de Calidad.}
  \item Gestión del tiempo.
\end{enumerate}

\subsection*{Pregunta 755}
Vas a comenzar la definición de actividades de su proyecto. Tienes ya el Enunciado del Trabajo, la EDT y su Diccionario y los procedimientos establecidos por la organización para Definir las actividades. ¿Qué más necesitas?

\begin{enumerate}[label=\alph*)]
  \item Identificar requerimientos que afecten a las actividades
  \item Identificar hipótesis de partida, importantes para identificar riesgos
  \item Obtener información de las partes interesadas
  \item Recoger información histórica de proyectos similares
\end{enumerate}

\subsection*{Pregunta 756}
Vas a tener tu primera reunión con los miembros del equipo que trabajará en una oferta transnacional y que, por ello, pertenecerán a diferentes culturas, religiones y países. ¿Qué es lo mejor que puedes hacer?

\begin{enumerate}[label=\alph*)]
  \item Informarse de las culturas del país, preparar un dossier informativo a entregar a todo el equipo de manera que sepan lo que deben y no deben hacer.
  \item Evitar los choques culturales asignando paquetes de trabajo independientes a cada grupo cultural de manera que las relaciones y coordinación las lleve directamente el Project manager.
  \item Informar que cada miembro del equipo puede expresarse libremente tal y como les indica su cultura y que el resto traten de respetar y asumir estas diferencias.
  \item Evitar las comunicaciones cara a cara y tratar de que las relaciones sean por escrito para evitar problemas de comunicación y comportamientos no verbales tan diferentes entre culturas.
\end{enumerate}

\subsection*{Pregunta 757}
Una vez analizado el cronograma que has creado, ves la necesidad de reducir la duración total del proyecto y optas por el Fast-tracking lo que significa:

\begin{enumerate}[label=\alph*)]
  \item Que comienzas el proyecto antes y trabajas horas extra
  \item Asignas más tareas aumentando el coste total, especialmente en las actividades del camino crítico
  \item Comienzas las actividades antes y las solapas
  \item Reduces la duración de las actividades y solicitas al equipo de trabajo que trabaje horas extra
\end{enumerate}

\subsection*{Pregunta 758}
Una vez firmado un contrato , es legalmente vinculante salvo que

\begin{enumerate}[label=\alph*)]
  \item El proveedor no es capaz técnicamente de ejecutar el trabajo
  \item El comprador no efectúa los pagos en los plazos previstos
  \item El contrato contravenga alguna ley estatal o local
  \item Una de las partes no está de acuerdo con el contenido.
\end{enumerate}

\subsection*{Pregunta 759}
Una vez identificadas las partes interesadas de tu proyecto, debes proceder a:

\begin{enumerate}[label=\alph*)]
  \item Actualizar el acta de constitución
  \item Crear el plan para la dirección del proyecto
  \item Crear el plan de gestión de las comunicaciones
  \item Crear el registro de partes interesadas
\end{enumerate}

\subsection*{Pregunta 760}
Una vez que el contratista entrega toda la documentación correspondiente al producto en formato digital y perfectamente completa, el Project manager le solicita dos copias de la documentación en formato papel a lo que el proveedor se niega aduciendo que el contrato indicaba entregar la información completa en un formato compatible, legible y utilizable. ¿En cuál de estas fases debía haber indicado el Project manager que deseaba tener dos copias en papel?

\begin{enumerate}[label=\alph*)]
  \item El Enunciado del Alcance del producto
  \item En la Verificación del Alcance
  \item En el Plan de Adquisición y Compras
  \item En el Plan de Control de Calidad
\end{enumerate}

\subsection*{Pregunta 761}
Una vez que tienes tu Lista de actividades y vas a preparar el diagrama de dependencias entre actividades utilizando la técnica PDM (Método de Diagrama de Precedencias), cuál de las siguientes afirmaciones es cierta

\begin{enumerate}[label=\alph*)]
  \item PDM es igual al método de diagrama AON y utiliza únicamente una estimación de duración.
  \item PDM es igual al método de diagramación AOA y utiliza relaciones lógicas
  \item PDM es igual al método ADM y su relación de precedencia más habitual es Final a inicio
  \item PDM es igual que el método GERT lo que permite condiciones, ramas e iteraciones
\end{enumerate}

\end{document}
